% Dosages d'oxydoréduction — Chimie Tle TI — Cours signé OnBuch+
% Programme officiel MINESEC, Terminale TI. Compiler avec : tectonic L11-dosages-oxydoreduction.tex
\newcommand{\DOCMATIERE}{Chimie}
\newcommand{\DOCNIVEAU}{Tle TI}
\newcommand{\DOCMODULE}{Module 3 · Oxydoréduction}
\newcommand{\DOCLECON}{11}
\newcommand{\DOCTITRE}{Dosages d'oxydoréduction}
% OnBuch+ — préambule « cours signés » ENRICHI (compilé avec tectonic / XeLaTeX)
% Système visuel « Pop » d'OnBuch+ : fond crème (jamais blanc), bordures encre
% épaisses, ombres « dures » décalées, titres Archivo Black en capitales.
% Version enrichie pour les cours de chimie : chimie (mhchem, chemfig),
% graphes (pgfplots), boîtes pédagogiques supplémentaires (définition,
% propriété, méthode, attention, le savais-tu, exercice résolu, exercice,
% TP, bilan), page de garde refaite et signature OnBuch+ ancrée en bas.
%
% Macros attendues dans main.tex AVANT \input{preamble} :
%   \DOCMATIERE  \DOCNIVEAU  \DOCMODULE  \DOCLECON (numéro)  \DOCTITRE
\documentclass[11pt]{article}

\usepackage{fontspec}
\usepackage[french]{babel}
\usepackage[a4paper,margin=2cm,top=3.4cm,bottom=2.8cm,headheight=1.4cm,footskip=1.3cm]{geometry}
\usepackage{amsmath,amssymb,amsfonts}
\usepackage[table]{xcolor} % [table] : fournit aussi \rowcolors (alternance de lignes)
\usepackage{pagecolor}
\usepackage{tikz}
\usetikzlibrary{arrows.meta,positioning,calc,shapes.geometric,shapes.misc,shapes.arrows,decorations.pathmorphing,decorations.markings,patterns,shadows,mindmap,trees,fit,backgrounds,matrix,intersections}
\usepackage{pgfplots}
\pgfplotsset{compat=1.18}
\usepackage[version=4]{mhchem}
\usepackage{chemfig}
\usepackage{enumitem}
\usepackage{fancyhdr}
% [most] charge toutes les bibliothèques tcolorbox (listings, minted, raster,
% documentation...) et gonfle inutilement la mémoire TeX sur un document long
% (~300 boîtes) : on ne charge que ce qui est réellement utilisé.
\usepackage[skins,breakable]{tcolorbox}
\usepackage{graphicx}
\usepackage{colortbl}
\usepackage{booktabs}
\usepackage{tabularx}
\usepackage{array}
\usepackage{multicol}
\usepackage{siunitx}
% parse-numbers=false : accepte aussi \SI{2,5 \times 10^{-3}}{...}
\sisetup{locale=FR,output-decimal-marker={,},group-minimum-digits=4,parse-numbers=false}
\DeclareSIUnit\ohm{\ensuremath{\symup{\Omega}}} % U+2126 absent de la police
\usepackage{qrcode}
% \degree : commande fréquemment utilisée par les modèles pour un angle ou une
% température en dehors de \SI{}{} (non fournie par les paquets chargés).
\providecommand{\degree}{^{\circ}}
\usepackage{lastpage}
\usepackage{hyperref}
\hypersetup{hidelinks}

% ============================================================
% Palette « Pop » OnBuch+ — mêmes hex que la classe P (plus_ui.dart)
% ============================================================
\definecolor{poporange}{HTML}{F59321}
\definecolor{popdark}{HTML}{E07A0C}
\definecolor{popcream}{HTML}{FFF6E8}
\definecolor{popink}{HTML}{1C1714}
\definecolor{poppurple}{HTML}{7A5AE0}
\definecolor{popgreen}{HTML}{2E9E6A}
\definecolor{poppink}{HTML}{E0457B}
\definecolor{popblue}{HTML}{2F7DE1}
\definecolor{popgold}{HTML}{C98A12}
\definecolor{popmuted}{HTML}{8A7F76}
\definecolor{popline}{HTML}{EADFD2}
\definecolor{popgray}{HTML}{8A7F76}
\colorlet{popgrey}{popgray}
% Alias tolérants (le modèle nomme parfois les couleurs différemment)
\colorlet{popwhite}{white}
\colorlet{popblack}{popink}
\colorlet{popred}{poppink}
\colorlet{popyellow}{popgold}
% Teintes claires pour fonds de tableaux / zones
\colorlet{popblueL}{popblue!10!white}
\colorlet{poporangeL}{poporange!14!white}
\colorlet{popgreenL}{popgreen!12!white}
\colorlet{poppurpleL}{poppurple!12!white}
\colorlet{poppinkL}{poppink!12!white}
\colorlet{popgoldL}{popgold!14!white}

\pagecolor{popcream}

% ============================================================
% Typographie
% ============================================================
\defaultfontfeatures{Ligatures=TeX}
\setmainfont{PlusJakartaSans-Medium.ttf}[Path=../../../fonts/,BoldFont=PlusJakartaSans-Bold.ttf]
% Maths en sans-serif (Fira Math) avec chiffres et lettres droites de Plus
% Jakarta, pour que les formules se fondent dans le texte.
\usepackage[math-style=ISO,bold-style=ISO]{unicode-math}
\setmathfont{FiraMath-Regular.otf}[Path=../../../fonts/]
\setmathfont{PlusJakartaSans-Medium.ttf}[Path=../../../fonts/,range={up/num,up/{Latin,latin},\mathup}]
\usepackage{icomma} % virgule décimale sans espace en mode math
% Caractères absents de la police de texte (sinon carré vide) : rendus par
% leur équivalent mathématique, en texte comme en formule.
\usepackage{newunicodechar}
\newunicodechar{α}{\ensuremath{\alpha}}
\newunicodechar{Α}{\ensuremath{\alpha}}
\newunicodechar{β}{\ensuremath{\beta}}
\newunicodechar{δ}{\ensuremath{\delta}}
\newunicodechar{✓}{\popcheck{popgreen}}
\newfontfamily\popdisplay{ArchivoBlack-Regular.ttf}[Path=../../../fonts/]
\newfontfamily\poplogo{SpaceGrotesk-Bold.ttf}[Path=../../../fonts/]
\newfontfamily\popsemi{PlusJakartaSans-SemiBold.ttf}[Path=../../../fonts/]
\newfontfamily\popxbold{PlusJakartaSans-ExtraBold.ttf}[Path=../../../fonts/]
\newfontfamily\popxboldls{PlusJakartaSans-ExtraBold.ttf}[Path=../../../fonts/,LetterSpace=3.0]

\setlist[enumerate]{leftmargin=1.7em,itemsep=5pt,topsep=4pt}
\setlist[itemize]{leftmargin=1.7em,itemsep=4pt,topsep=4pt,label={\color{poporange}\textbullet}}
\setlength{\parindent}{0pt}
\setlength{\parskip}{5pt}
\renewcommand{\arraystretch}{1.25}

% chemfig : liaisons un peu plus épaisses, encre Pop
\setchemfig{atom sep=2em,bond style={line width=0.9pt,popink},cram width=3pt}

% pgfplots : style Pop par défaut
\pgfplotsset{
  every axis/.append style={axis line style={popink,line width=1pt},tick style={popink},
    grid style={popline,line width=0.5pt},label style={font=\small},tick label style={font=\footnotesize}},
  popcurve/.style={poporange,line width=1.8pt,no markers,smooth},
}

% ============================================================
% Pill / squiggle
% ============================================================
\newcommand{\poppill}[3]{%
  \tikz[baseline=-0.55ex]\node[draw=popink,line width=1.2pt,rounded corners=2.6pt,%
    fill=#2,inner xsep=6.5pt,inner ysep=3pt]{%
    \popxboldls\fontsize{7}{7}\selectfont\color{#3}\MakeUppercase{#1}};%
}
\newcommand{\popsquiggle}[1]{%
  \tikz[baseline=-0.4ex]{
    \draw[#1,line width=2.6pt,line cap=round]
      plot [smooth] coordinates {(0,0.09) (0.34,0.01) (0.68,0.21) (1.02,0.01) (1.36,0.10)};
  }%
}
% Mot-clé surligné (marqueur orange sous le texte)
\newcommand{\cle}[1]{{\popxbold\color{popdark}#1}}

% ============================================================
% En-tête / pied
% ============================================================
\pagestyle{fancy}
\fancyhf{}
\renewcommand{\headrulewidth}{0pt}
\renewcommand{\footrulewidth}{0pt}
\lhead{\raisebox{-0.15ex}{\poplogo\fontsize{13}{13}\selectfont\color{popink}OnBuch\color{poporange}+}}
\rhead{\poppill{\DOCMATIERE\ \textbullet\ \DOCNIVEAU\ \textbullet\ Leçon \DOCLECON}{white}{popink}}
\lfoot{\popxbold\fontsize{7.6}{7.6}\selectfont\color{popmuted}\MakeUppercase{Cours signé OnBuch+}\ \textbullet\ {\popsemi\DOCTITRE}}
\rfoot{\tikz[baseline=-0.6ex]\node[rounded corners=8.5pt,fill=popink,minimum height=17pt,minimum width=17pt,inner xsep=5pt,inner sep=0]{\popxbold\fontsize{8}{8}\selectfont\color{popcream}\thepage\,/\,\pageref*{LastPage}};}

% ============================================================
% Titres
% ============================================================
\newcommand{\chaptitre}[2]{%
  \begin{tcolorbox}[enhanced,colback=poporange,colframe=popink,boxrule=2.6pt,arc=11pt,
      drop shadow=popink,left=15pt,right=15pt,top=12pt,bottom=11pt,before skip=3pt,after skip=18pt]
    \poppill{#2}{popink}{popcream}\par\vspace{9pt}
    {\raggedright\hyphenpenalty=10000\exhyphenpenalty=10000\popdisplay\fontsize{22}{24}\selectfont\color{popcream}\MakeUppercase{#1}\par}
  \end{tcolorbox}
}
% Section numérotée : pastille orange avec le numéro + titre Archivo
\newcounter{popsec}
\newcommand{\coursec}[1]{%
  \stepcounter{popsec}\setcounter{popsub}{0}%
  \par\vspace{16pt}\noindent
  \begin{minipage}{\linewidth}
  \tikz[baseline=-0.7ex]\node[draw=popink,line width=1.6pt,fill=poporange,rounded corners=4pt,minimum size=22pt,inner sep=0]{\popdisplay\fontsize{12}{12}\selectfont\color{popcream}\thepopsec};%
  \hspace{8pt}{\popdisplay\fontsize{15}{17}\selectfont\color{popink}\MakeUppercase{#1}}\par
  \vspace{4pt}\noindent{\color{poporange}\rule{36pt}{3.6pt}}
  \end{minipage}\par\nopagebreak\vspace{8pt}
}
\newcounter{popsub}
\newcommand{\courssub}[1]{%
  \stepcounter{popsub}%
  \par\vspace{9pt}\noindent{\popxbold\fontsize{11.5}{13}\selectfont\color{popink}{\color{poporange}\thepopsec.\thepopsub}\hspace{6pt}#1}\par\nopagebreak\vspace{4pt}
}

% ============================================================
% Boîtes pédagogiques — même recette : fond blanc, bordure épaisse colorée,
% ombre dure encre, badge plein. Titre optionnel [#1].
% ============================================================
\tcbset{popbase/.style={breakable,enhanced,colback=white,boxrule=2.3pt,arc=9pt,
  shadow={3pt}{-3pt}{0pt}{popink},left=14pt,right=14pt,top=11pt,bottom=11pt,before skip=11pt,after skip=15pt,
  fontupper=\popsemi\fontsize{10.6}{14.5}\selectfont\color{popink},
  fontlower=\popsemi\fontsize{10.6}{14.5}\selectfont\color{popink}}}
% Coche dessinée (le glyphe ✓ n'existe pas dans les polices du document)
\newcommand{\popcheck}[1]{\tikz[baseline=-0.5ex]\draw[#1,line width=1.6pt,line cap=round,line join=round](-0.13,0.0)--(-0.04,-0.09)--(0.13,0.1);}
\newcommand{\popboxhead}[3]{\poppill{#1}{#2}{white}\ifx\relax#3\relax\else\hspace{6pt}{\popxbold\fontsize{10.6}{12}\selectfont\color{popink}#3}\fi\par\vspace{7pt}}

\newtcolorbox{aretenir}[1][]{popbase,colframe=popgreen,before upper={\popboxhead{À retenir}{popgreen}{#1}}}
\newtcolorbox{exemplebox}[1][]{popbase,colframe=poporange,before upper={\popboxhead{Exemple}{poporange}{#1}}}
\newtcolorbox{definition}[1][]{popbase,colframe=popblue,before upper={\popboxhead{Définition}{popblue}{#1}}}
\newtcolorbox{propriete}[1][]{popbase,colframe=poppurple,before upper={\popboxhead{Propriété}{poppurple}{#1}}}
\newtcolorbox{methode}[1][]{popbase,colframe=popgold,before upper={\popboxhead{Méthode}{popgold}{#1}}}
\newtcolorbox{attention}[1][]{popbase,colframe=poppink,before upper={\popboxhead{Attention}{poppink}{#1}}}
\newtcolorbox{experience}[1][]{popbase,colframe=popblue,colback=popblueL,before upper={\popboxhead{Expérience}{popblue}{#1}}}
% « Le savais-tu ? » : carte violette pleine (contexte, culture, Cameroun)
\newtcolorbox{savaistu}[1][]{popbase,colback=poppurple,colframe=popink,
  fontupper=\popsemi\fontsize{10.4}{14.2}\selectfont\color{white},
  before upper={\poppill{Le savais-tu\,?}{white}{poppurple}\ifx\relax#1\relax\else\hspace{6pt}{\popxbold\color{white}#1}\fi\par\vspace{7pt}}}
% Exercice résolu : énoncé en haut, \tcblower puis la solution sur fond crème
\newcounter{popexo}\newcounter{popexr}
\newtcolorbox{exoresolu}[1][]{popbase,colframe=poporange,colbacklower=popcream,
  segmentation style={popink,line width=1.2pt,dashed},
  before upper={\stepcounter{popexr}\popboxhead{Exercice résolu \thepopexr}{poporange}{#1}},
  before lower={\poppill{Solution}{popink}{popcream}\par\vspace{6pt}}}
% Exercice (non résolu en place) — la correction est dans la partie Corrigés
\newtcolorbox{exercice}[1][]{popbase,colframe=popink,
  before upper={\stepcounter{popexo}\popboxhead{Exercice \thepopexo}{popink}{#1}}}
% Corrigé
\newtcolorbox{corrige}[1][]{popbase,colframe=popgreen,colback=popgreenL,
  before upper={\popboxhead{Corrigé}{popgreen}{#1}}}
% Objectifs / prérequis / bilan
\newtcolorbox{objectifs}{popbase,colframe=popink,colback=poporangeL,
  before upper={\popboxhead{Objectifs de la leçon}{popink}{}}}
\newtcolorbox{prerequis}{popbase,colframe=popink,colback=popblueL,
  before upper={\popboxhead{Prérequis}{popblue}{}}}
\newtcolorbox{bilan}{popbase,colframe=popink,colback=white,boxrule=2.8pt,
  before upper={\popboxhead{Fiche bilan}{poporange}{L'essentiel à connaître pour l'examen}}}
% Figure encadrée avec légende
\newtcolorbox{popfigure}[1][]{enhanced,breakable=false,colback=white,colframe=popline,boxrule=1.4pt,arc=8pt,
  left=10pt,right=10pt,top=10pt,bottom=8pt,before skip=10pt,after skip=12pt,halign=center,
  lower separated=false,halign lower=center,
  fontlower=\popsemi\fontsize{9}{11.5}\selectfont\color{popmuted},
  #1}
\newcommand{\legende}[1]{\par\vspace{6pt}{\popsemi\fontsize{9}{11.5}\selectfont\color{popmuted}#1}}

% ============================================================
% Signature OnBuch+
% ============================================================
\newcommand{\poplogoinline}[1]{{\poplogo\fontsize{#1}{#1}\selectfont\color{popink}OnBuch\color{poporange}+}}
\newcommand{\signatureonbuch}{%
  \begin{tcolorbox}[enhanced,colback=popink,colframe=popink,boxrule=2.2pt,arc=10pt,
      drop shadow=poporange,left=14pt,right=14pt,top=10pt,bottom=10pt,before skip=0pt,after skip=0pt]
    \tikz[baseline=-0.7ex]\node[circle,fill=popgreen,draw=popcream,line width=1.3pt,minimum size=22pt,inner sep=0]{\popcheck{white}};%
    \hspace{9pt}\begin{minipage}[c]{0.62\linewidth}
      {\popxbold\fontsize{12}{13}\selectfont\color{popcream}Signé\ \textbullet\ {\poplogo OnBuch\color{poporange}+}}\par\vspace{2pt}
      {\raggedright\popsemi\fontsize{8}{10.5}\selectfont\color{popline}\MakeUppercase{Équipe pédagogique OnBuch+}\\\MakeUppercase{Conforme au programme officiel MINESEC}\par}
    \end{minipage}\hfill
    {\poplogo\fontsize{22}{22}\selectfont\color{popcream}OnBuch\color{poporange}+}
  \end{tcolorbox}%
}

% ============================================================
% Page de garde
% #1 titre · #2 matière · #3 niveau · #4 module · #5 n° leçon · #6 durée ·
% #7 description / objectifs (texte ou itemize)
% ============================================================
\newcommand{\pagegarde}[7]{%
  \thispagestyle{empty}
  \newgeometry{a4paper,margin=1.7cm,top=1.6cm,bottom=1.4cm}%
  \begin{tikzpicture}[remember picture,overlay]
    \fill[poporange!18] ([xshift=-3.2cm,yshift=-3.6cm]current page.north east) circle (4.6cm);
    \fill[poppurple!14] ([xshift=2.4cm,yshift=7.6cm]current page.south west) circle (3.8cm);
  \end{tikzpicture}%
  {\poplogo\fontsize{30}{30}\selectfont\color{popink}OnBuch\color{poporange}+}\hfill
  \raisebox{0.6ex}{\poppill{Cours signé \textbullet\ Premium}{popink}{popcream}}\par
  \vspace{1pt}\popsquiggle{poppurple}\par
  \vspace{8pt}
  \begin{tcolorbox}[enhanced,colback=poporange,colframe=popink,boxrule=2.8pt,arc=13pt,
      drop shadow=popink,left=18pt,right=18pt,top=12pt,bottom=13pt,after skip=10pt]
    \poppill{#2 \textbullet\ #3}{popink}{popcream}\hspace{5pt}\poppill{#4}{white}{popink}\par\vspace{8pt}
    {\popdisplay\fontsize{14}{14}\selectfont\color{popink}LEÇON #5}\par\vspace{4pt}
    {\raggedright\hyphenpenalty=10000\exhyphenpenalty=10000\popdisplay\fontsize{25}{27}\selectfont\color{popcream}\MakeUppercase{#1}\par}
    \vspace{8pt}
    {\popxbold\fontsize{9.5}{9.5}\selectfont\color{popink}\MakeUppercase{Durée conseillée : #6}}
  \end{tcolorbox}
  \begin{tcolorbox}[enhanced,colback=white,colframe=popink,boxrule=2.2pt,arc=10pt,
      drop shadow=popink,left=16pt,right=16pt,top=10pt,bottom=10pt,after skip=8pt]
    \poppill{Dans cette leçon}{poppurple}{white}\par\vspace{5pt}
    \popsemi\fontsize{9.3}{12.6}\selectfont\color{popink}#7
  \end{tcolorbox}
  \noindent\begin{minipage}[t]{0.487\linewidth}
    \begin{tcolorbox}[enhanced,colback=popgreen,colframe=popink,boxrule=2.2pt,arc=10pt,
        drop shadow=popink,left=11pt,right=11pt,top=10pt,bottom=10pt,height=3.1cm,valign=center]
      \tcbox[colback=white,colframe=popink,boxrule=1.3pt,arc=5pt,boxsep=0pt,nobeforeafter,
        left=3pt,right=3pt,top=3pt,bottom=3pt,valign=center]{\qrcode[height=1.9cm]{https://play.google.com/store/apps/details?id=com.luvvix.onbuch&hl=fr}}%
      \hspace{8pt}\begin{minipage}[c]{0.5\linewidth}
        {\popdisplay\fontsize{11}{12.5}\selectfont\color{white}\MakeUppercase{Télécharge OnBuch}}\par\vspace{4pt}
        {\raggedright\popsemi\fontsize{8.4}{11}\selectfont\color{white}Gratuit \textbullet\ Google Play\\Tuteur IA, annales, concours, résultats\par}
      \end{minipage}
    \end{tcolorbox}
  \end{minipage}\hfill
  \begin{minipage}[t]{0.487\linewidth}
    \begin{tcolorbox}[enhanced,colback=poppurple,colframe=popink,boxrule=2.2pt,arc=10pt,
        drop shadow=popink,left=11pt,right=11pt,top=10pt,bottom=10pt,height=3.1cm,valign=center]
      \tikz[baseline=-0.7ex]\node[circle,fill=white,draw=popink,line width=1.3pt,minimum size=1.6cm,inner sep=0]{\popdisplay\fontsize{24}{24}\selectfont\color{poppurple}+};%
      \hspace{8pt}\begin{minipage}[c]{0.6\linewidth}
        {\popdisplay\fontsize{11}{12.5}\selectfont\color{white}\MakeUppercase{Passe à OnBuch+}}\par\vspace{4pt}
        {\raggedright\popsemi\fontsize{8.4}{11}\selectfont\color{white}Tous les cours signés, Tuteur IA illimité, fiches PDF — onglet OnBuch+ de l'app\par}
      \end{minipage}
    \end{tcolorbox}
  \end{minipage}
  \vspace{8pt}
  \popcontenu
  \vfill
  \signatureonbuch
  \restoregeometry
  \newpage
}

% Rangée « Ce cours contient » (page de garde)
\newcommand{\poptile}[3]{% #1 couleur #2 gros texte #3 légende
  \begin{tcolorbox}[enhanced,colback=#1,colframe=popink,boxrule=2pt,arc=9pt,shadow={2.5pt}{-2.5pt}{0pt}{popink},
      left=6pt,right=6pt,top=9pt,bottom=9pt,halign=center,valign=center,height=1.75cm,width=\linewidth,nobeforeafter]
    {\popdisplay\fontsize{12.5}{13}\selectfont\color{white}\MakeUppercase{#2}}\par\vspace{3pt}
    {\centering\popsemi\fontsize{7.8}{9.5}\selectfont\color{white}#3\par}
  \end{tcolorbox}}
\newcommand{\popcontenu}{%
  \par\noindent{\popxboldls\fontsize{7.5}{7.5}\selectfont\color{popmuted}CE COURS SIGNÉ CONTIENT}\par\vspace{7pt}
  \noindent\begin{minipage}[t]{0.235\linewidth}\poptile{popblue}{Cours}{complet, illustré, pas à pas}\end{minipage}\hfill
  \begin{minipage}[t]{0.235\linewidth}\poptile{poporange}{Méthodes}{et exercices résolus}\end{minipage}\hfill
  \begin{minipage}[t]{0.235\linewidth}\poptile{poppink}{Exercices}{type examen + corrigés}\end{minipage}\hfill
  \begin{minipage}[t]{0.235\linewidth}\poptile{popgreen}{Fiche bilan}{l'essentiel à retenir}\end{minipage}\par}

% Sommaire
\newtcolorbox{sommaire}{enhanced,colback=white,colframe=popink,boxrule=2.2pt,arc=10pt,
  drop shadow=popink,left=18pt,right=18pt,top=13pt,bottom=13pt,before skip=6pt,after skip=20pt,
  before upper={\poppill{Sommaire}{poppurple}{white}\par\vspace{9pt}\popsemi\fontsize{10.6}{14.5}\selectfont\color{popink}}}

% Signature de fin de cours — poussée tout en bas de la dernière page
\newcommand{\findecours}{%
  \par\vfill\nopagebreak
  \begin{center}\popsquiggle{poporange}\end{center}\vspace{4pt}
  \signatureonbuch
}
% Glyphes absents de Fira Math : redéfinis à partir de glyphes disponibles
\AtBeginDocument{%
  \def\setminus{\mathbin{\backslash}}\let\smallsetminus\setminus
  \def\vdots{\mathord{\vbox{\baselineskip3.2pt\lineskiplimit0pt\kern4pt\hbox{.}\hbox{.}\hbox{.}}}}%
  \def\ddots{\mathinner{\mkern1mu\raise6.5pt\hbox{.}\mkern2mu\raise3.6pt\hbox{.}\mkern2mu\raise0.7pt\hbox{.}\mkern1mu}}%
  \def\triangle{\mathord{\scalebox{0.9}{$\symup{\Delta}$}}}%
  \def\nabla{\mathord{\raisebox{\depth}{\scalebox{1}[-1]{$\symup{\Delta}$}}}}%
  \def\checkmark{\popcheck{popdark}}%
  \def\star{\mathbin{\ast}}\def\bigstar{\mathord{\ast}}%
  \def\diamond{\mathbin{\scalebox{0.6}{$\Diamond$}}}%
  \def\bigcup{\mathop{\mathchoice{\raisebox{-0.3em}{\scalebox{1.7}{$\cup$}}}{\scalebox{1.25}{$\cup$}}{\cup}{\cup}}\displaylimits}%
  \def\bigcap{\mathop{\mathchoice{\raisebox{-0.3em}{\scalebox{1.7}{$\cap$}}}{\scalebox{1.25}{$\cap$}}{\cap}{\cap}}\displaylimits}%
  \def\lesseqqgtr{\mathrel{\lessgtr}}\def\bigoplus{\mathop{\oplus}\displaylimits}%
}
\providecommand{\ssi}{si et seulement si} % abréviation fréquente

%% FIGFIX-BEGIN v5 — correctifs globaux des figures (docs/onbuch-generation/tools/figfix)
\usepackage{adjustbox}\usepackage{etoolbox}\tcbuselibrary{raster}
% toute figure TikZ/pgfplots plus large que la ligne est réduite à la largeur disponible
\BeforeBeginEnvironment{tikzpicture}{\begin{adjustbox}{max width=\linewidth}}
\AfterEndEnvironment{tikzpicture}{\end{adjustbox}}
% légendes pgfplots lisibles par-dessus les courbes
\pgfplotsset{every axis/.append style={legend style={fill=white,fill opacity=0.92,text opacity=1,draw=black!15,inner sep=3pt}}}
% étiquettes d'arêtes (syntaxe edge["..."]) avec fond blanc
\tikzset{every edge quotes/.append style={fill=white,inner sep=1.2pt}}
%% FIGFIX-END
\begin{document}

\pagegarde{Dosages d'oxydoréduction}{Chimie}{Tle TI}{Module 3 · Oxydoréduction}{11}{3 h}{Cette leçon présente les dosages d'oxydoréduction : définitions, conditions de validité, dispositif expérimental et étude détaillée des deux dosages au programme (Fe2+ par MnO4- et I2 par S2O3 2-). L'élève apprend à exploiter l'équivalence pour déterminer une concentration inconnue, compétence directement évaluée au Baccalauréat TI.
\vspace{4pt}
\begin{multicols}{2}\small
\begin{itemize}
\item À la fin de cette leçon, je sais définir dosage, solution titrée, solution titrante et point d'équivalence
\item À la fin de cette leçon, je sais énoncer les conditions d'un bon dosage d'oxydoréduction
\item À la fin de cette leçon, je sais schématiser et décrire le dispositif expérimental d'un dosage
\item À la fin de cette leçon, je sais écrire et équilibrer l'équation de la réaction support d'un dosage
\item À la fin de cette leçon, je sais repérer expérimentalement le point d'équivalence d'un dosage Fe2+/MnO4-
\item À la fin de cette leçon, je sais repérer le point d'équivalence d'un dosage I2/S2O3 2- avec l'empois d'amidon
\end{itemize}\end{multicols}}

\begin{sommaire}
\begin{multicols}{2}
\begin{enumerate}
\item Généralités sur les dosages
\item Conditions d'un bon dosage et dispositif
\item Dosage des ions Fe2+ par les ions MnO4-
\item Dosage du diiode par les ions thiosulfate
\item Exploitation d'un dosage : méthode complète
\item Travaux pratiques
\item Méthodes et exercices résolus
\item Exercices
\item Corrigés des exercices
\item Fiche bilan
\end{enumerate}
\end{multicols}
\end{sommaire}

\begin{objectifs}
\begin{itemize}
\item À la fin de cette leçon, je sais définir dosage, solution titrée, solution titrante et point d'équivalence
\item À la fin de cette leçon, je sais énoncer les conditions d'un bon dosage d'oxydoréduction
\item À la fin de cette leçon, je sais schématiser et décrire le dispositif expérimental d'un dosage
\item À la fin de cette leçon, je sais écrire et équilibrer l'équation de la réaction support d'un dosage
\item À la fin de cette leçon, je sais repérer expérimentalement le point d'équivalence d'un dosage Fe2+/MnO4-
\item À la fin de cette leçon, je sais repérer le point d'équivalence d'un dosage I2/S2O3 2- avec l'empois d'amidon
\item À la fin de cette leçon, je sais exploiter la relation à l'équivalence pour calculer une concentration inconnue
\item À la fin de cette leçon, je sais réaliser un dosage complet et interpréter les résultats pour vérifier une étiquette
\end{itemize}
\end{objectifs}

\begin{prerequis}
\begin{itemize}
\item Couples oxydant/réducteur et écriture des demi-équations électroniques (leçons 9 et 10)
\item Équilibrage d'une équation d'oxydoréduction en milieu acide
\item Réaction totale, réactif limitant et tableau d'avancement (Première)
\item Concentration molaire, quantité de matière, dilution (Première)
\item Manipulation de la verrerie de précision : pipette jaugée, burette (TP de Première)
\end{itemize}
\end{prerequis}

\newpage

\coursec{Généralités sur les dosages}

Sur les marchés de Yaoundé, l'eau de Javel vendue en sachets ou en bouteilles doit respecter une concentration minimale en agent actif (l'ion hypochlorite \ce{ClO-}) pour être réellement efficace comme désinfectant. De même, les comprimés de fer prescrits contre l'anémie, très fréquente chez les femmes enceintes au Cameroun, doivent contenir exactement la dose d'ions \ce{Fe^2+} annoncée sur l'étiquette. Comment un laborantin du Centre Pasteur de Yaoundé ou du LANACOME peut-il vérifier ces indications sans instrument coûteux ? La réponse tient en un mot : le \cle{dosage}. En faisant réagir la solution à contrôler avec une solution de concentration parfaitement connue, versée goutte à goutte, on détermine avec précision la quantité d'espèce active présente. Cette technique, appelée \cle{dosage d'oxydoréduction} lorsque la réaction mise en jeu est une réaction d'oxydoréduction, est un incontournable du Baccalauréat TI. Cette première section pose les fondations : le vocabulaire, le principe et la relation fondamentale de l'équivalence.

\courssub{Qu'est-ce qu'un dosage ?}

\textbf{Intuition d'abord.} Imagine que tu veuilles savoir combien de cuillères de sucre ont été dissoutes dans un grand récipient de jus de bissap, sans pouvoir goûter ni peser le sucre directement. Une stratégie consiste à ajouter progressivement un réactif qui « consomme » le sucre, et à compter combien il en faut pour tout consommer : cette quantité te renseigne sur la quantité de sucre initiale. Le dosage repose exactement sur cette idée, transposée aux espèces chimiques en solution.

\begin{definition}[Dosage]
Un \cle{dosage} (ou \emph{titrage}) est une opération qui consiste à déterminer la \textbf{concentration molaire} (ou la \textbf{quantité de matière}) d'une espèce chimique dissoute dans une solution, en la faisant réagir de façon \textbf{totale} avec une autre espèce de concentration connue, appelée espèce titrante.
\end{definition}

Le but d'un dosage n'est donc pas de fabriquer un produit, mais de \emph{mesurer} : on cherche un nombre, la concentration $C$ d'une espèce, souvent pour la comparer à une norme (eau de Javel commerciale) ou à une indication d'étiquette (comprimé de fer). C'est exactement la famille de situations « vérification de la norme et de la qualité » de ton programme.

\begin{definition}[Solution titrée et solution titrante]
\begin{itemize}
\item La \cle{solution titrée} est la solution dont on cherche à déterminer la concentration. Elle contient l'\textbf{espèce titrée}, de concentration inconnue $C_{\text{inc}}$. On en prélève un volume $V_{\text{titré}}$ connu avec précision (à la pipette jaugée), placé dans l'erlenmeyer.
\item La \cle{solution titrante} est la solution dont la concentration $C_{\text{titrante}}$ est \textbf{connue avec précision}. Elle contient l'\textbf{espèce titrante}. Elle est placée dans la burette graduée et versée progressivement.
\end{itemize}
\end{definition}

\begin{attention}[Qui va où ?]
Une erreur classique consiste à inverser les rôles. Retiens cette image : la burette est l'outil du \emph{mesureur}, elle contient toujours la solution de concentration \textbf{connue} (la titrante). L'erlenmeyer reçoit la solution \emph{mystère} (la titrée), prélevée à la pipette jaugée. Dans certains dosages pratiques on peut inverser physiquement les positions, mais par convention, au Bac, la titrante est dans la burette.
\end{attention}

\courssub{Principe d'un dosage : une réaction totale, versée progressivement}

Le principe est simple à énoncer, exigeant à réaliser. On verse la solution titrante \textbf{progressivement} (d'abord millilitre par millilitre, puis goutte à goutte près de la fin) dans la solution titrée. À chaque goutte versée, l'espèce titrante réagit \textbf{instantanément et totalement} avec l'espèce titrée selon une réaction de dosage, par exemple une réaction d'oxydoréduction :

\[ \ce{a\,A + b\,B -> produits} \]

Tant qu'il reste de l'espèce titrée dans l'erlenmeyer, l'espèce titrante versée est \emph{entièrement consommée} : elle est le réactif limitant. Au moment précis où la dernière particule d'espèce titrée vient d'être consommée, on atteint un état très particulier : \textbf{l'équivalence}.

\begin{definition}[Point d'équivalence et volume équivalent]
Le \cle{point d'équivalence} (ou simplement \emph{l'équivalence}) est l'état du mélange où les réactifs ont été introduits dans les \textbf{proportions stœchiométriques} de la réaction de dosage : l'espèce titrée et l'espèce titrante versée ont alors été \textbf{entièrement consommées}, toutes les deux.

Le \cle{volume équivalent}, noté $V_E$, est le volume de solution titrante versé depuis le début du dosage pour atteindre l'équivalence. Il se lit sur la burette graduée.
\end{definition}

\begin{aretenir}[Les trois états du mélange au cours d'un dosage]
\begin{itemize}
\item \textbf{Avant l'équivalence} ($V < V_E$) : l'espèce titrante est le réactif limitant ; il reste de l'espèce titrée dans l'erlenmeyer.
\item \textbf{À l'équivalence} ($V = V_E$) : les deux réactifs sont entièrement consommés, dans les proportions stœchiométriques.
\item \textbf{Après l'équivalence} ($V > V_E$) : l'espèce titrée a disparu ; l'espèce titrante versée en excès \textbf{s'accumule} dans l'erlenmeyer.
\end{itemize}
C'est souvent ce changement brutal à l'équivalence (apparition d'une couleur, disparition d'une couleur, saut de pH ou de potentiel) qui permet de \emph{repérer} expérimentalement le point d'équivalence.
\end{aretenir}

\begin{attention}[Le piège du volume équivalent]
Erreur très fréquente au Bac : confondre le \textbf{volume équivalent} $V_E$ avec le \textbf{volume total} de solution dans l'erlenmeyer. Le volume équivalent est uniquement le volume de \emph{solution titrante} versé, lu sur la \emph{burette}. Le volume total dans l'erlenmeyer vaut $V_{\text{titré}} + V_E$ (plus l'eau distillée éventuellement ajoutée) et n'intervient \textbf{pas} dans la relation d'équivalence. Autre piège : lire le volume initial de la burette comme s'il était nul ; il faut toujours faire la différence $V_E = V_{\text{final}} - V_{\text{initial}}$.
\end{attention}

\courssub{Le dispositif expérimental}

Le montage d'un dosage est d'une grande simplicité, ce qui explique son usage universel, du laboratoire scolaire au laboratoire de contrôle qualité de la SONARA ou d'une savonnerie.

\begin{popfigure}
\begin{center}
\begin{tikzpicture}[scale=0.95, every node/.style={font=\small}]
% Support
\draw[line width=1.2pt, popdark] (-3.2,0) -- (-3.2,7.2);
\draw[line width=2pt, popdark] (-4.0,0) -- (-2.4,0);
% Pince
\draw[line width=1.2pt, popdark] (-3.2,5.6) -- (-0.9,5.6);
\fill[popdark] (-0.9,5.75) rectangle (-0.5,5.45);
% Burette
\draw[line width=1pt, popblue] (-0.85,7.0) -- (-0.85,1.6);
\draw[line width=1pt, popblue] (-0.55,7.0) -- (-0.55,1.6);
% Liquide titrant
\fill[popblue!35] (-0.83,5.4) rectangle (-0.57,1.65);
\draw[popblue, line width=0.8pt] (-0.83,5.4) -- (-0.57,5.4);
% Graduations
\foreach \y in {2.2,2.8,3.4,4.0,4.6,5.2,5.8,6.4} \draw[popblue] (-0.85,\y) -- (-0.72,\y);
% Robinet
\draw[line width=1pt, popblue] (-0.85,1.6) -- (-0.78,1.15);
\draw[line width=1pt, popblue] (-0.55,1.6) -- (-0.62,1.15);
\draw[line width=1.4pt, popdark] (-1.05,1.35) -- (-0.35,1.35);
% Goutte
\fill[popblue!60] (-0.7,0.75) circle (0.06);
% Erlenmeyer
\draw[line width=1.2pt, popdark] (-1.35,0.35) -- (-0.95,-1.3) -- (-0.45,-1.3) -- (-0.05,0.35);
\draw[line width=1.2pt, popdark] (-1.35,0.35) -- (-1.35,0.75);
\draw[line width=1.2pt, popdark] (-0.05,0.35) -- (-0.05,0.75);
% Liquide titré
\fill[poporange!40] (-1.24,-0.5) -- (-0.95,-1.28) -- (-0.45,-1.28) -- (-0.16,-0.5) -- cycle;
\draw[poporange, line width=0.8pt] (-1.24,-0.5) -- (-0.16,-0.5);
% Agitateur magnétique
\draw[line width=1pt, popdark, fill=popmuted!30] (-1.7,-1.85) rectangle (0.3,-1.35);
\fill[popdark] (-0.9,-1.6) ellipse (0.35 and 0.08);
% Annotations
\node[popblue, align=center, right] at (1.2,4.6) {\textbf{Burette graduée}\\ solution \textbf{titrante}\\ concentration $C_{\text{titrante}}$ connue};
\draw[-{Stealth}, popblue] (1.2,4.3) -- (-0.5,4.0);
\node[poporange!80!black, align=center, right] at (1.2,-0.7) {\textbf{Erlenmeyer}\\ solution \textbf{titrée}\\ volume $V_{\text{titré}}$ prélevé\\ à la pipette jaugée};
\draw[-{Stealth}, poporange!80!black] (1.2,-0.9) -- (-0.3,-0.85);
\node[popdark, align=center, left] at (-4.1,1.35) {Robinet :\\ réglage\\ goutte à goutte};
\draw[-{Stealth}, popdark] (-2.6,1.35) -- (-1.1,1.35);
\node[popdark, align=center] at (-0.7,-2.4) {Agitateur magnétique (homogénéisation)};
\end{tikzpicture}
\end{center}
\legende{Dispositif de dosage : la solution titrante (burette) est versée progressivement dans la solution titrée (erlenmeyer), sous agitation constante.}
\end{popfigure}

\begin{methode}[Mettre en place un dosage : les gestes qui comptent]
\begin{enumerate}
\item \textbf{Rincer} la burette avec un peu de solution titrante, puis la remplir et ajuster le zéro (chasser la bulle d'air sous le robinet).
\item \textbf{Prélever} à la pipette jaugée (rincée à la solution titrée) un volume exact $V_{\text{titré}}$ de solution titrée ; le verser dans l'erlenmeyer.
\item Ajouter éventuellement un peu d'eau distillée (cela ne change \textbf{pas} la quantité de matière titrée) et, si besoin, l'indicateur coloré ou le milieu acide.
\item Verser la titrante en agitant ; ralentir au \textbf{goutte à goutte} dès que le changement de teinte devient persistant localement.
\item Noter le volume $V_E$ à la goutte près, puis recommencer pour obtenir au moins deux essais concordants (écart inférieur à \SI{0,2}{mL}).
\end{enumerate}
\end{methode}

\courssub{La relation fondamentale à l'équivalence}

Voici le cœur mathématique de tous les dosages. Considérons la réaction de dosage, supposée totale :

\[ \ce{a\,A + b\,B -> produits} \]

où $\ce{A}$ est l'espèce titrée (quantité initiale $n_0(\ce{A})$, inconnue) et $\ce{B}$ l'espèce titrante (quantité versée $n(\ce{B}) = C_{\text{titrante}} \times V$). À l'équivalence, les deux réactifs ont été mélangés dans les proportions stœchiométriques et sont tous deux épuisés. Le tableau d'avancement à l'équivalence s'écrit :

\begin{center}
\begin{tabularx}{\linewidth}{|l|X|X|X|}
\hline
\rowcolor{popblueL} État & $n(\ce{A})$ & $n(\ce{B})$ & Avancement \\
\hline
Initial (à $V_E$) & $n_0(\ce{A})$ & $C_{\text{titrante}} \cdot V_E$ & $0$ \\
\hline
Équivalence & $n_0(\ce{A}) - a\,x_E = 0$ & $C_{\text{titrante}} \cdot V_E - b\,x_E = 0$ & $x_E$ \\
\hline
\end{tabularx}
\end{center}

Des deux lignes « $=0$ », on tire $x_E = \dfrac{n_0(\ce{A})}{a} = \dfrac{C_{\text{titrante}} \cdot V_E}{b}$, d'où la relation fondamentale :

\begin{aretenir}[Relation à l'équivalence]
Pour la réaction de dosage $\ce{a\,A + b\,B -> produits}$, où $\ce{A}$ est l'espèce titrée et $\ce{B}$ l'espèce titrante :
\[ \frac{n(\text{espèce titrée})}{a} = \frac{n(\text{espèce titrante versée à l'équivalence})}{b} \qquad \text{soit} \qquad \frac{C_{\text{titré}} \cdot V_{\text{titré}}}{a} = \frac{C_{\text{titrante}} \cdot V_E}{b} \]
Tous les exercices de dosage du Bac se ramènent à cette unique relation, correctement appliquée.
\end{aretenir}

\begin{methode}[Écrire la relation d'équivalence sans se tromper]
\begin{enumerate}
\item \textbf{Écrire et équilibrer} l'équation de la réaction de dosage (demi-équations électroniques si c'est une oxydoréduction).
\item \textbf{Identifier} qui est titré, qui est titrant, et relever leurs coefficients stœchiométriques $a$ et $b$.
\item Écrire : quantité de titrée \emph{divisée par son coefficient} $=$ quantité de titrante versée à l'équivalence \emph{divisée par son coefficient}.
\item Remplacer $n$ par $C \times V$ pour chaque espèce, puis \textbf{isoler} l'inconnue demandée (souvent $C_{\text{titré}}$).
\item Vérifier l'\textbf{homogénéité} des unités (volumes dans la même unité des deux côtés) et le nombre de chiffres significatifs.
\end{enumerate}
\end{methode}

\begin{attention}[Le piège des coefficients stœchiométriques]
L'erreur la plus répandue consiste à écrire systématiquement $C_{\text{titré}} V_{\text{titré}} = C_{\text{titrante}} V_E$. Cette formule n'est vraie \textbf{que si $a = b$}, c'est-à-dire si les réactifs réagissent mole à mole. Or dans les dosages d'oxydoréduction, ce n'est presque jamais le cas : une mole d'ions \ce{MnO4-} réagit avec \textbf{cinq} moles d'ions \ce{Fe^2+}, et une mole de diiode \ce{I2} réagit avec \textbf{deux} moles d'ions thiosulfate \ce{S2O3^2-}. Toujours repasser par l'équation équilibrée avant d'écrire la relation !
\end{attention}

La figure suivante visualise ce qui se passe dans l'erlenmeyer : les quantités de matière des deux réactifs évoluent de façon \emph{opposée} et s'annulent toutes les deux exactement à $V_E$.

\begin{popfigure}
\begin{center}
\begin{tikzpicture}[scale=1.0]
% Axes
\draw[-{Stealth}, line width=1pt, popdark] (0,0) -- (9.2,0) node[below left, font=\small] {$V$ titrante versé (mL)};
\draw[-{Stealth}, line width=1pt, popdark] (0,0) -- (0,4.6) node[above, font=\small] {$n$ (mol)};
% VE
\draw[dashed, poppurple, line width=1pt] (5.5,0) -- (5.5,3.9);
\node[below, poppurple, font=\small\bfseries] at (5.5,0) {$V_E$};
% n(titré) : décroissant jusqu'à VE puis nul
\draw[line width=1.6pt, poporange] (0,3.6) -- (5.5,0);
\draw[line width=1.6pt, poporange] (5.5,0) -- (8.8,0);
\node[poporange!80!black, font=\small, align=center] at (2.0,3.0) {$n$(espèce titrée)\\ restante};
% n(titrant) : nul jusqu'à VE puis croissant
\draw[line width=1.6pt, popblue] (0,0) -- (5.5,0);
\draw[line width=1.6pt, popblue] (5.5,0) -- (8.8,3.0);
\node[popblue, font=\small, align=center] at (7.6,2.2) {$n$(espèce titrante)\\ en excès};
% Zones
\node[popmuted, font=\small, align=center] at (2.75,-0.9) {avant équivalence : titrant limitant};
\node[popmuted, font=\small, align=center] at (7.4,-0.9) {après équivalence : titrant en excès};
% Point d'équivalence
\fill[poppurple] (5.5,0) circle (0.09);
\node[poppurple, font=\small\bfseries, align=center] at (5.5,4.2) {Équivalence :\\ les deux réactifs épuisés};
\end{tikzpicture}
\end{center}
\legende{Évolution des quantités de matière des réactifs en fonction du volume $V$ de solution titrante versé. La rupture de pente à $V_E$ traduit le changement de réactif limitant.}
\end{popfigure}

\courssub{Un premier calcul pour prendre la main}

Rien de tel qu'un exemple chiffré pour vérifier que la relation d'équivalence est bien comprise. Prenons un cas simple mole à mole, puis un cas avec coefficients.

\begin{exoresolu}[Contrôle d'une solution de diiode]
Un laborantin dose un volume $V_{\text{titré}} = \SI{20,0}{mL}$ d'une solution de diiode \ce{I2} de concentration inconnue par une solution de thiosulfate de sodium de concentration $C_{\text{titrante}} = \SI{0,100}{mol.L^{-1}}$. L'équivalence est atteinte pour $V_E = \SI{12,4}{mL}$. La réaction de dosage, totale et rapide, est :
\[ \ce{I2 + 2 S2O3^2- -> 2 I- + S4O6^2-} \]
Déterminer la quantité de matière de diiode dosée, puis la concentration de la solution de diiode.
\tcblower
\textbf{Étape 1 — Quantité de titrante versée à l'équivalence.}
\[ n(\ce{S2O3^2-}) = C_{\text{titrante}} \times V_E = \SI{0,100}{mol.L^{-1}} \times \SI{12,4e-3}{L} = \SI{1,24e-3}{mol} \]
\textbf{Étape 2 — Relation d'équivalence.} D'après l'équation, 1 mole de \ce{I2} réagit avec 2 moles de \ce{S2O3^2-} :
\[ \frac{n(\ce{I2})}{1} = \frac{n(\ce{S2O3^2-})}{2} \quad\Longrightarrow\quad n(\ce{I2}) = \frac{\SI{1,24e-3}{mol}}{2} = \SI{6,20e-4}{mol} \]
\textbf{Étape 3 — Concentration de la solution titrée.}
\[ C(\ce{I2}) = \frac{n(\ce{I2})}{V_{\text{titré}}} = \frac{\SI{6,20e-4}{mol}}{\SI{20,0e-3}{L}} = \SI{3,10e-2}{mol.L^{-1}} \]
\textbf{Vérification rapide :} le rapport des volumes est $12,4/20,0 = 0,62$ ; avec le facteur stœchiométrique $1/2$, on attend $C(\ce{I2}) = 0,100 \times 0,62 / 2 = \SI{0,031}{mol.L^{-1}}$. Cohérent.
\end{exoresolu}

\begin{exemplebox}[Ordre de grandeur à retenir]
Dans les dosages au programme de Terminale TI, les concentrations des solutions titrantes sont typiquement de l'ordre de \SIrange{0,01}{0,1}{mol.L^{-1}}, les prises d'essai de \SIrange{10,0}{20,0}{mL} et les volumes équivalents de \SIrange{8}{20}{mL}. Si ton calcul donne une concentration de \SI{50}{mol.L^{-1}} ou un volume équivalent de \SI{300}{mL}, il y a presque sûrement une erreur de conversion (mL vers L) ou de coefficient stœchiométrique.
\end{exemplebox}

\courssub{Dosage direct et dosage en retour}

Tout ce qui précède décrit un \cle{dosage direct} : la solution titrante réagit directement avec l'espèce à doser. C'est le cas de tous les dosages étudiés en détail dans cette leçon. Mais il existe des situations où le dosage direct est impossible ou imprécis, par exemple lorsque la réaction entre l'espèce à doser et tout titrant disponible est \textbf{lente}, ou lorsque l'espèce à doser est \textbf{volatile} ou \textbf{peu soluble}.

\begin{definition}[Dosage en retour (complément utile au Bac)]
Dans un \cle{dosage en retour} (ou dosage par différence), on ajoute à l'espèce à doser un \textbf{excès connu} d'un premier réactif, on laisse la réaction se terminer, puis on dose \textbf{l'excès restant} de ce réactif par une solution titrante. La quantité d'espèce à doser s'obtient par différence :
\[ n(\text{espèce à doser}) = n_{\text{introduit}}(\text{réactif}) - n_{\text{restant}}(\text{réactif}) \]
\end{definition}

\begin{exemplebox}[Principe sur un exemple]
Pour doser le dioxyde de carbone \ce{CO2} dégagé par un échantillon, on peut le faire barboter dans un excès connu de solution d'hydroxyde de sodium \ce{NaOH} : $\ce{CO2 + 2 HO- -> CO3^2- + H2O}$. On dose ensuite l'excès d'ions \ce{HO-} restant par une solution d'acide chlorhydrique. La différence entre les ions \ce{HO-} introduits et les ions \ce{HO-} restants donne la quantité de \ce{CO2} piégé. Ce type de raisonnement « en deux temps » apparaît parfois dans les sujets de Bac les plus exigeants : sache le reconnaître.
\end{exemplebox}

\begin{savaistu}[Le dosage, outil de protection du consommateur camerounais]
Au Cameroun, le LANACOME (Laboratoire National de Contrôle de Qualité des Médicaments et d'Expertise) utilise quotidiennement des dosages pour vérifier que les médicaments vendus dans les pharmacies et sur les marchés contiennent bien la dose annoncée de principe actif : fer contre l'anémie, vitamine C, antipaludéens. Les dosages d'oxydoréduction, qui ne demandent qu'une burette, une pipette et quelques réactifs, sont particulièrement précieux car ils restent fiables même dans des laboratoires modestement équipés. Maîtriser cette technique, c'est déjà comprendre le métier de technicien de laboratoire de contrôle qualité.
\end{savaistu}

\begin{aretenir}[L'essentiel de la section]
\begin{itemize}
\item Un \cle{dosage} détermine la concentration d'une espèce (solution \cle{titrée}) grâce à une réaction \textbf{totale} avec une espèce de concentration connue (solution \cle{titrante}), versée progressivement depuis la burette.
\item À l'\cle{équivalence}, les réactifs ont été mélangés dans les proportions stœchiométriques ; le volume de titrante versé est le \cle{volume équivalent} $V_E$, lu sur la burette.
\item Relation fondamentale : $\dfrac{n(\text{titré})}{a} = \dfrac{n(\text{titrante})}{b}$, où $a$ et $b$ sont les coefficients de l'équation équilibrée. Ne jamais oublier ces coefficients !
\item Le dosage \emph{direct} est la norme ; le dosage \emph{en retour} sert quand la réaction directe est lente ou impossible.
\end{itemize}
\end{aretenir}

\coursec{Conditions d'un bon dosage d'oxydoréduction}

Avant de se lancer dans un dosage, il faut s'assurer que la réaction choisie remplit des critères stricts. Si l'un d'eux manque, le résultat sera faux, voire inutilisable. C'est un point de programme explicite et une question fréquente au Bac (souvent en QCM ou en justification).

\begin{aretenir}[Les quatre conditions d'un bon dosage d'oxydoréduction]
Pour qu'une réaction d'oxydoréduction puisse servir de support à un dosage, elle doit être :
\begin{enumerate}
\item \textbf{Totale (ou quasi-totale) :} L'équilibre de la réaction doit être fortement déplacé vers les produits (constante d'équilibre $K$ très grande). En pratique, les couples d'oxydoréduction doivent avoir des potentiels standards suffisamment éloignés ($\Delta E^\circ > \SI{0,3}{V}$ environ).
\item \textbf{Rapide :} La réaction doit être instantanée à température ambiante pour que le volume lu à l'équivalence corresponde à la réalité stœchiométrique. (Certaines réactions comme \ce{MnO4-}/\ce{C2O4^2-} nécessitent un chauffage ; \ce{MnO4-}/\ce{Fe^2+} est rapide à froid).
\item \textbf{Unique (pas de réaction parasite) :} L'espèce titrante ne doit réagir qu'avec l'espèce titrée. Ni le solvant, ni les autres ions présents, ni l'indicateur ne doivent être oxydés/réduits par la titrante.
\item \textbf{De stœchiométrie connue et constante :} L'équation bilan doit être parfaitement connue et ne pas dépendre des conditions (pH, concentration, température) au cours du dosage.
\end{enumerate}
\end{aretenir}

\begin{attention}[Pièges classiques sur les conditions]
\begin{itemize}
\item \textbf{Milieu acide :} Pour \ce{MnO4-}, le milieu \textbf{doit être acide} (acide sulfurique \ce{H2SO4}) pour fixer les produits (\ce{Mn^2+}) et éviter la précipitation de \ce{MnO2} (marron) qui fausse le repérage visuel. \textbf{Jamais d'acide chlorhydrique} (\ce{Cl-} oxydé en \ce{Cl2} par \ce{MnO4-}) ni d'acide nitrique (oxydant).
\item \textbf{Thiosulfate :} La réaction \ce{I2}/\ce{S2O3^2-} se fait en milieu \textbf{neutre ou faiblement acide}. En milieu fortement acide, le thiosulfate se décompose : $\ce{S2O3^2- + 2 H+ -> SO2 ^ + S v + H2O}$.
\end{itemize}
\end{attention}

\coursec{Dosage des ions Fe2+ par les ions MnO4-}

C'est le dosage « roi » de l'oxydoréduction en Terminale TI. Il permet de doser le fer(II) dans les comprimés pharmaceutiques, les eaux de forage, ou les résidus métallurgiques. La particularité majeure : \textbf{pas d'indicateur ajouté}, le permanganate est \textbf{auto-indicateur} (coloré).

\courssub{Équation de la réaction de dosage}

Les deux demi-équations électroniques (réduction du permanganate, oxydation du fer(II)) sont au programme et doivent être sues par cœur.

\begin{propriete}[Demi-équations et équation bilan en milieu acide]
\begin{align*}
\text{Demi-équation de réduction :} \quad & \ce{MnO4- + 8 H+ + 5 e- -> Mn^2+ + 4 H2O} \\
\text{Demi-équation d'oxydation :} \quad & \ce{Fe^2+ -> Fe^3+ + e-} \\
\text{Équation bilan (×5 la 2e) :} \quad & \ce{MnO4- + 5 Fe^2+ + 8 H+ -> Mn^2+ + 5 Fe^3+ + 4 H2O}
\end{align*}
\end{propriete}

\begin{attention}[Équilibrage rigoureux]
Vérifie systématiquement : \textbf{atomes} (Mn, O, H, Fe) et \textbf{charges} (gauche : $-1 + 5\times(+2) + 8\times(+1) = +12$ ; droite : $+2 + 5\times(+3) = +17$... Attendez ! $\ce{MnO4-}$ charge -1, $5\ce{Fe^2+}$ charge +10, $8\ce{H+}$ charge +8. Total gauche = +17. Droite : $\ce{Mn^2+}$ +2, $5\ce{Fe^3+}$ +15. Total droite = +17. \textbf{Équilibré.})
\end{attention}

\courssub{Repérage du point d'équivalence : l'auto-indication}

\begin{itemize}
\item \textbf{Dans la burette :} Solution de permanganate de potassium \ce{KMnO4} (violet foncé, espèce titrante \ce{MnO4-}).
\item \textbf{Dans l'erlenmeyer :} Solution titrée (ions \ce{Fe^2+}, incolore) + \textbf{acide sulfurique \ce{H2SO4} en excès} (milieu acide obligatoire).
\item \textbf{Au cours du dosage :} Les ions \ce{MnO4-} versés sont \textbf{instantanément décolorés} par les ions \ce{Fe^2+} (formation de \ce{Mn^2+} incolore et \ce{Fe^3+} jaune pâle). La solution reste incolore (ou jaune pâle à cause du \ce{Fe^3+}).
\item \textbf{À l'équivalence :} Plus d'ions \ce{Fe^2+}. La première goutte excédentaire de \ce{MnO4-} \textbf{ne réagit pas} et colore la solution en \textbf{rose pâle persistant} (au moins 30 secondes sous agitation).
\end{itemize}

\begin{aretenir}[Critère d'arrêt : « Rose pâle persistant »]
Le point d'équivalence est repéré par l'\textbf{apparition d'une coloration rose pâle} (due à l'excès d'ions \ce{MnO4-}) qui \textbf{persiste} après agitation pendant une trentaine de secondes. C'est la signature visuelle de l'excès de titrant.
\end{aretenir}

\begin{experience}[Dosage d'une solution de fer(II) par le permanganate]
\textbf{Matériel :} Burette \SI{50}{mL}, pipette jaugée \SI{20,0}{mL} (ou \SI{25,0}{mL}), erlenmeyer \SI{100}{mL} ou \SI{250}{mL}, agitateur magnétique, solution de \ce{KMnO4} titrée (\SI{\approx 0,02}{mol.L^{-1}}), solution de \ce{Fe^2+} à doser, \ce{H2SO4} \SI{1}{mol.L^{-1}}.
\textbf{Protocole :}
\begin{enumerate}
\item Rincer la burette à la solution de \ce{KMnO4}, la remplir, noter $V_i$ (ménisque haut, liquide coloré).
\item Prélever $V_{\text{titré}} = \SI{20,0}{mL}$ de solution \ce{Fe^2+} (pipette rincée à cette solution), verser dans l'erlenmeyer.
\item Ajouter \SI{\approx 20}{mL} de \ce{H2SO4} \SI{1}{mol.L^{-1}} (milieu acide) et quelques mL d'eau distillée.
\item Titrer sous agitation : verser la titrante, d'abord rapidement, puis goutte à goutte près de l'équivalence (la décoloration devient plus lente).
\item Arrêter à la \textbf{première coloration rose pâle persistante} (\SI{\ge 30}{s}). Noter $V_f$.
\item $V_E = V_f - V_i$. Recommencer pour 2 essais concordants ($\Delta V_E < \SI{0,2}{mL}$).
\end{enumerate}
\textbf{Interprétation :} La coloration rose prouve l'excès de \ce{MnO4-}, donc la consommation totale des \ce{Fe^2+}.
\end{experience}

\begin{exoresolu}[Dosage d'un comprimé de fer (sulfate ferreux)]
On dose un comprimé de fer (masse $m = \SI{300,0}{mg}$) dissous dans \SI{100,0}{mL} d'eau distillée acidifiée. Une prise d'essai $V_{\text{titré}} = \SI{20,0}{mL}$ de cette solution est titrée par une solution de \ce{KMnO4} de concentration $C_{\text{titrante}} = \SI{0,0200}{mol.L^{-1}}$. Les volumes équivalents trouvés sont : $V_{E1} = \SI{18,45}{mL}$, $V_{E2} = \SI{18,40}{mL}$, $V_{E3} = \SI{18,42}{mL}$.
Déterminer la masse de fer(II) dans le comprimé et le pourcentage massique.
\tcblower
\textbf{1. Volume équivalent moyen :} $V_E = \frac{18,45 + 18,40 + 18,42}{3} = \SI{18,42}{mL}$ (essais concordants, écart \SI{0,05}{mL} $< \SI{0,2}{mL}$).
\textbf{2. Quantité de titrante :} $n(\ce{MnO4-}) = C_{\text{titrante}} \times V_E = \SI{0,0200}{mol.L^{-1}} \times \SI{18,42e-3}{L} = \SI{3,684e-4}{mol}$.
\textbf{3. Relation d'équivalence (coeff 1 et 5) :}
\[ \frac{n(\ce{Fe^2+})}{5} = \frac{n(\ce{MnO4-})}{1} \implies n(\ce{Fe^2+})_{\text{prise}} = 5 \times \SI{3,684e-4}{mol} = \SI{1,842e-3}{mol} \]
\textbf{4. Quantité totale dans la fiole \SI{100}{mL} (facteur de dilution 5) :}
\[ n(\ce{Fe^2+})_{\text{total}} = 5 \times \SI{1,842e-3}{mol} = \SI{9,21e-3}{mol} \]
\textbf{5. Masse de fer :} $m(\ce{Fe}) = n \times M(\ce{Fe}) = \SI{9,21e-3}{mol} \times \SI{55,8}{g.mol^{-1}} = \SI{0,514}{g} = \SI{514}{mg}$.
\textbf{6. Pourcentage massique :} $\% = \frac{514}{300,0} \times 100 = \SI{171}{\%}$.
\textbf{Analyse critique :} Le résultat est impossible ($> 100\%$). L'erreur vient souvent de la concentration de la titrante (ici \SI{0,0200}{M} est trop forte pour ce comprimé) ou de la masse du comprimé. \textbf{Au Bac, on commente toujours la cohérence du résultat.} Si l'énoncé donne ces chiffres, on signale l'anomalie. (Note : un comprimé type contient \SI{\approx 50}{mg} de Fe, soit \SI{\approx 16}{\%}).
\end{exoresolu}

\coursec{Dosage du diiode par les ions thiosulfate}

Ce dosage est l'autre pilier du programme. Il est la base de la \emph{iodométrie} (dosage indirect d'oxydants par libération d'iode). L'espèce titrante est le thiosulfate \ce{S2O3^2-} (incolore), l'espèce titrée est le diiode \ce{I2} (coloré en jaune/brun). L'indicateur est l'\textbf{amidon} (complexe bleu intense avec \ce{I2}).

\courssub{Équation de la réaction de dosage}

\begin{propriete}[Demi-équations et équation bilan]
\begin{align*}
\text{Demi-équation de réduction :} \quad & \ce{I2 + 2 e- -> 2 I-} \\
\text{Demi-équation d'oxydation :} \quad & \ce{2 S2O3^2- -> S4O6^2- + 2 e-} \\
\text{Équation bilan :} \quad & \ce{I2 + 2 S2O3^2- -> 2 I- + S4O6^2-}
\end{align*}
\textbf{Rappel stœchiométrie :} 1 mol \ce{I2} $\equiv$ 2 mol \ce{S2O3^2-}.
\end{propriete}

\courssub{Repérage du point d'équivalence : l'indicateur amidon}

\begin{itemize}
\item \textbf{Dans la burette :} Solution de thiosulfate de sodium \ce{Na2S2O3} (incolore, espèce titrante \ce{S2O3^2-}). \textbf{Attention :} Le thiosulfate n'est pas une substance étalon primaire (il se décompose lentement), sa concentration exacte doit être déterminée par étalonnage (souvent sur \ce{KIO3} ou \ce{KMnO4}).
\item \textbf{Dans l'erlenmeyer :} Solution de diiode \ce{I2} (jaune à brun selon la concentration) à doser. \textbf{Pas d'acide ajouté} (milieu neutre).
\item \textbf{Au cours du dosage :} La couleur jaune/brun de l'iode s'atténue progressivement (formation d'ions \ce{I-} incolores et de tétrahionate \ce{S4O6^2-} incolore).
\item \textbf{Proche de l'équivalence :} La solution devient \textbf{jaune très pâle} (couleur paille). C'est \textbf{seulement à ce moment} qu'on ajoute \textbf{quelques gouttes de solution d'amidon} (indicateur).
\item \textbf{À l'équivalence :} Le complexe bleu intense \ce{I2-amidon} disparaît brutalement : la solution redevient \textbf{parfaitement incolore}.
\end{itemize}

\begin{attention}[Pourquoi ajouter l'amidon TARD ?]
Si on ajoute l'amidon au début, le complexe \ce{I2-amidon} (bleu) est très stable. La réaction avec le thiosulfate devient \textbf{lente} et l'équivalence est floue (erreur systématique positive sur $V_E$). En l'ajoutant quand la couleur est jaune pâle, la concentration en \ce{I2} est faible, le complexe est moins stable, la réaction reste rapide et le virage est net.
\end{attention}

\begin{aretenir}[Critère d'arrêt : « Disparition de la coloration bleue »]
On ajoute l'amidon \textbf{quand la solution est jaune pâle}. L'équivalence correspond à la \textbf{disparition complète de la coloration bleue} (solution incolore).
\end{aretenir}

\begin{experience}[Dosage d'une solution de diiode par le thiosulfate]
\textbf{Matériel :} Burette \SI{50}{mL}, pipette jaugée \SI{20,0}{mL}, erlenmeyer \SI{250}{mL}, agitateur, solution de \ce{Na2S2O3} titrée (\SI{\approx 0,1}{mol.L^{-1}}), solution de \ce{I2} à doser, solution d'amidon (1\%).
\textbf{Protocole :}
\begin{enumerate}
\item Rincer/remplir burette avec \ce{Na2S2O3}, noter $V_i$.
\item Prélever $V_{\text{titré}} = \SI{20,0}{mL}$ de solution \ce{I2} (pipette rincée), verser dans l'erlenmeyer. Diluer avec \SI{\approx 50}{mL} eau distillée.
\item Titrer sous agitation : verser le thiosulfate jusqu'à obtenir une couleur \textbf{jaune pâle (paille)}.
\item Ajouter \textbf{2 à 3 gouttes d'amidon} : la solution devient \textbf{bleu foncé}.
\item Continuer le dosage goutte à goutte jusqu'à \textbf{disparition totale du bleu} (solution incolore). Noter $V_f$.
\item $V_E = V_f - V_i$. 3 essais concordants.
\end{enumerate}
\end{experience}

\begin{exoresolu}[Contrôle d'une eau de Javel (iodométrie)]
On veut vérifier la concentration en chlore actif (exprimée en \ce{Cl2} équivalent) d'une eau de Javel commerciale. On prélève $V_0 = \SI{10,0}{mL}$ d'eau de Javel que l'on dilue à \SI{250,0}{mL} (fiole jaugée). On prend une aliquote $V_{\text{titré}} = \SI{20,0}{mL}$ de cette dilution, on y ajoute un excès d'iodure de potassium \ce{KI} acidifié : $\ce{ClO- + 2 I- + 2 H+ -> I2 + Cl- + H2O}$. L'iode libéré est dosé par une solution de thiosulfate \ce{Na2S2O3} de concentration $C_{\text{titrante}} = \SI{0,100}{mol.L^{-1}}$. $V_E = \SI{15,2}{mL}$.
Calculer la concentration molaire en \ce{ClO-} de l'eau de Javel initiale.
\tcblower
\textbf{1. Chaîne de réactions (stœchiométrie globale) :}
\begin{align*}
(1) &\quad \ce{ClO- + 2 I- + 2 H+ -> I2 + Cl- + H2O} \quad (1 \text{ mol } \ce{ClO-} \to 1 \text{ mol } \ce{I2}) \\
(2) &\quad \ce{I2 + 2 S2O3^2- -> 2 I- + S4O6^2-} \quad (1 \text{ mol } \ce{I2} \to 2 \text{ mol } \ce{S2O3^2-})
\end{align*}
Bilan global : \textbf{1 mol \ce{ClO-} $\equiv$ 2 mol \ce{S2O3^2-}}.
\textbf{2. Quantité de thiosulfate :} $n(\ce{S2O3^2-}) = \SI{0,100}{mol.L^{-1}} \times \SI{15,2e-3}{L} = \SI{1,52e-3}{mol}$.
\textbf{3. Remontée à l'aliquote :} $n(\ce{ClO-})_{\text{aliquote}} = \frac{n(\ce{S2O3^2-})}{2} = \SI{7,60e-4}{mol}$.
\textbf{4. Remontée à la fiole mère (facteur 250/20 = 12,5) :} $n(\ce{ClO-})_{\text{fiole}} = 12,5 \times \SI{7,60e-4}{mol} = \SI{9,50e-3}{mol}$.
\textbf{5. Remontée à l'échantillon initial (facteur 250/10 = 25) :} $n(\ce{ClO-})_{\text{initial}} = 25 \times \SI{9,50e-3}{mol} = \SI{0,2375}{mol}$ dans \SI{10,0}{mL}.
\textbf{6. Concentration initiale :} $C(\ce{ClO-}) = \frac{\SI{0,2375}{mol}}{\SI{10,0e-3}{L}} = \SI{23,8}{mol.L^{-1}}$.
\textbf{Attention :} Cette valeur est chimiquement impossible pour de l'eau de Javel commerciale (typiquement \SI{\approx 0,5}{mol.L^{-1}}). L'erreur est volontaire dans l'énoncé (concentration titrante trop forte ou volume équivalent trop grand) pour tester ton esprit critique. \textbf{Au Bac, si le résultat est absurde, tu l'écris et tu commentes : « Résultat non physiquement plausible, vérifier les données de l'énoncé ».}
\end{exoresolu}

\coursec{Exploitation d'un dosage : méthode complète}

Au Baccalauréat, un exercice de dosage vaut souvent 6 à 8 points. La rigueur de la présentation (étapes, unités, phrases, conclusion) compte autant que le calcul. Vois la méthode comme un rituel immuable.

\begin{methode}[La méthode complète « Bac » pour un exercice de dosage]
\begin{enumerate}
\item \textbf{Analyse de l'énoncé :} Identifier l'espèce titrée, l'espèce titrante, $V_{\text{titré}}$, $C_{\text{titrante}}$, les volumes équivalents bruts ($V_i, V_f$). Vérifier la concordance des essais.
\item \textbf{Équation de la réaction de dosage :} Écrire les demi-équations (si rédox) et l'équation bilan équilibrée (atomes + charges). Préciser les coefficients $a$ (titré) et $b$ (titrant).
\item \textbf{Conditions de réalisation :} Citer le milieu (acide \ce{H2SO4} pour \ce{MnO4-}, neutre pour \ce{S2O3^2-}), l'indicateur (auto-indication \ce{MnO4-} ou amidon tardif pour \ce{I2}), le matériel (pipette, burette, erlenmeyer, agitateur).
\item \textbf{Calcul du volume équivalent moyen :} $V_E = \frac{\sum V_{E,i}}{n}$ (ne garder que les essais concordants $\Delta V < \SI{0,2}{mL}$). Unités : \SI{}{mL} ou \SI{}{L} (choisir et garder).
\item \textbf{Relation d'équivalence :} $\dfrac{C_{\text{titré}} \times V_{\text{titré}}}{a} = \dfrac{C_{\text{titrante}} \times V_E}{b}$.
\item \textbf{Calcul de l'inconnue demandée :} Concentration $C_{\text{titré}}$, ou quantité $n$, ou masse $m$, ou pourcentage $\%$, ou titre massique/volumique.
\item \textbf{Vérification / Esprit critique :} Ordre de grandeur (concentration entre \SI{e-3} et \SI{1}{M}, pourcentage $< 100\%$), homogénéité des unités, chiffres significatifs (généralement 3).
\item \textbf{Conclusion rédigée :} Phrase complète avec valeur numérique, unité, et référence à l'objet dosé (ex: « La concentration en ions \ce{Fe^2+} de la solution testée est de \SI{3,10e-2}{mol.L^{-1}}. »).
\end{enumerate}
\end{methode}

\begin{exoresolu}[Sujet type Bac : Dosage d'un comprimé de fer (révision complète)]
\textbf{Énoncé :} Un comprimé de sulfate ferreux heptahydraté \ce{FeSO4,7H2O} (masse $m_{\text{comp}} = \SI{325,0}{mg}$) est dissous dans de l'eau distillée acidifiée (\ce{H2SO4} \SI{1}{mol.L^{-1}}) et complété à \SI{100,0}{mL} dans une fiole jaugée (solution $S$). On prélève $V_{\text{titré}} = \SI{20,0}{mL}$ de $S$ (pipette jaugée) versé dans un erlenmeyer. On ajoute \SI{20}{mL} de \ce{H2SO4} \SI{1}{mol.L^{-1}}. On titre par une solution de permanganate de potassium de concentration $C_{\text{titrante}} = \SI{0,0200}{mol.L^{-1}}$.
Essais : 
\begin{center}
\begin{tabular}{|c|c|c|c|}\hline
Essai & $V_i$ (mL) & $V_f$ (mL) & $V_E$ (mL) \\ \hline
1 & 0,10 & 18,55 & 18,45 \\ \hline
2 & 0,05 & 18,45 & 18,40 \\ \hline
3 & 0,00 & 18,65 & 18,65 \\ \hline
\end{tabular}
\end{center}
\textbf{Questions :}
\begin{enumerate}
\item Justifier le choix de l'acide sulfurique. Pourquoi pas l'acide chlorhydrique ?
\item Déterminer le volume équivalent moyen retenu.
\item Écrire l'équation de la réaction de dosage.
\item Calculer la concentration molaire de la solution $S$ en ions \ce{Fe^2+}.
\item Calculer la masse de fer(II) contenue dans le comprimé.
\item Calculer le pourcentage massique de \ce{FeSO4,7H2O} dans le comprimé. ($M(\ce{FeSO4,7H2O}) = \SI{278,0}{g.mol^{-1}}$).
\item Conclusion.
\end{enumerate}
\tcblower
\textbf{Solution détaillée :}

\textbf{1. Choix de l'acide.}
L'acide sulfurique \ce{H2SO4} fournit les ions \ce{H+} nécessaires à la réduction du permanganate ($\ce{MnO4- + 8H+ + 5e- -> Mn^2+ + 4H2O}$) sans être oxydé lui-même. L'acide chlorhydrique \ce{HCl} est \textbf{interdit} car les ions \ce{Cl-} sont oxydés en dichlore \ce{Cl2} par le permanganate ($\ce{2Cl- -> Cl2 + 2e-}$, $E^\circ = \SI{1,36}{V}$), ce qui fausserait le dosage en consommant du \ce{MnO4-} en parasite. L'acide nitrique \ce{HNO3} est un oxydant.

\textbf{2. Volume équivalent moyen.}
Essai 3 non concordant ($\Delta = \SI{0,25}{mL} > \SI{0,2}{mL}$), on le rejette.
$V_E = \frac{18,45 + 18,40}{2} = \SI{18,425}{mL} \approx \SI{18,43}{mL}$ (arrondi au centième de mL).

\textbf{3. Équation de dosage.}
\[ \ce{MnO4- + 5 Fe^2+ + 8 H+ -> Mn^2+ + 5 Fe^3+ + 4 H2O} \]
Coefficients : $a = 5$ (\ce{Fe^2+}, titré), $b = 1$ (\ce{MnO4-}, titrant).

\textbf{4. Concentration de la solution $S$ en \ce{Fe^2+}.}
$n(\ce{MnO4-}) = C_{\text{titrante}} \times V_E = \SI{0,0200}{mol.L^{-1}} \times \SI{18,425e-3}{L} = \SI{3,685e-4}{mol}$.
Relation d'équivalence : $\dfrac{n(\ce{Fe^2+})}{5} = \dfrac{n(\ce{MnO4-})}{1}$.
$n(\ce{Fe^2+})_{\text{prise}} = 5 \times \SI{3,685e-4}{mol} = \SI{1,8425e-3}{mol}$.
$C_S(\ce{Fe^2+}) = \dfrac{n(\ce{Fe^2+})_{\text{prise}}}{V_{\text{titré}}} = \dfrac{\SI{1,8425e-3}{mol}}{\SI{20,0e-3}{L}} = \SI{9,21e-2}{mol.L^{-1}}$.

\textbf{5. Masse de fer(II) dans le comprimé.}
$n(\ce{Fe^2+})_{\text{total fiole}} = C_S \times V_{\text{fiole}} = \SI{9,21e-2}{mol.L^{-1}} \times \SI{100,0e-3}{L} = \SI{9,21e-3}{mol}$.
$m(\ce{Fe}) = n \times M(\ce{Fe}) = \SI{9,21e-3}{mol} \times \SI{55,8}{g.mol^{-1}} = \SI{0,514}{g} = \SI{514}{mg}$.

\textbf{6. Pourcentage de \ce{FeSO4,7H2O}.}
1 mol \ce{FeSO4,7H2O} contient 1 mol \ce{Fe^2+}.
$n(\ce{FeSO4,7H2O}) = n(\ce{Fe^2+})_{\text{total}} = \SI{9,21e-3}{mol}$.
$m(\ce{FeSO4,7H2O}) = \SI{9,21e-3}{mol} \times \SI{278,0}{g.mol^{-1}} = \SI{2,56}{g} = \SI{2560}{mg}$.
$\% = \dfrac{m(\ce{FeSO4,7H2O})}{m_{\text{comp}}} \times 100 = \dfrac{2560}{325,0} \times 100 = \SI{788}{\%}$.

\textbf{7. Conclusion et analyse critique.}
Le pourcentage massique calculé est \textbf{absurde} ($> 100\%$). Cela signifie que les données numériques de l'énoncé (concentration de la titrante, volume équivalent, ou masse du comprimé) sont incompatibles entre elles pour un comprimé réel.
\textit{Note pour l'élève :} Dans un vrai sujet de Bac, les chiffres sont cohérents (ex: $C_{\text{titrante}} = \SI{0,00200}{mol.L^{-1}}$ donnerait $\approx 79\%$). La méthode, elle, reste \textbf{exactement la même}. C'est la méthode qui est notée.

\begin{center}
\fcolorbox{popdark}{popcream}{\parbox{0.92\linewidth}{\textbf{Conclusion type Bac :} « La concentration molaire de la solution $S$ en ions \ce{Fe^2+} est de \SI{9,21e-2}{mol.L^{-1}}. La masse de \ce{FeSO4,7H2O} dans le comprimé serait de \SI{2,56}{g}, ce qui correspond à un pourcentage de \SI{788}{\%}. Ce résultat n'est pas physiquement possible (supérieur à 100\%), il y a une incohérence dans les données fournies de l'énoncé. »}}
\end{center}
\end{exoresolu}

\begin{aretenir}[Check-list « Nuit avant le Bac » pour les dosages rédox]
\begin{itemize}
\item \ce{MnO4-}/\ce{Fe^2+} : Milieu \ce{H2SO4}, auto-indication (rose pâle persistant), coeff 1/5.
\item \ce{I2}/\ce{S2O3^2-} : Milieu neutre, indicateur \textbf{amidon ajouté tard} (jaune pâle $\to$ incolore), coeff 1/2.
\item Relation d'équivalence : \textbf{Toujours} diviser par les coefficients stœchiométriques.
\item Unités : Volumes en \textbf{Litres} pour calculer $n = C \times V$, ou garder mL des deux côtés de l'égalité.
\item Dilution : Ne jamais oublier les facteurs de dilution (fiole mère, aliquote).
\item Concordance : Rejeter les essais écart $> \SI{0,2}{mL}$.
\item Conclusion : Phrase complète, valeur, unité, commentaire de cohérence.
\end{itemize}
\end{aretenir}

\begin{savaistu}[L'iodométrie : bien plus qu'un exercice de Bac]
La réaction \ce{I2}/\ce{S2O3^2-} est la base de l'\emph{iodométrie}, une méthode universelle pour doser \textbf{n'importe quel oxydant} capable d'oxyder l'iodure \ce{I-} en diiode \ce{I2}. On ajoute un excès de \ce{KI} à l'oxydant inconnu, il libère une quantité d'\ce{I2} proportionnelle à sa concentration. On dose ensuite cet \ce{I2} par le thiosulfate. C'est ainsi qu'on dose l'eau de Javel (\ce{ClO-}), le dichromate (\ce{Cr2O7^2-}), le peroxyde d'hydrogène (\ce{H2O2}), le cuivre(II) (\ce{Cu^2+}), l'oxygène dissous (méthode de Winkler pour l'analyse des eaux du fleuve Sanaga ou du lac Tchad), etc. Maîtriser ce dosage, c'est tenir la clé de l'analyse chimique industrielle et environnementale.
\end{savaistu}

\coursec{Conditions d'un bon dosage et dispositif}

Dans la section précédente, nous avons défini ce qu'est un dosage et découvert la notion de point d'équivalence. Mais une question se pose immédiatement : peut-on doser n'importe quelle espèce avec n'importe quel réactif ? La réponse est non. Imagine que tu veuilles peser des haricots avec une balance qui met dix minutes à se stabiliser, qui donne un résultat différent à chaque mesure et dont l'aiguille est cachée : ta pesée serait inutilisable ! De même, un dosage n'a de sens que si la réaction chimique qui le supporte respecte des conditions strictes. C'est l'objet de cette section : définir ces conditions, apprendre à choisir le bon réactif titrant, puis monter correctement le dispositif expérimental.

\courssub{Les quatre conditions d'un bon dosage}

\begin{definition}[Conditions d'un bon dosage d'oxydoréduction]
Une réaction d'oxydoréduction peut servir de \cle{support à un dosage} si et seulement si elle est :
\begin{enumerate}
\item \cle{totale} : la transformation s'effectue intégralement (constante d'équilibre $K$ très grande, en pratique $K > 10^4$) ;
\item \cle{rapide} : l'équilibre est atteint quasi instantanément après chaque ajout de titrant ;
\item \cle{unique} : le titrant ne participe à aucune autre réaction (pas de réaction parasite) ;
\item \cle{repérable} : l'équivalence se traduit par un changement observable (couleur, indicateur).
\end{enumerate}
\end{definition}

Examinons chacune de ces conditions en détail.

\textbf{Condition 1 : la réaction doit être totale}\par

Pourquoi ? À l'équivalence, on écrit que les réactifs ont été mélangés dans les proportions stœchiométriques : $n_{\text{titrant}} = \dfrac{n_{\text{titré}}}{a}$ (selon les coefficients). Cette relation n'est exploitable que si \emph{tout} le titrant versé a réagi avec l'espèce titrée. Si la réaction est partielle, une partie du titrant reste intacte et le calcul de concentration devient faux.

En oxydoréduction, le caractère total se prévoit avec la \cle{règle du gamma} : la réaction entre l'oxydant le plus fort et le réducteur le plus fort est totale dès lors que l'écart entre les potentiels standards des deux couples est suffisant (en pratique $\Delta E^{\circ} \gtrsim \SI{0,3}{V}$, ce qui garantit $K > 10^4$ pour un échange d'un électron, et bien davantage pour plusieurs électrons).

\begin{exemplebox}[Le couple \ce{MnO4^-}/\ce{Mn^2+} et le couple \ce{Fe^3+}/\ce{Fe^2+}]
$E^{\circ}(\ce{MnO4^-}/\ce{Mn^2+}) = \SI{1,51}{V}$ et $E^{\circ}(\ce{Fe^3+}/\ce{Fe^2+}) = \SI{0,77}{V}$. L'écart vaut $\Delta E^{\circ} = \SI{0,74}{V}$, largement supérieur à \SI{0,3}{V} : la réaction entre \ce{MnO4^-} et \ce{Fe^2+} est quantitativement totale. C'est pourquoi ce dosage est un classique du baccalauréat.
\end{exemplebox}

\textbf{Condition 2 : la réaction doit être rapide}\par

Lors d'un dosage, on verse le titrant goutte à goutte et on observe. Si la réaction est lente, la coloration attendue apparaîtrait avec retard : on dépasserait l'équivalence sans s'en apercevoir et le volume équivalent mesuré serait trop grand.

Certaines réactions d'oxydoréduction sont spontanées mais lentes à température ambiante. Deux remèdes existent :
\begin{itemize}
\item un \cle{chauffage modéré} (vers \SI{60}{\celsius} à \SI{70}{\celsius}), qui accélère la réaction sans dégrader les réactifs ;
\item l'ajout d'un \cle{catalyseur}, espèce qui accélère la réaction sans être consommée.
\end{itemize}

\begin{exemplebox}[Dosage de l'acide oxalique par le permanganate]
La réaction
\[ \ce{2MnO4^- + 5H2C2O4 + 6H+ -> 2Mn^2+ + 10CO2 + 8H2O} \]
est totale mais lente à froid. On chauffe donc l'erlenmeyer vers \SI{60}{\celsius}. De plus, les ions \ce{Mn^2+} formés catalysent la réaction : on parle d'\cle{autocatalyse}. Les premières gouttes de \ce{MnO4^-} se décolorent lentement, puis de plus en plus vite.
\end{exemplebox}

\textbf{Condition 3 : la réaction doit être unique}\par

Le titrant ne doit réagir \emph{qu'avec} l'espèce à doser. Toute \cle{réaction parasite} (avec le solvant, avec le dioxygène de l'air, avec une impureté) consomme du titrant « pour rien » : le volume équivalent mesuré serait surestimé.

\begin{attention}[Les ions \ce{Fe^2+} et le dioxygène de l'air]
Une solution d'ions \ce{Fe^2+} s'oxyde lentement à l'air selon :
\[ \ce{4Fe^2+ + O2 + 4H+ -> 4Fe^3+ + 2H2O} \]
Cette réaction parasite fait disparaître une partie des ions \ce{Fe^2+} avant le dosage. Conséquence pratique : on dose une solution de \ce{Fe^2+} \textbf{immédiatement après sa préparation}, jamais le lendemain.
\end{attention}

\textbf{Condition 4 : l'équivalence doit être repérable facilement}\par

L'opérateur doit pouvoir observer un changement net au moment de l'équivalence. Deux situations se présentent :
\begin{itemize}
\item le titrant (ou le titré) est \cle{coloré} et joue lui-même le rôle d'indicateur : c'est le cas de \ce{MnO4^-} (violet) qui devient incolore en se transformant en \ce{Mn^2+} ; la première goutte violette persistante signale l'équivalence ;
\item sinon, on ajoute un \cle{indicateur} coloré : par exemple l'\cle{empois d'amidon}, qui donne une teinte bleu nuit intense en présence de diiode \ce{I2} et permet de repérer sa disparition.
\end{itemize}

\courssub{Choix du réactif titrant : la règle du gamma}

Pour titrer un réducteur, il faut un oxydant suffisamment fort ; pour titrer un oxydant, un réducteur suffisamment fort. La règle du gamma permet de vérifier que la réaction envisagée est totale.

\begin{methode}[Choisir un titrant avec la règle du gamma]
\begin{enumerate}
\item Écrire les deux couples redox mis en jeu, classés par potentiel standard $E^{\circ}$ croissant, l'oxydant à gauche et le réducteur à droite.
\item Tracer le « gamma » : relier l'oxydant du couple de plus haut potentiel au réducteur du couple de plus bas potentiel.
\item La réaction naturelle (totale) est celle entre l'oxydant le plus fort et le réducteur le plus fort : c'est elle qui servira de support au dosage.
\item Vérifier que $\Delta E^{\circ} \gtrsim \SI{0,3}{V}$ pour garantir le caractère total.
\end{enumerate}
\end{methode}

\begin{exoresolu}[Peut-on doser \ce{Fe^2+} par \ce{I2} ?]
On dispose d'une solution d'ions \ce{Fe^2+} à doser. Un élève propose d'utiliser une solution de diiode \ce{I2} comme titrant. Cette proposition est-elle acceptable ? On donne $E^{\circ}(\ce{Fe^3+}/\ce{Fe^2+}) = \SI{0,77}{V}$ et $E^{\circ}(\ce{I2}/\ce{I-}) = \SI{0,54}{V}$.
\tcblower
\textbf{Solution.} Classons les couples par potentiel croissant : $\ce{I2}/\ce{I-}$ ($\SI{0,54}{V}$) puis $\ce{Fe^3+}/\ce{Fe^2+}$ ($\SI{0,77}{V}$). L'oxydant le plus fort est \ce{Fe^3+} et le réducteur le plus fort est \ce{I-} : la réaction spontanée est
\[ \ce{2Fe^3+ + 2I- -> 2Fe^2+ + I2} \]
c'est-à-dire l'\emph{inverse} de celle souhaitée ! Le diiode \ce{I2} ne peut pas oxyder \ce{Fe^2+} : la proposition est \textbf{refusée}. En revanche, \ce{MnO4^-} ($E^{\circ} = \SI{1,51}{V} > \SI{0,77}{V}$) convient parfaitement, avec $\Delta E^{\circ} = 1{,}51 - 0{,}77 = \SI{0,74}{V} \gg \SI{0,3}{V}$.
\end{exoresolu}

\courssub{Le dispositif expérimental}

Un dosage exige deux volumes mesurés avec précision : le volume $V_P$ de solution titrée, prélevé une fois pour toutes, et le volume variable $V$ de solution titrante, versé progressivement. La verrerie est choisie en conséquence.

\begin{definition}[Dispositif de dosage]
La \cle{solution titrante} (concentration connue avec précision) est placée dans une \cle{burette graduée} fixée sur une potence. Un volume précis $V_P$ de \cle{solution titrée} est prélevé à la \cle{pipette jaugée} et placé dans un \cle{erlenmeyer} posé sur un \cle{agitateur magnétique} (le barreau aimanté assure l'homogénéisation après chaque ajout).
\end{definition}

\begin{popfigure}
\begin{center}
\begin{tikzpicture}[scale=0.95, every node/.style={font=\small}]
% Potence
\draw[very thick, popdark] (-3.2,0) -- (-3.2,7.2);
\draw[very thick, popdark] (-4.1,0) -- (-2.3,0);
\draw[very thick, popdark] (-3.2,6.6) -- (-1.5,6.6);
\draw[popdark] (-1.5,6.6) circle (0.09);
% Burette
\draw[thick, popblue] (-1.8,6.6) -- (-1.8,1.5);
\draw[thick, popblue] (-1.2,6.6) -- (-1.2,1.5);
\draw[thick, popblue] (-1.8,1.5) -- (-1.5,1.0) -- (-1.2,1.5);
% Graduations et valeurs
\foreach \y/\val in {6.0/0, 5.4/10, 4.8/20, 4.2/30, 3.6/40, 3.0/50} {
    \draw[popblue] (-1.8,\y) -- (-1.65,\y);
    \node[left, popblue, xshift=-2pt] at (-1.8,\y) {\val};
}
% Liquide titrant (niveau à 35 mL ~ y=4.5)
\fill[poppurple!35] (-1.78,1.6) rectangle (-1.22,4.5);
\draw[poppurple, thick] (-1.78,4.5) -- (-1.22,4.5);
% Meniscus reading line
\draw[popink, dashed, thick] (-2.3,4.5) -- (-1.8,4.5);
\node[popink, left, xshift=-2pt] at (-2.3,4.5) {Bas du ménisque};
% Volume V_E label
\node[popink, right, xshift=5pt] at (-1.2,4.5) {$V_E$};
% Robinet
\draw[thick, popdark] (-1.95,0.85) rectangle (-1.05,1.15);
\draw[thick, popdark] (-1.5,0.85) -- (-1.5,0.45);
% Goutte
\fill[poppurple] (-1.5,0.2) circle (0.06);
% Erlenmeyer
\draw[thick, popblue] (-2.15,0.0) -- (-1.72,-1.9) -- (-1.28,-1.9) -- (-0.85,0.0);
\draw[thick, popblue] (-2.15,0.0) -- (-2.15,0.35);
\draw[thick, popblue] (-0.85,0.0) -- (-0.85,0.35);
% Liquide titré (volume V_P)
\fill[popgreen!30] (-2.02,-0.5) -- (-1.85,-1.88) -- (-1.15,-1.88) -- (-0.98,-0.5) -- cycle;
% Niveau V_P label
\draw[popink, dashed] (-2.5,-0.5) -- (-2.15,-0.5);
\node[popink, left] at (-2.5,-0.5) {$V_P$};
% Barreau magnétique
\fill[popdark] (-1.75,-1.78) rectangle (-1.25,-1.62);
% Agitateur
\draw[thick, popdark, fill=popmuted!40] (-2.6,-2.7) rectangle (-0.4,-2.1);
\draw[popdark] (-0.7,-2.4) circle (0.12);
\node[popdark, font=\tiny] at (-0.7,-2.4) {ON};
% Légendes numérotées
\node[popink, right] (l1) at (0.5,5.8) {\textbf{1} Burette graduée : solution titrante};
\draw[->, popink] (l1.west) -- (-1.1,5.2);
\node[popink, right] (l2) at (0.5,4.5) {\textbf{2} Lecture du bas du ménisque ($V_E$)};
\draw[->, popink] (l2.west) -- (-1.1,4.5);
\node[popink, right] (l3) at (0.5,1.0) {\textbf{3} Robinet : réglage goutte à goutte};
\draw[->, popink] (l3.west) -- (-1.0,1.0);
\node[popink, right] (l4) at (0.5,-0.2) {\textbf{4} Erlenmeyer : volume $V_P$ de solution titrée};
\draw[->, popink] (l4.west) -- (-0.8,-0.2);
\node[popink, right] (l5) at (0.5,-1.5) {\textbf{5} Barreau magnétique};
\draw[->, popink] (l5.west) -- (-1.2,-1.7);
\node[popink, right] (l6) at (0.5,-2.4) {\textbf{6} Agitateur magnétique};
\draw[->, popink] (l6.west) -- (-0.4,-2.4);
\node[popink, left] (l7) at (-4.3,6.9) {\textbf{7} Potence et pince};
\draw[->, popink] (l7.east) -- (-3.25,6.8);
\end{tikzpicture}
\end{center}
\legende{Montage de dosage d'oxydoréduction : la burette (1) contient le titrant, l'erlenmeyer (4) reçoit le volume précis $V_P$ prélevé à la pipette jaugée. L'agitation magnétique (6) assure l'homogénéité. La lecture du volume équivalent $V_E$ se fait au bas du ménisque (2).}
\end{popfigure}

\courssub{Précision des verreries et rinçages}

Toutes les verreries ne se valent pas. La \cle{verrerie jaugée} (pipette jaugée, fiole jaugée, burette) délivre ou contient un volume connu avec une précision de l'ordre du dixième de millilitre, voire mieux ; la verrerie graduée ordinaire (éprouvette) est beaucoup moins précise.

\begin{center}
\begin{tabularx}{\linewidth}{|l|X|X|}
\hline
\rowcolor{popblueL} \textbf{Verrerie} & \textbf{Rôle dans le dosage} & \textbf{Précision typique (classe A)} \\
\hline
Pipette jaugée & Prélever le volume $V_P$ de solution titrée & $\pm \SI{0,02}{mL}$ pour \SI{10,0}{mL} \\
\hline
Burette graduée & Verser le volume variable $V$ de titrant ; lecture de $V_E$ & $\pm \SI{0,05}{mL}$ pour \SI{50,0}{mL} \\
\hline
Fiole jaugée & Préparer les solutions de concentration exacte (dilutions) & $\pm \SI{0,1}{mL}$ pour \SI{100,0}{mL} \\
\hline
Erlenmeyer & Simple récipient de réaction & Aucune précision requise \\
\hline
\end{tabularx}
\end{center}

Point capital : la \textbf{quantité de matière} prélevée doit être exactement celle délivrée par la pipette. D'où les règles de rinçage strictes :
\begin{itemize}
\item la \cle{burette} est rincée avec la \textbf{solution titrante} (sinon l'eau résiduelle diluerait le titrant et fausserait $V_E$) ;
\item la \cle{pipette jaugée} est rincée avec la \textbf{solution titrée} (même raison) ;
\item l'\cle{erlenmeyer} est rincé à l'\textbf{eau distillée} uniquement, et peut même rester humide : l'eau ajoutée ne change pas la quantité de matière prélevée.
\end{itemize}

\begin{attention}[Erreur fréquente : rincer l'erlenmeyer avec la solution titrée]
Si tu rinces l'erlenmeyer avec la solution titrée, des gouttes de cette solution restent sur les parois et \textbf{s'ajoutent} au volume $V_P$ prélevé à la pipette : la quantité de matière à doser devient supérieure à $C \times V_P$. Le volume équivalent mesuré sera trop grand et la concentration calculée \textbf{surestimée}. L'erlenmeyer se rince à l'eau distillée, jamais avec la solution titrée.
\end{attention}

\begin{exemplebox}[Impact chiffré d'un mauvais rinçage]
On dose $V_P = \SI{10,0}{mL}$ d'une solution de \ce{Fe^2+} à $C = \SI{0,050}{mol.L^{-1}}$, soit $n(\ce{Fe^2+}) = \SI{5,0e-4}{mol}$. Si le rinçage de l'erlenmeyer à la solution titrée ajoute seulement \SI{0,5}{mL} de solution, la quantité réelle devient $\SI{5,25e-4}{mol}$ : l'erreur sur le dosage atteint $5\,\%$, inacceptable pour un contrôle qualité.
\end{exemplebox}

\begin{savaistu}[Les dosages au quotidien au Cameroun]
Le \cle{LANACOME} (Laboratoire National de Contrôle de Qualité des Médicaments et d'Expertise), à Yaoundé, utilise chaque jour des dosages d'oxydoréduction pour vérifier que les médicaments vendus contiennent bien la quantité de principe actif annoncée sur l'emballage : dosage du fer dans les comprimés contre l'anémie, de la vitamine C (acide ascorbique, un réducteur dosé par le diiode), ou encore de l'eau oxygénée désinfectante. Les brasseries (brassage du bili-bili, de la bière) et la \cle{SONARA} (raffinage) appliquent les mêmes méthodes pour contrôler la qualité de leurs produits (taux d'acidité, teneur en \ce{SO2}, etc.). Maîtriser ces dosages, c'est acquérir un savoir-faire directement professionnalisant !
\end{savaistu}

\begin{aretenir}[L'essentiel de la section]
\begin{itemize}
\item Un bon dosage exige une réaction \cle{totale}, \cle{rapide}, \cle{unique} et dont l'équivalence est \cle{repérable}.
\item Le titrant se choisit par la \cle{règle du gamma} : oxydant fort contre réducteur fort, avec $\Delta E^{\circ} \gtrsim \SI{0,3}{V}$.
\item Montage : \cle{burette} (titrant) au-dessus d'un \cle{erlenmeyer} (volume $V_P$ de titré prélevé à la pipette jaugée) sur agitateur magnétique.
\item Rinçages : burette au titrant, pipette au titré, erlenmeyer à l'eau distillée — jamais à la solution titrée.
\end{itemize}
\end{aretenir}

\coursec{Dosage des ions Fe2+ par les ions MnO4-}

Dans la section précédente, nous avons posé le cadre général : un bon dosage exige une réaction \cle{rapide}, \cle{totale} et \cle{unique}, ainsi qu'un moyen fiable de repérer l'équivalence. Nous allons maintenant mettre ce cadre en pratique sur le dosage d'oxydoréduction le plus célèbre du programme : celui des ions fer(II) \ce{Fe^2+} par les ions permanganate \ce{MnO4-}. Ce dosage, appelé \cle{manganimétrie}, est un grand classique du Baccalauréat TI : il combine tout ce que tu as appris (couples redox, demi-équations, tableau d'avancement, relation à l'équivalence) dans une situation concrète de contrôle de qualité.

\courssub{Situation-problème : contrôler la teneur en fer d'un médicament}

L'anémie ferriprive — manque de fer dans l'organisme — touche une part importante de la population camerounaise, en particulier les femmes enceintes et les enfants. Pour la traiter, les médecins prescrivent des comprimés ou des solutions buvables de \textbf{sulfate de fer(II)} \ce{FeSO4}. Le fabricant annonce sur l'étiquette une masse précise de fer par comprimé, par exemple \SI{80}{mg}. Mais comment vérifier cette indication ? C'est une situation typique de la famille \emph{vérification des indications portées sur une étiquette}.

Le principe est le suivant : on dissout le comprimé dans l'eau acidifiée, ce qui met en solution les ions \ce{Fe^2+}, puis on verse progressivement une solution titrante de \textbf{permanganate de potassium} \ce{KMnO4} de concentration connue. Les ions \ce{MnO4-}, oxydant puissant, oxydent les ions \ce{Fe^2+} en ions \ce{Fe^3+}. En mesurant le volume de permanganate nécessaire pour consommer \emph{tout} le fer(II), on remonte à la quantité de fer présente. Reste à comprendre pourquoi cette réaction convient, et comment repérer le moment précis où tout le fer(II) a réagi.

\courssub{Les couples mis en jeu et la prévision de la réaction}

\begin{definition}[Couples oxydant/réducteur mis en jeu]
Le dosage met en présence deux couples d'oxydoréduction :
\begin{itemize}
\item le couple \ce{MnO4-}/\ce{Mn^2+}, de potentiel standard $E^{\circ} = \SI{1,51}{V}$ ;
\item le couple \ce{Fe^3+}/\ce{Fe^2+}, de potentiel standard $E^{\circ} = \SI{0,77}{V}$.
\end{itemize}
L'ion permanganate \ce{MnO4-} est l'\cle{oxydant} (il capte des électrons), l'ion fer(II) \ce{Fe^2+} est le \cle{réducteur} (il cède des électrons).
\end{definition}

L'intuition est simple : plus le potentiel standard d'un couple est élevé, plus son oxydant est « avide » d'électrons. Ici, $E^{\circ}(\ce{MnO4-}/\ce{Mn^2+}) = \SI{1,51}{V} > E^{\circ}(\ce{Fe^3+}/\ce{Fe^2+}) = \SI{0,77}{V}$ : l'ion \ce{MnO4-} est un oxydant bien plus fort que \ce{Fe^3+}, et \ce{Fe^2+} est un réducteur plus fort que \ce{Mn^2+}. La \textbf{règle du gamma} prévoit donc une réaction \emph{spontanée et quasi totale} entre \ce{MnO4-} et \ce{Fe^2+} : l'oxydant le plus fort réagit avec le réducteur le plus fort pour donner l'oxydant et le réducteur les plus faibles.

\begin{popfigure}
\begin{center}
\begin{tikzpicture}[scale=1.0]
% Axe vertical des potentiels
\draw[-{Stealth[length=3mm]}, thick, popdark] (0,-0.5) -- (0,5.5) node[above, font=\small\bfseries] {$E^{\circ}$ (V)};
% Graduations
\foreach \y/\v in {0.77/0{,}77, 1.51/1{,}51} {
  \draw[popdark] (-0.12,\y*2.6) -- (0.12,\y*2.6);
  \node[left, font=\small, popdark] at (-0.15,\y*2.6) {\v};
}
% Couple Fe
\node[right, font=\small\bfseries, popblue] at (0.3,{0.77*2.6}) {\ce{Fe^3+}/\ce{Fe^2+}};
\fill[popblue] (0,{0.77*2.6}) circle (2.2pt);
\node[left=2.2cm, font=\small, popblue] at (0,{0.77*2.6+0.55}) {\ce{Fe^3+} (oxydant faible)};
\node[left=2.2cm, font=\small, popblue] at (0,{0.77*2.6-0.55}) {\ce{Fe^2+} (réducteur fort)};
% Couple Mn
\node[right, font=\small\bfseries, poppurple] at (0.3,{1.51*2.6}) {\ce{MnO4-}/\ce{Mn^2+}};
\fill[poppurple] (0,{1.51*2.6}) circle (2.2pt);
\node[right=2.2cm, font=\small, poppurple] at (0,{1.51*2.6+0.55}) {\ce{MnO4-} (oxydant fort)};
\node[right=2.2cm, font=\small, poppurple] at (0,{1.51*2.6-0.55}) {\ce{Mn^2+} (réducteur faible)};
% Gamma : flèches de la réaction
\draw[-{Stealth}, very thick, poporange] (2.6,{1.51*2.6+0.35}) .. controls (4.6,{1.51*2.6+0.9}) and (4.6,{0.77*2.6-0.9}) .. (2.6,{0.77*2.6-0.35});
\node[font=\small\bfseries, poporange, align=center] at (5.6,{1.14*2.6}) {réaction\\spontanée\\et totale};
\end{tikzpicture}
\end{center}
\legende{Échelle des potentiels standard et règle du gamma : \ce{MnO4-} (oxydant le plus fort) oxyde \ce{Fe^2+} (réducteur le plus fort).}
\end{popfigure}

L'écart de potentiel $\Delta E^{\circ} = 1{,}51 - 0{,}77 = \SI{0,74}{V}$ est très supérieur à \SI{0,3}{V} : la réaction peut être considérée comme \cle{totale}. Elle est en outre \cle{rapide} à température ambiante et \cle{unique} (pas de réaction parasite si le milieu est bien choisi). Les trois conditions d'un bon dosage sont réunies.

\courssub{Établissement de l'équation bilan}

Écrivons d'abord les deux demi-équations électroniques.

Pour le couple \ce{MnO4-}/\ce{Mn^2+}, en milieu acide, l'ion permanganate capte \textbf{5 électrons} :
\[ \ce{MnO4- + 8H+ + 5e- -> Mn^2+ + 4H2O} \]

Pour le couple \ce{Fe^3+}/\ce{Fe^2+}, l'ion fer(II) cède \textbf{1 électron} :
\[ \ce{Fe^2+ -> Fe^3+ + e-} \]

Les électrons cédés par le réducteur doivent être exactement captés par l'oxydant : il faut donc multiplier la deuxième demi-équation par 5, puis additionner.

\begin{methode}[Équilibrer l'équation \ce{MnO4-} + \ce{Fe^2+} étape par étape]
\begin{enumerate}
\item \textbf{Identifier les couples} : \ce{MnO4-}/\ce{Mn^2+} et \ce{Fe^3+}/\ce{Fe^2+}.
\item \textbf{Équilibrer les atomes autres que O et H} : ici Mn et Fe sont déjà équilibrés dans chaque demi-équation.
\item \textbf{Équilibrer l'oxygène avec \ce{H2O}} : 4 O dans \ce{MnO4-} $\Rightarrow$ on ajoute \ce{4H2O} à droite.
\item \textbf{Équilibrer l'hydrogène avec \ce{H+}} : 8 H à droite $\Rightarrow$ on ajoute \ce{8H+} à gauche.
\item \textbf{Équilibrer les charges avec les électrons} : à gauche $-1 + 8 = +7$, à droite $+2$ ; il faut $5\,e^-$ à gauche : \ce{MnO4- + 8H+ + 5e- -> Mn^2+ + 4H2O}.
\item \textbf{Écrire la demi-équation du fer} : \ce{Fe^2+ -> Fe^3+ + e-}.
\item \textbf{Égaliser les électrons échangés} : multiplier la demi-équation du fer par 5.
\item \textbf{Additionner et simplifier} : les $5\,e^-$ disparaissent des deux côtés.
\item \textbf{Vérifier} : conservation des atomes (1 Mn, 5 Fe, 4 O, 8 H) et des charges ($-1+10+8 = +17$ à gauche ; $+2+15 = +17$ à droite).
\end{enumerate}
\end{methode}

On obtient l'équation bilan du dosage :

\[ \ce{MnO4- + 5Fe^2+ + 8H+ -> Mn^2+ + 5Fe^3+ + 4H2O} \]

\begin{aretenir}[L'équation du dosage à connaître par cœur]
\[ \ce{MnO4- + 5Fe^2+ + 8H+ -> Mn^2+ + 5Fe^3+ + 4H2O} \]
Caractéristiques : réaction \textbf{rapide}, \textbf{totale} ($\Delta E^{\circ} = \SI{0,74}{V}$), \textbf{unique}, en \textbf{milieu acide}. Le coefficient clé est le \cle{5} devant \ce{Fe^2+} : une mole de permanganate oxyde cinq moles d'ions fer(II).
\end{aretenir}

\courssub{Pourquoi un milieu acide obligatoire — et quel acide choisir ?}

L'équation bilan consomme des ions \ce{H+} : le dosage doit donc être réalisé en \textbf{milieu acide}. Sans acidification, l'ion permanganate n'est pas réduit jusqu'à \ce{Mn^2+} (incolore) mais seulement jusqu'au dioxyde de manganèse \ce{MnO2}, un \emph{précipité brun} qui trouble la solution et fausse complètement le repérage de l'équivalence. On ajoute donc systématiquement quelques millilitres d'\textbf{acide sulfurique dilué} \ce{H2SO4} dans l'erlenmeyer avant de commencer le dosage.

\begin{attention}[Jamais d'acide chlorhydrique !]
Pour acidifier, on utilise l'acide sulfurique dilué, \textbf{jamais} l'acide chlorhydrique \ce{HCl}. En effet, les ions chlorure \ce{Cl-} sont des réducteurs : l'oxydant puissant \ce{MnO4-} les oxyderait en dichlore \ce{Cl2} selon une réaction parasite :
\[ \ce{2MnO4- + 10Cl- + 16H+ -> 2Mn^2+ + 5Cl2 + 8H2O} \]
Une partie du permanganate versé serait consommée par les ions \ce{Cl-} au lieu de réagir avec \ce{Fe^2+} : le volume équivalent mesuré serait trop grand, la concentration calculée en fer serait \textbf{fausse par excès}, et il se dégagerait de plus un gaz toxique. L'acide sulfurique, lui, fournit des ions \ce{SO4^2-} inertes vis-à-vis de \ce{MnO4-}.
\end{attention}

\courssub{Repérage de l'équivalence : l'auto-indication}

Voici le point le plus élégant de ce dosage : il ne nécessite \textbf{aucun indicateur coloré}. Observons les couleurs des espèces en présence :

\begin{center}
\begin{tabularx}{\linewidth}{|l|X|}
\hline
\rowcolor{popblueL} \textbf{Espèce} & \textbf{Couleur en solution} \\
\hline
\ce{MnO4-} (permanganate) & violette intense (presque noire en solution concentrée) \\
\hline
\ce{Mn^2+} (ion manganèse) & quasiment incolore (rose très pâle, imperceptible en solution diluée) \\
\hline
\ce{Fe^2+} (ion fer II) & verdâtre pâle, presque imperceptible en solution diluée \\
\hline
\ce{Fe^3+} (ion fer III) & jaune pâle à rouille selon la concentration \\
\hline
\end{tabularx}
\end{center}

Raisonnons goutte à goutte :

\begin{itemize}
\item \textbf{Avant l'équivalence} : chaque goutte de permanganate violet qui tombe dans l'erlenmeyer rencontre encore des ions \ce{Fe^2+}. Elle est immédiatement \textbf{décolorée} : \ce{MnO4-} est réduit en \ce{Mn^2+} incolore. La solution reste donc incolore à verdâtre (ou légèrement jaunâtre à cause des \ce{Fe^3+} formés).
\item \textbf{À l'équivalence} : il ne reste plus aucun ion \ce{Fe^2+}. La première goutte de \ce{MnO4-} versée en excès ne trouve plus rien à oxyder : elle reste intacte et colore la solution en \textbf{rose persistant}.
\end{itemize}

\begin{definition}[Auto-indication de l'équivalence]
On dit que le dosage est \cle{auto-indiqué} (ou que l'équivalence est repérée par \cle{auto-indication}) lorsque l'un des réactifs joue lui-même le rôle d'indicateur de fin de réaction. Ici, l'ion permanganate \ce{MnO4-} est à la fois le \textbf{réactif titrant} et l'\textbf{indicateur} : sa couleur violette disparaît tant qu'il reste du réducteur \ce{Fe^2+}, et la première goutte en excès donne une coloration \textbf{rose persistante} qui signale l'équivalence. Le volume alors lu à la burette est le volume équivalent $V_E$.
\end{definition}

\begin{experience}[Protocole du dosage manganimétrique des ions Fe2+]
\textbf{Dispositif} : burette graduée contenant la solution de \ce{KMnO4} de concentration $C(\ce{MnO4-})$ connue, placée au-dessus d'un erlenmeyer contenant un volume $V_P$ de solution de \ce{Fe^2+} prélevé à la pipette jaugée, additionné d'environ \SI{5}{mL} d'acide sulfurique dilué, le tout posé sur un agitateur magnétique au-dessus d'un fond blanc.
\begin{enumerate}
\item Rincer la burette avec la solution de \ce{KMnO4}, la remplir et ajuster le zéro.
\item Prélever $V_P = \SI{20,0}{mL}$ de la solution de \ce{Fe^2+} à la pipette jaugée ; verser dans l'erlenmeyer ; ajouter l'acide sulfurique dilué.
\item Verser le permanganate par fractions de plus en plus petites en agitant ; la coloration violette disparaît d'abord instantanément, puis de plus en plus lentement.
\item Dès que la teinte rose \textbf{persiste plus de 30 secondes}, fermer le robinet et lire $V_E$.
\item Refaire un deuxième dosage rapide puis goutte à goutte pour affiner la valeur de $V_E$.
\end{enumerate}
\textbf{Observation} : solution initialement verdâtre/incolore ; décoloration de chaque goutte avant $V_E$ ; virage au rose persistant à $V_E$.
\end{experience}

\begin{popfigure}
\begin{center}
\begin{tikzpicture}[scale=0.95]
% ===== AVANT EQUIVALENCE =====
\begin{scope}
% Burette
\draw[thick, popdark] (-0.25,3.2) -- (-0.25,5.6) (0.25,3.2) -- (0.25,5.6);
\draw[thick, popdark] (-0.25,3.2) -- (-0.06,2.9) (0.25,3.2) -- (0.06,2.9);
\draw[thick, popdark] (-0.06,2.9) -- (-0.06,2.55) (0.06,2.9) -- (0.06,2.55);
\draw[thick, popdark] (-0.45,2.75) -- (0.45,2.75);
\fill[poppurple!60] (-0.22,3.35) rectangle (0.22,5.5);
\node[font=\scriptsize, poppurple, right] at (0.35,4.6) {\ce{KMnO4} violet};
% Goutte
\fill[poppurple] (0,2.1) circle (2.5pt);
\draw[-{Stealth}, poppurple, thick] (0,2.45) -- (0,2.2);
% Erlenmeyer
\draw[thick, popdark] (-0.55,1.5) -- (-1.15,0) -- (1.15,0) -- (0.55,1.5);
\draw[thick, popdark] (-0.55,1.5) -- (-0.55,1.9) (0.55,1.5) -- (0.55,1.9);
\fill[popgreen!25] (-1.02,0.06) -- (1.02,0.06) -- (0.72,0.95) -- (-0.72,0.95) -- cycle;
\node[font=\scriptsize, align=center] at (0,0.5) {\ce{Fe^2+} + \ce{H+}\\incolore/verdâtre};
\node[font=\small\bfseries, popdark, align=center] at (0,-0.7) {Avant l'équivalence\\la goutte violette se décolore};
\end{scope}
% Flèche de transition
\draw[-{Stealth}, very thick, poporange] (2.0,1.2) -- (3.4,1.2) node[midway, above, font=\scriptsize\bfseries] {$V = V_E$};
% ===== APRES EQUIVALENCE =====
\begin{scope}[xshift=6.4cm]
\draw[thick, popdark] (-0.25,3.2) -- (-0.25,5.6) (0.25,3.2) -- (0.25,5.6);
\draw[thick, popdark] (-0.25,3.2) -- (-0.06,2.9) (0.25,3.2) -- (0.06,2.9);
\draw[thick, popdark] (-0.06,2.9) -- (-0.06,2.55) (0.06,2.9) -- (0.06,2.55);
\draw[thick, popdark] (-0.45,2.75) -- (0.45,2.75);
\fill[poppurple!60] (-0.22,3.35) rectangle (0.22,5.5);
\fill[poppurple] (0,2.1) circle (2.5pt);
\draw[thick, popdark] (-0.55,1.5) -- (-1.15,0) -- (1.15,0) -- (0.55,1.5);
\draw[thick, popdark] (-0.55,1.5) -- (-0.55,1.9) (0.55,1.5) -- (0.55,1.9);
\fill[poppink!45] (-1.02,0.06) -- (1.02,0.06) -- (0.72,0.95) -- (-0.72,0.95) -- cycle;
\node[font=\scriptsize, align=center] at (0,0.5) {\ce{Fe^3+} + \ce{Mn^2+}\\+ excès de \ce{MnO4-}};
\node[font=\small\bfseries, poppink!80!black, align=center] at (0,-0.7) {À l'équivalence\\rose persistant : stop !};
\end{scope}
\end{tikzpicture}
\end{center}
\legende{Repérage de l'équivalence par auto-indication : avant $V_E$, chaque goutte de \ce{MnO4-} est décolorée ; à $V_E$, la première goutte en excès colore la solution en rose persistant.}
\end{popfigure}

\courssub{Relation à l'équivalence}

À l'équivalence, les réactifs ont été mélangés dans les proportions stœchiométriques de l'équation bilan. Dressons le tableau d'avancement à cet instant précis :

\begin{center}
\begin{tabularx}{\linewidth}{|l|X|X|X|}
\hline
\rowcolor{popblueL} \textbf{Équation} & \ce{MnO4-} & $+$ \quad \ce{5Fe^2+} & $\rightarrow$ produits \\
\hline
Quantités à l'équivalence & $n_E(\ce{MnO4-})$ & $n_P(\ce{Fe^2+})$ & -- \\
\hline
Avancement final & $n_E(\ce{MnO4-}) - x_E = 0$ & $n_P(\ce{Fe^2+}) - 5x_E = 0$ & -- \\
\hline
\end{tabularx}
\end{center}

De $x_E = n_E(\ce{MnO4-})$ et $x_E = \dfrac{n_P(\ce{Fe^2+})}{5}$, on tire la relation fondamentale :

\begin{aretenir}[Relation à l'équivalence]
\[ n(\ce{Fe^2+}) = 5 \times n_E(\ce{MnO4-}) \qquad \text{soit} \qquad \boxed{\,C(\ce{Fe^2+}) = \dfrac{5 \, C(\ce{MnO4-}) \times V_E}{V_P}\,} \]
avec $C(\ce{MnO4-})$ en \si{mol.L^{-1}}, $V_E$ et $V_P$ dans la \textbf{même unité} (mL ou L), et $C(\ce{Fe^2+})$ en \si{mol.L^{-1}}.
\end{aretenir}

\begin{attention}[Erreur fréquente : oublier le facteur 5]
C'est \textbf{le} piège classique du Bac. Beaucoup d'élèves écrivent machinalement $C_1V_1 = C_2V_2$, comme en dosage acide-base, et obtiennent une concentration cinq fois trop faible. Souviens-toi : l'équivalence signifie des quantités \emph{proportionnelles aux coefficients stœchiométriques}, pas des quantités égales. Ici, 1 mole de \ce{MnO4-} consomme 5 moles de \ce{Fe^2+} : il faut donc \textbf{multiplier par 5} la quantité de permanganate versée. Vérifie toujours la cohérence : le réducteur dosé est 5 fois plus concentré que l'oxydant versé (à volumes comparables).
\end{attention}

\courssub{Application numérique type Bac}

\begin{exoresolu}[Dosage des ions Fe2+ d'une solution de sulfate de fer(II)]
\textbf{Énoncé.} Pour contrôler une solution de sulfate de fer(II), on prélève $V_P = \SI{20,0}{mL}$ de solution que l'on place dans un erlenmeyer avec de l'acide sulfurique dilué. On dose par une solution de permanganate de potassium de concentration $C(\ce{MnO4-}) = \SI{0,020}{mol.L^{-1}}$. Le virage au rose persistant est obtenu pour $V_E = \SI{12,4}{mL}$.
\begin{enumerate}
\item Écrire l'équation du dosage.
\item Calculer la concentration $C(\ce{Fe^2+})$ de la solution.
\item Calculer la concentration massique en fer de la solution. On donne $M(\text{Fe}) = \SI{55,8}{g.mol^{-1}}$.
\end{enumerate}
\tcblower
\textbf{Solution détaillée.}
\begin{enumerate}
\item Couples : \ce{MnO4-}/\ce{Mn^2+} et \ce{Fe^3+}/\ce{Fe^2+}. En égalisant les électrons échangés (5 et 1) :
\[ \ce{MnO4- + 5Fe^2+ + 8H+ -> Mn^2+ + 5Fe^3+ + 4H2O} \]
\item À l'équivalence : $n(\ce{Fe^2+}) = 5 \times n_E(\ce{MnO4-})$, donc :
\begin{align*}
C(\ce{Fe^2+}) &= \dfrac{5 \times C(\ce{MnO4-}) \times V_E}{V_P} \\
&= \dfrac{5 \times 0{,}020 \times 12{,}4}{20{,}0} = \SI{0,062}{mol.L^{-1}}
\end{align*}
\item La concentration massique en fer s'obtient par $C_m = C \times M$ :
\[ C_m(\text{Fe}) = 0{,}062 \times 55{,}8 \approx \SI{3,46}{g.L^{-1}} \]
La solution contient donc environ \SI{3,5}{g} de fer par litre.
\end{enumerate}
\textbf{Vérification du piège} : sans le facteur 5, on trouverait \SI{0,0124}{mol.L^{-1}}, valeur cinq fois trop faible — erreur éliminatoire au Bac.
\end{exoresolu}

\begin{exemplebox}[Du comprimé à la masse de fer : démarche complète]
Un comprimé de sulfate de fer(II) est dissous dans une fiole jaugée de \SI{100,0}{mL} (solution S). On dose $V_P = \SI{10,0}{mL}$ de S par du permanganate à \SI{0,010}{mol.L^{-1}} ; l'équivalence est atteinte à $V_E = \SI{10,8}{mL}$.
\begin{itemize}
\item Concentration : $C(\ce{Fe^2+}) = \dfrac{5 \times 0{,}010 \times 10{,}8}{10{,}0} = \SI{0,054}{mol.L^{-1}}$.
\item Quantité de fer dans la fiole : $n = C \times V = 0{,}054 \times 0{,}1000 = \SI{5,4e-3}{mol}$.
\item Masse de fer par comprimé : $m = n \times M = 5{,}4 \times 10^{-3} \times 55{,}8 \approx \SI{0,301}{g}$, soit environ \SI{301}{mg}.
\end{itemize}
Si l'étiquette annonce \SI{300}{mg} de fer par comprimé, le contrôle qualité est concluant.
\end{exemplebox}

\courssub{Complément : allure de la courbe de suivi potentiométrique}

\begin{savaistu}[Pour aller plus loin : le suivi potentiométrique (complément Bac)]
Au lieu de repérer l'équivalence à l'œil, on peut suivre le \textbf{potentiel $E$ de la solution} (mesuré avec une électrode de platine et une électrode de référence) en fonction du volume $V$ de permanganate versé. Avant $V_E$, le potentiel est fixé par le couple \ce{Fe^3+}/\ce{Fe^2+} (voisinage de \SI{0,77}{V}) ; après $V_E$, il est fixé par le couple \ce{MnO4-}/\ce{Mn^2+} (voisinage de \SI{1,51}{V}). Entre les deux : un \textbf{saut de potentiel} brutal dont le milieu correspond exactement à $V_E$. Cette méthode, plus précise que l'œil, est utilisée dans les laboratoires de contrôle industriels.
\end{savaistu}

\begin{popfigure}
\begin{center}
\begin{tikzpicture}
\begin{axis}[
  width=0.85\linewidth, height=6.2cm,
  xlabel={$V(\ce{MnO4-})$ versé (mL)}, ylabel={$E$ (V)},
  xmin=0, xmax=25, ymin=0.6, ymax=1.7,
  xtick={0,5,10,12.4,15,20,25}, ytick={0.77,1.51},
  yticklabels={0{,}77,1{,}51},
  grid=both, grid style={popline!40},
  label style={font=\small}, ticklabel style={font=\scriptsize},
  clip=true
]
\addplot[popcurve, domain=0:12.39, samples=200] {0.77 + 0.06*ln(12.4/(12.4-x))/2.3};
\addplot[popcurve, domain=12.41:25, samples=200] {1.51 - 0.06*ln(12.4/(x-12.4))/2.3};
\draw[dashed, thick, poporange] (axis cs:12.4,0.6) -- (axis cs:12.4,1.7);
\node[font=\small\bfseries, poporange, anchor=south west] at (axis cs:12.6,0.65) {$V_E = \SI{12,4}{mL}$};
\node[font=\scriptsize, popblue] at (axis cs:5,0.83) {couple \ce{Fe^3+}/\ce{Fe^2+}};
\node[font=\scriptsize, poppurple] at (axis cs:19,1.44) {couple \ce{MnO4-}/\ce{Mn^2+}};
\end{axis}
\end{tikzpicture}
\end{center}
\legende{Courbe qualitative $E = f(V)$ du dosage potentiométrique de \ce{Fe^2+} par \ce{MnO4-} : le saut de potentiel permet de repérer $V_E$ avec précision.}
\end{popfigure}

\begin{savaistu}[Un dosage qui sauve des vies au Cameroun]
L'anémie ferriprive est l'un des problèmes de santé publique les plus répandus au Cameroun : elle touche une proportion importante des enfants de moins de cinq ans et des femmes enceintes, souvent à cause d'une alimentation pauvre en fer assimilable et du paludisme. Les comprimés et sirops de sulfate de fer(II) distribués dans les hôpitaux et pharmacies doivent respecter strictement la teneur annoncée : un sous-dosage rend le traitement inefficace, un surdosage peut être toxique. Le dosage manganimétrique que tu viens d'étudier est précisément l'une des méthodes utilisées par les laboratoires de contrôle pharmaceutique pour vérifier, lot après lot, que chaque comprimé contient bien la masse de fer inscrite sur l'étiquette. La chimie analytique est ici une alliée directe de la médecine.
\end{savaistu}

\begin{exercice}[À toi de jouer]
Une solution de sulfate de fer(II) acidifiée est dosée par une solution de \ce{KMnO4} à \SI{0,015}{mol.L^{-1}}. Pour une prise d'essai $V_P = \SI{25,0}{mL}$, l'équivalence est repérée à $V_E = \SI{16,2}{mL}$.
\begin{enumerate}
\item Écrire les deux demi-équations puis l'équation bilan du dosage.
\item Calculer $C(\ce{Fe^2+})$.
\item Calculer la masse de fer contenue dans \SI{250}{mL} de cette solution.
\end{enumerate}
\end{exercice}

\begin{corrige}[Exercice 1]
\begin{enumerate}
\item \ce{MnO4- + 8H+ + 5e- -> Mn^2+ + 4H2O} et \ce{Fe^2+ -> Fe^3+ + e-} ($\times 5$) ; bilan : \ce{MnO4- + 5Fe^2+ + 8H+ -> Mn^2+ + 5Fe^3+ + 4H2O}.
\item $C(\ce{Fe^2+}) = \dfrac{5 \times 0{,}015 \times 16{,}2}{25{,}0} = \SI{0,0486}{mol.L^{-1}} \approx \SI{0,049}{mol.L^{-1}}$.
\item $n = 0{,}0486 \times 0{,}250 = \SI{1,215e-2}{mol}$ ; $m = n \times M = 1{,}215 \times 10^{-2} \times 55{,}8 \approx \SI{0,68}{g}$ de fer.
\end{enumerate}
\end{corrige}

Tu maîtrises désormais le dosage manganimétrique : couples, équation, milieu acide sulfurique, auto-indication au rose persistant, et la fameuse relation avec son facteur 5. Dans la section suivante, nous étudierons un deuxième dosage redox tout aussi classique — celui du diiode par les ions thiosulfate — où, cette fois, le repérage de l'équivalence nécessitera un véritable indicateur coloré : l'empois d'amidon.

\coursec{Dosage du diiode par les ions thiosulfate}

Dans la section précédente, tu as découvert le dosage des ions fer(II) \ce{Fe^2+} par les ions permanganate \ce{MnO4-} : un dosage direct, où le titrant est lui-même coloré et joue le rôle d'indicateur. Toutes les réactions d'oxydoréduction n'offrent pas ce confort. Comment doser une espèce dont la couleur disparaît progressivement, sans virage franc ? C'est exactement le problème que résout le dosage du diiode par les ions thiosulfate, probablement le dosage redox le plus utilisé au laboratoire — et le plus riche en applications concrètes, de la pharmacie à l'eau de Javel.

\courssub{Pourquoi doser le diiode ? Contexte et enjeux}

Le diiode \ce{I2} est un solide gris-noir qui, dissous dans une solution d'iodure de potassium, donne une solution \cle{brune} (parfois dite « jaune-orangé » quand elle est diluée). On le rencontre dans de nombreuses situations du quotidien et de l'industrie :

\begin{itemize}
\item la \cle{teinture d'iode}, antiseptique vendu en pharmacie, dont la teneur en \ce{I2} doit respecter une norme précise (trop concentrée, elle brûle la peau ; trop diluée, elle est inefficace) ;
\item l'\cle{eau de Javel}, dont l'ingrédient actif (l'ion hypochlorite \ce{ClO-}) ne se dose pas directement, mais \emph{indirectement} en le faisant réagir avec des ions iodure : le diiode formé est ensuite dosé ;
\item de nombreux dosages dits \emph{iodométriques} en agroalimentaire (vitamine C, sulfites dans le vin, etc.).
\end{itemize}

La vérification des indications portées sur une étiquette — par exemple « teinture d'iode à 5 \% » — repose donc souvent sur ce dosage. C'est une situation typique de \textbf{contrôle de qualité}.

\courssub{Les couples mis en jeu et la réaction de dosage}

\textbf{Intuition d'abord.} Le diiode est un oxydant modéré. Il lui faut un réducteur qui réagisse avec lui de façon \emph{totale}, \emph{rapide} et \emph{unique}. L'ion thiosulfate \ce{S2O3^2-} remplit parfaitement ce contrat : c'est un réducteur qui s'oxyde en ion tétrathionate \ce{S4O6^2-}.

Les deux couples redox mis en jeu sont :

\begin{center}
\begin{tabularx}{\linewidth}{|l|X|c|}
\hline
\rowcolor{popblueL} \textbf{Couple} & \textbf{Demi-équation électronique} & $E^\circ$ (V) \\
\hline
\ce{I2}/\ce{I-} & $\ce{I2} + 2e^- \ce{->} 2\ce{I-}$ & $0{,}62$ \\
\hline
\ce{S4O6^2-}/\ce{S2O3^2-} & $\ce{S4O6^2-} + 2e^- \ce{->} 2\ce{S2O3^2-}$ & $0{,}08$ \\
\hline
\end{tabularx}
\end{center}

Comme $E^\circ(\ce{I2}/\ce{I-}) > E^\circ(\ce{S4O6^2-}/\ce{S2O3^2-})$, la règle du gamma prévoit que le diiode (oxydant du couple de potentiel le plus élevé) oxyde les ions thiosulfate (réducteur du couple de potentiel le plus faible). La réaction est \textbf{quasi totale} ($K \gg 1$), \textbf{rapide} à température ambiante et ne nécessite \textbf{aucun catalyseur} : elle satisfait toutes les conditions d'un bon dosage.

\begin{popfigure}
\begin{center}
\begin{tikzpicture}[scale=1.0]
% Axe vertical des potentiels
\draw[->, thick, popdark] (0,0) -- (0,6.5) node[above] {$E^\circ$ (V)};
\foreach \y/\lab in {1/0, 2/0{,}2, 3/0{,}4, 4/0{,}6, 5/0{,}8} {
  \draw[popmuted] (-0.1,\y) -- (0.1,\y);
  \node[left, popmuted, font=\small] at (-0.15,\y) {\lab};
}
% Couple thiosulfate (bas)
\draw[very thick, popgreen] (0.8,0.4) -- (3.2,0.4);
\node[right, popgreen, font=\small] at (3.3,0.4) {\ce{S4O6^2-}/\ce{S2O3^2-} \; $E^\circ = 0{,}08$ V};
\node[above, popgreen, font=\footnotesize] at (2.0,0.45) {Ox : \ce{S4O6^2-}};
\node[below, popgreen, font=\footnotesize] at (2.0,0.35) {R\'ed : \ce{S2O3^2-}};
% Couple diiode (haut)
\draw[very thick, poporange] (0.8,3.1) -- (3.2,3.1);
\node[right, poporange, font=\small] at (3.3,3.1) {\ce{I2}/\ce{I-} \; $E^\circ = 0{,}62$ V};
\node[above, poporange, font=\footnotesize] at (2.0,3.15) {Ox : \ce{I2}};
\node[below, poporange, font=\footnotesize] at (2.0,3.05) {R\'ed : \ce{I-}};
% Règle du gamma
\draw[->, very thick, poppink] (2.6,3.1) .. controls (4.6,3.1) and (4.6,0.4) .. (2.6,0.4);
\node[poppink, font=\small\bfseries] at (5.6,1.75) {règle du $\gamma$};
\node[poppink, font=\footnotesize, align=center] at (5.6,1.2) {\ce{I2} oxyde\\ \ce{S2O3^2-}};
\end{tikzpicture}
\end{center}
\legende{\'Echelle des potentiels standard et règle du gamma : l'oxydant le plus fort (\ce{I2}) réagit avec le réducteur le plus fort (\ce{S2O3^2-}).}
\end{popfigure}

\textbf{Établissement de l'équation bilan.} On écrit la demi-équation de réduction du diiode et celle d'oxydation du thiosulfate (écrite dans le sens de l'oxydation), puis on additionne en éliminant les électrons :

\begin{align*}
\text{Réduction : } & \ce{I2} + 2e^- \ce{->} 2\ce{I-} \\
\text{Oxydation : } & 2\ce{S2O3^2-} \ce{->} \ce{S4O6^2-} + 2e^-
\end{align*}

Les électrons s'éliminent directement (2 échangés de part et d'autre) :

\[ \ce{I2 + 2S2O3^2- -> 2I- + S4O6^2-} \]

\begin{aretenir}[Réaction de dosage]
Le diiode est réduit en ions iodure pendant que les ions thiosulfate sont oxydés en ions tétrathionate :
\[ \ce{I2 + 2S2O3^2- -> 2I- + S4O6^2-} \]
Réaction \textbf{totale}, \textbf{rapide}, \textbf{unique}, à température ambiante, sans catalyseur. Le nombre stœchiométrique clé : \textbf{1 mol de \ce{I2} réagit avec 2 mol de \ce{S2O3^2-}}.
\end{aretenir}

\courssub{Protocole et repérage de l'équivalence}

\textbf{Dispositif.} Le montage est identique à celui du dosage fer(II)/permanganate : la solution de thiosulfate de sodium (\ce{2Na+} + \ce{S2O3^2-}), de concentration connue, est placée dans la \cle{burette} (solution \cle{titrante}) ; un volume précis $V_P$ de la solution de diiode, prélevé à la pipette jaugée, est placé dans l'erlenmeyer (solution \cle{titrée}), avec un barreau aimanté et un agitateur magnétique.

\textbf{Observation au cours du dosage.} La solution de diiode est brune. Au fur et à mesure de l'ajout de thiosulfate, le diiode est consommé : la teinte passe du brun foncé au brun pâle, puis au jaune paille très clair. L'équivalence correspond à la \textbf{décoloration complète} : la solution devient incolore (les ions \ce{I-} et \ce{S4O6^2-} sont incolores).

Le problème, c'est que l'œil perçoit mal la disparition d'une teinte jaune très pâle : on risque de dépasser l'équivalence ou de s'arrêter trop tôt. D'où l'usage d'un auxiliaire précieux : l'\cle{empois d'amidon}.

\begin{definition}[Empois d'amidon : indicateur coloré]
L'\cle{empois d'amidon} est une solution d'amidon qui forme avec le diiode un \textbf{complexe coloré bleu foncé} (presque noir), extrêmement visible même à très faible concentration en \ce{I2}. Il joue le rôle d'\textbf{indicateur coloré de fin de réaction} : tant qu'il reste du diiode, la solution est bleu foncé ; dès que tout le diiode est consommé (équivalence), la solution devient \textbf{incolore}. On l'ajoute \textbf{en fin de dosage}, lorsque la solution n'est plus que jaune pâle.
\end{definition}

\begin{attention}[Erreur fréquente : l'empois d'amidon trop tôt]
N'ajoute \textbf{jamais} l'empois d'amidon dès le début du dosage ! En présence d'une grande quantité de diiode, le complexe amidon--iode formé est \textbf{trop stable} : il « emprisonne » le diiode, qui ne réagit plus correctement avec le thiosulfate. Le virage devient lent et imprécis, et le volume équivalent mesuré est \textbf{faux}. Bon réflexe : on dose d'abord jusqu'au jaune pâle, \emph{puis} on ajoute quelques gouttes d'empois d'amidon, et on termine goutte à goutte jusqu'à la décoloration nette du bleu.
\end{attention}

\begin{experience}[Dosage d'une teinture d'iode]
\textbf{Dispositif} : burette contenant la solution de thiosulfate de sodium à $C = \SI{0,10}{mol.L^{-1}}$ ; erlenmeyer contenant $V_P = \SI{10,0}{mL}$ de solution de diiode diluée, barreau aimanté, agitateur.\par
\textbf{Protocole} :
\begin{enumerate}
\item Verser le thiosulfate en agitant : la teinte brune s'éclaircit progressivement.
\item Quand la solution est jaune pâle, ajouter 5 gouttes d'empois d'amidon : la solution vire au \textbf{bleu foncé}.
\item Continuer l'ajout \textbf{goutte à goutte} jusqu'à la décoloration complète : la solution devient \textbf{incolore}.
\item Noter le volume versé à l'équivalence $V_E$ et refaire un dosage concordant.
\end{enumerate}
\textbf{Interprétation} : à l'équivalence, tout le diiode a été consommé selon $\ce{I2 + 2S2O3^2- -> 2I- + S4O6^2-}$.
\end{experience}

\begin{popfigure}
\begin{center}
\begin{tikzpicture}[scale=0.95]
% Étape 1 : solution brune
\draw[thick, popdark] (-0.8,0) -- (-0.8,-1.8) .. controls (-0.8,-2.5) and (0.8,-2.5) .. (0.8,-1.8) -- (0.8,0);
\fill[poporange!70] (-0.72,-0.3) -- (-0.72,-1.75) .. controls (-0.72,-2.35) and (0.72,-2.35) .. (0.72,-1.75) -- (0.72,-0.3) -- cycle;
\node[align=center, font=\small] at (0,-3.0) {\textbf{Début} : solution\\ \textbf{brune} (\ce{I2} en excès)};
% Étape 2 : brun pâle + amidon = bleu foncé
\begin{scope}[xshift=4.2cm]
\draw[thick, popdark] (-0.8,0) -- (-0.8,-1.8) .. controls (-0.8,-2.5) and (0.8,-2.5) .. (0.8,-1.8) -- (0.8,0);
\fill[popblue!75] (-0.72,-0.3) -- (-0.72,-1.75) .. controls (-0.72,-2.35) and (0.72,-2.35) .. (0.72,-1.75) -- (0.72,-0.3) -- cycle;
\node[align=center, font=\small] at (0,-3.0) {\textbf{Fin approchée} : jaune p\^ale\\ + amidon $\Rightarrow$ \textbf{bleu foncé}};
\end{scope}
% Étape 3 : incolore
\begin{scope}[xshift=8.4cm]
\draw[thick, popdark] (-0.8,0) -- (-0.8,-1.8) .. controls (-0.8,-2.5) and (0.8,-2.5) .. (0.8,-1.8) -- (0.8,0);
\fill[popcream] (-0.72,-0.3) -- (-0.72,-1.75) .. controls (-0.72,-2.35) and (0.72,-2.35) .. (0.72,-1.75) -- (0.72,-0.3) -- cycle;
\node[align=center, font=\small] at (0,-3.0) {\textbf{\'Equivalence} :\\ solution \textbf{incolore}};
\end{scope}
% Flèches
\draw[->, very thick, popmuted] (1.2,-1.2) -- (3.0,-1.2) node[midway, above, font=\footnotesize] {+ \ce{S2O3^2-}};
\draw[->, very thick, popmuted] (5.4,-1.2) -- (7.2,-1.2) node[midway, above, font=\footnotesize] {goutte \`a goutte};
\end{tikzpicture}
\end{center}
\legende{Évolution des teintes au cours du dosage : brun $\rightarrow$ bleu foncé (après ajout d'empois d'amidon en fin de dosage) $\rightarrow$ incolore à l'équivalence.}
\end{popfigure}

\courssub{Relation à l'équivalence et exploitation}

À l'équivalence, les réactifs ont été mélangés dans les proportions stœchiométriques de l'équation $\ce{I2 + 2S2O3^2- -> 2I- + S4O6^2-}$ :

\[ \frac{n(\ce{I2})}{1} = \frac{n(\ce{S2O3^2-})}{2} \quad \Longleftrightarrow \quad n(\ce{I2}) = \frac{n(\ce{S2O3^2-})}{2} \]

En notant $C$ et $V_E$ la concentration et le volume équivalent du thiosulfate, $C_P$ et $V_P$ la concentration inconnue et le volume prélevé de diiode :

\begin{aretenir}[Relation d'équivalence]
\[ C_P \times V_P = \frac{C \times V_E}{2} \quad \Longleftrightarrow \quad \boxed{C(\ce{I2}) = \frac{C(\ce{S2O3^2-}) \times V_E}{2 \times V_P}} \]
Le facteur $\dfrac{1}{2}$ vient des coefficients stœchiométriques : 1 mol de \ce{I2} pour 2 mol de \ce{S2O3^2-}. C'est \textbf{l'erreur de stœchiométrie la plus fréquente} au Bac : vérifie toujours les coefficients avant d'écrire la relation !
\end{aretenir}

\begin{exemplebox}[Calcul de la concentration d'une solution de diiode]
On dose $V_P = \SI{10,0}{mL}$ d'une solution de diiode par une solution de thiosulfate de sodium de concentration $C = \SI{0,10}{mol.L^{-1}}$. L'équivalence (décoloration du bleu d'amidon) est obtenue pour $V_E = \SI{8,0}{mL}$.\par
\textbf{Calcul} :
\[ C(\ce{I2}) = \frac{C \times V_E}{2 \times V_P} = \frac{0{,}10 \times 8{,}0}{2 \times 10{,}0} = \SI{0,040}{mol.L^{-1}} \]
La solution de diiode a une concentration de $\SI{4,0e-2}{mol.L^{-1}}$, soit une concentration massique $t = C \times M(\ce{I2}) = 0{,}040 \times 253{,}8 \approx \SI{10,2}{g.L^{-1}}$.
\end{exemplebox}

\courssub{Complément Bac : dosage en retour de l'eau de Javel}

L'eau de Javel doit son pouvoir désinfectant à l'ion hypochlorite \ce{ClO-}, oxydant du couple \ce{ClO-}/\ce{Cl2}. Problème : on ne peut pas doser \ce{ClO-} directement avec le thiosulfate de façon nette. On utilise alors une stratégie en deux temps, dite \textbf{dosage indirect} (ou \textbf{en retour}) :

\begin{enumerate}
\item \textbf{Étape 1 — formation de diiode.} On ajoute à un volume connu d'eau de Javel un \textbf{excès} d'ions iodure \ce{I-} en milieu acide. Les ions hypochlorite oxydent les iodures :
\[ \ce{ClO- + 2I- + 2H+ -> Cl- + I2 + H2O} \]
Chaque mole de \ce{ClO-} produit exactement \textbf{1 mole de \ce{I2}}.
\item \textbf{Étape 2 — dosage du diiode formé.} On dose le diiode ainsi libéré par la solution de thiosulfate, comme vu plus haut : $\ce{I2 + 2S2O3^2- -> 2I- + S4O6^2-}$.
\end{enumerate}

En remontant la chaîne : $n(\ce{ClO-}) = n(\ce{I2})_{\text{formé}} = \dfrac{n(\ce{S2O3^2-})}{2}$.

\begin{methode}[Mener un dosage en retour (ou indirect)]
\begin{enumerate}
\item \textbf{Identifier} pourquoi le dosage direct est impossible (réaction lente, repérage impossible, espèce instable).
\item \textbf{Faire réagir} l'espèce à doser avec un réactif en \textbf{excès connu} (ou qui la transforme quantitativement en un produit dosable, ici \ce{I2}).
\item \textbf{Doser} le produit formé (ou l'excès restant) par une réaction de dosage classique.
\item \textbf{Remonter} les relations stœchiométriques étape par étape pour trouver la quantité initiale. Écris chaque équation et chaque relation d'équivalence séparément : c'est la clé pour ne pas te tromper.
\end{enumerate}
\end{methode}

\begin{exoresolu}[Degré chlorométrique d'une eau de Javel]
\textbf{Énoncé.} On prélève $V_0 = \SI{10,0}{mL}$ d'une eau de Javel diluée 10 fois. On ajoute un excès de KI et d'acide. Le diiode formé est dosé par une solution de thiosulfate à $C = \SI{0,10}{mol.L^{-1}}$ ; l'équivalence est atteinte pour $V_E = \SI{12,0}{mL}$. Déterminer la concentration en \ce{ClO-} de l'eau de Javel commerciale.
\tcblower
\textbf{Solution.}\par
\textbf{Étape 1} : quantité de thiosulfate versée :
\[ n(\ce{S2O3^2-}) = C \times V_E = 0{,}10 \times 12{,}0 \times 10^{-3} = \SI{1,2e-3}{mol} \]
\textbf{Étape 2} : d'après $\ce{I2 + 2S2O3^2- -> 2I- + S4O6^2-}$ :
\[ n(\ce{I2}) = \frac{n(\ce{S2O3^2-})}{2} = \SI{6,0e-4}{mol} \]
\textbf{Étape 3} : d'après $\ce{ClO- + 2I- + 2H+ -> Cl- + I2 + H2O}$, $n(\ce{ClO-}) = n(\ce{I2}) = \SI{6,0e-4}{mol}$ dans \SI{10,0}{mL} de solution diluée :
\[ C_{\text{diluée}} = \frac{6{,}0 \times 10^{-4}}{10{,}0 \times 10^{-3}} = \SI{0,060}{mol.L^{-1}} \]
\textbf{Étape 4} : l'eau de Javel commerciale est 10 fois plus concentrée :
\[ C(\ce{ClO-}) = 10 \times 0{,}060 = \SI{0,60}{mol.L^{-1}} \]
\end{exoresolu}

\begin{attention}[Pièges classiques de ce dosage]
\begin{itemize}
\item Oublier le facteur $1/2$ dans la relation d'équivalence (1 \ce{I2} pour 2 \ce{S2O3^2-}).
\item Ajouter l'empois d'amidon au début du dosage (complexe trop stable).
\item Dans un dosage en retour, oublier de tenir compte de la \textbf{dilution} de l'échantillon.
\item Confondre le volume prélevé $V_P$ (pipette) et le volume équivalent $V_E$ (burette) dans la formule.
\end{itemize}
\end{attention}

\begin{savaistu}[La teinture d'iode contrôlée à Douala]
La teinture d'iode vendue dans les pharmacies de Douala, Yaoundé ou Bafoussam est un antiseptique dont la teneur en diiode est strictement réglementée. Les laboratoires de contrôle de qualité vérifient cette teneur précisément par le dosage iodométrique que tu viens d'étudier : un prélèvement de teinture, une solution étalon de thiosulfate de sodium, quelques gouttes d'empois d'amidon — et la décoloration du bleu donne la réponse. Le même principe sert à contrôler l'eau de Javel produite localement et même la teneur en vitamine C de certains jus de fruits !
\end{savaistu}

\begin{exercice}[À toi de jouer]
On dose $V_P = \SI{20,0}{mL}$ d'une solution de diiode par du thiosulfate de sodium à $C = \SI{0,050}{mol.L^{-1}}$. L'équivalence est repérée à $V_E = \SI{14,4}{mL}$. Calcule la concentration molaire puis la concentration massique de la solution de diiode ($M(\ce{I2}) = \SI{253,8}{g.mol^{-1}}$).
\end{exercice}

\begin{corrige}[Exercice]
$n(\ce{S2O3^2-}) = 0{,}050 \times 14{,}4 \times 10^{-3} = \SI{7,2e-4}{mol}$. D'où $n(\ce{I2}) = \dfrac{7{,}2 \times 10^{-4}}{2} = \SI{3,6e-4}{mol}$ dans \SI{20,0}{mL} :
\[ C(\ce{I2}) = \frac{3{,}6 \times 10^{-4}}{20{,}0 \times 10^{-3}} = \SI{1,8e-2}{mol.L^{-1}} \]
Concentration massique : $t = 0{,}018 \times 253{,}8 \approx \SI{4,6}{g.L^{-1}}$.
\end{corrige}

Tu maîtrises désormais les deux grands dosages redox du programme. Il reste à synthétiser tout cela dans une méthode générale d'exploitation : c'est l'objet de la dernière section.

\coursec{Exploitation d'un dosage : méthode complète}

Dans les sections précédentes, tu as appris à réaliser concrètement deux dosages d'oxydoréduction majeurs : celui des ions fer(II) \ce{Fe^2+} par les ions permanganate \ce{MnO4-} en milieu acide, puis celui du diiode \ce{I2} par les ions thiosulfate \ce{S2O3^2-}. Tu sais donc maintenant manipuler la burette, repérer l'équivalence grâce à un changement de couleur caractéristique et noter le volume versé $V_E$. Mais au Baccalauréat TI, l'examinateur ne te demande pas seulement de manipuler : il te demande surtout d'\emph{exploiter} les résultats, c'est-à-dire de transformer un volume lu sur une burette en une concentration, puis en une conclusion concrète (« l'étiquette est-elle conforme ? », « le produit est-il titré correctement ? »). C'est toute l'affaire de cette dernière section : te donner une méthode en six étapes, rigoureuse et universelle, qui fonctionne pour \emph{tous} les exercices de dosage d'oxydoréduction, quel que soit le couple redox mis en jeu.

\courssub{Vue d'ensemble : la démarche en six étapes}

Tout exercice de dosage se résout selon le même schéma logique. Retiens-le comme une recette de cuisine : si tu suis les étapes dans l'ordre, sans en sauter, tu ne peux pas te tromper.

\begin{methode}[Fiche récapitulative : les six étapes de tout exercice de dosage au Bac]
\begin{enumerate}
\item \textbf{Identifier} l'espèce titrée (dans le bécher/erlenmeyer, concentration inconnue), l'espèce titrante (dans la burette, concentration connue) et les deux couples redox mis en jeu.
\item \textbf{Écrire} les deux demi-équations électroniques, puis l'équation bilan de la réaction de dosage, équilibrée (atomes et charges).
\item \textbf{Décrire} le repérage de l'équivalence : quel changement de couleur observe-t-on, et quelle en est l'origine chimique (réactif auto-indicateur ou indicateur ajouté) ?
\item \textbf{Écrire} la relation stœchiométrique à l'équivalence, à partir des coefficients de l'équation bilan.
\item \textbf{Calculer} la concentration inconnue, avec les unités, en tenant compte d'une éventuelle dilution préalable (facteur de dilution $F$), et en respectant les chiffres significatifs.
\item \textbf{Conclure} en comparant le résultat à la valeur attendue (étiquette, norme, certificat d'analyse) grâce à l'écart relatif en pourcentage.
\end{enumerate}
\end{methode}

\begin{popfigure}
\begin{center}
\begin{tikzpicture}[
node distance=0.5cm,
etape/.style={rectangle, rounded corners=3pt, draw=popblue, fill=popblueL, text width=11cm, align=left, inner sep=5pt, font=\small},
etapeF/.style={rectangle, rounded corners=3pt, draw=popgreen, fill=popgreenL, text width=11cm, align=left, inner sep=5pt, font=\small},
fleche/.style={-{Stealth[length=2.5mm]}, thick, poporange}
]
\node[etape] (e1) {\textbf{Étape 1 — Identifier} : espèce titrée (bécher), espèce titrante (burette), couples redox (\emph{ex.} \ce{MnO4-}/\ce{Mn^2+} et \ce{Fe^3+}/\ce{Fe^2+}).};
\node[etape, below=of e1] (e2) {\textbf{Étape 2 — Équilibrer} : demi-équations $\rightarrow$ équation bilan de la réaction de dosage (totale, rapide, unique).};
\node[etape, below=of e2] (e3) {\textbf{Étape 3 — Repérer l'équivalence} : changement de couleur (réactif coloré \ce{MnO4-} ou indicateur amidon/\ce{I2}).};
\node[etape, below=of e3] (e4) {\textbf{Étape 4 — Relation à l'équivalence} : $\dfrac{n_{\text{titré}}}{a} = \dfrac{n_{\text{titrant}}}{b}$ ($a$, $b$ : coefficients stœchiométriques).};
\node[etape, below=of e4] (e5) {\textbf{Étape 5 — Calculer $C_{\text{inconnue}}$} : conversion unités, facteur de dilution $F = V_f/V_0$ si dilution, chiffres significatifs.};
\node[etapeF, below=of e5] (e6) {\textbf{Étape 6 — Conclure} : écart relatif $\varepsilon = \dfrac{|C_{\text{mesurée}} - C_{\text{étiquette}}|}{C_{\text{étiquette}}} \times 100$ (\%). $\varepsilon < 5\,\% \Rightarrow$ conforme.};
\draw[fleche] (e1) -- (e2);
\draw[fleche] (e2) -- (e3);
\draw[fleche] (e3) -- (e4);
\draw[fleche] (e4) -- (e5);
\draw[fleche] (e5) -- (e6);
\end{tikzpicture}
\end{center}
\legende{Organigramme des six étapes de résolution d'un exercice de dosage d'oxydoréduction.}
\end{popfigure}

\courssub{Étape 1 : identifier les espèces et les couples redox}

Avant tout calcul, lis l'énoncé en te posant trois questions fondamentales : \emph{quelle est la solution dans le bécher (ou l'erlenmeyer) ?} (c'est la solution \cle{titrée}, dont on cherche la concentration), \emph{quelle est la solution dans la burette ?} (c'est la solution \cle{titrante}, de concentration connue, l'« outil de mesure »), et \emph{quels sont les couples redox en présence ?}

Pour trouver les couples, repère les deux espèces qui réagissent et cherche pour chacune sa forme conjuguée (forme oxydée / forme réduite).

\begin{exemplebox}[Identification pour les deux dosages types du programme]
\begin{itemize}
\item \textbf{Dosage \ce{Fe^2+} par \ce{MnO4-} (milieu acide)} :
      \begin{itemize}
      \item Titré : \ce{Fe^2+} (réducteur) $\rightarrow$ couple \ce{Fe^3+}/\ce{Fe^2+}.
      \item Titrant : \ce{MnO4-} (oxydant) $\rightarrow$ couple \ce{MnO4-}/\ce{Mn^2+}.
      \end{itemize}
\item \textbf{Dosage \ce{I2} par \ce{S2O3^2-}} :
      \begin{itemize}
      \item Titré : \ce{I2} (oxydant) $\rightarrow$ couple \ce{I2}/\ce{I-}.
      \item Titrant : \ce{S2O3^2-} (réducteur) $\rightarrow$ couple \ce{S4O6^2-}/\ce{S2O3^2-}.
      \end{itemize}
\end{itemize}
\end{exemplebox}

\begin{attention}[Piège classique : inverser titré et titrant]
Une erreur fréquente consiste à inverser les rôles dans la relation finale. Retiens ceci : la solution titrante est \emph{toujours} celle dont la concentration est connue (c'est l'« outil de mesure »), et c'est elle qui est versée goutte à goutte. Le volume $V_E$ lu à la burette concerne donc \emph{toujours} le titrant. Le volume de solution titrée ($V_{\text{titré}}$) est celui prélevé à la pipette jaugée.
\end{attention}

\courssub{Étape 2 : écrire les demi-équations puis l'équation bilan}

La réaction de dosage doit être \cle{totale} (avancement final maximal), \cle{rapide} (cinétique instantanée à l'échelle de la manipulation) et \cle{unique} (pas de réaction parasite) : c'est elle qui servira de base à tous les calculs. Il faut donc l'équilibrer avec le plus grand soin, en milieu acide si nécessaire (ajout de \ce{H+} et \ce{H2O}).

\paragraph{Dosage des ions \ce{Fe^2+} par les ions \ce{MnO4-} (milieu acide \ce{H2SO4})}

Demi-équation de réduction du permanganate (couple \ce{MnO4-}/\ce{Mn^2+}) :
\[ \ce{MnO4- + 8H+ + 5e- -> Mn^2+ + 4H2O} \]

Demi-équation d'oxydation du fer(II) (couple \ce{Fe^3+}/\ce{Fe^2+}) :
\[ \ce{Fe^2+ -> Fe^3+ + e-} \]

On multiplie la seconde par 5 pour que les électrons s'éliminent, puis on additionne les deux demi-équations :
\[ \ce{MnO4- + 5Fe^2+ + 8H+ -> Mn^2+ + 5Fe^3+ + 4H2O} \]
Cette équation est totale, rapide à température ambiante et ne nécessite aucun catalyseur : c'est une excellente réaction de dosage. Les coefficients stœchiométriques (1 pour \ce{MnO4-}, 5 pour \ce{Fe^2+}) sont la clé de l'étape 4.

\paragraph{Dosage du diiode \ce{I2} par les ions thiosulfate \ce{S2O3^2-}}

Demi-équation de réduction du diiode (couple \ce{I2}/\ce{I-}) :
\[ \ce{I2 + 2e- -> 2I-} \]

Demi-équation d'oxydation du thiosulfate (couple \ce{S4O6^2-}/\ce{S2O3^2-}) :
\[ \ce{2S2O3^2- -> S4O6^2- + 2e-} \]

Équation bilan (électrons simplifiés) :
\[ \ce{I2 + 2S2O3^2- -> 2I- + S4O6^2-} \]
Cette réaction est totale, rapide et se déroule en milieu neutre ou faiblement acide. Coefficients : 1 pour \ce{I2}, 2 pour \ce{S2O3^2-}.

\begin{propriete}[Équations bilan des deux dosages de référence]
\begin{center}
\begin{tabularx}{\linewidth}{|l|X|X|}\hline
\textbf{Dosage} & \textbf{Équation bilan} & \textbf{Coeff. $(a,b)$} \\\hline
\ce{Fe^2+} / \ce{MnO4-} & \ce{MnO4- + 5Fe^2+ + 8H+ -> Mn^2+ + 5Fe^3+ + 4H2O} & $a=5,\; b=1$ \\\hline
\ce{I2} / \ce{S2O3^2-} & \ce{I2 + 2S2O3^2- -> 2I- + S4O6^2-} & $a=1,\; b=2$ \\\hline
\end{tabularx}
\end{center}
\end{propriete}

\courssub{Étape 3 : décrire le repérage de l'équivalence}

Il ne suffit pas d'écrire « changement de couleur » : il faut préciser \emph{lequel} (couleur initiale $\rightarrow$ couleur finale) et \emph{pourquoi} (quelle espèce est responsable). Deux situations se présentent au Bac TI :

\begin{itemize}
\item \textbf{Réactif auto-indicateur (cas \ce{MnO4-})} : l'ion \ce{MnO4-} est intensément violet, alors que le produit \ce{Mn^2+} est incolore (très pâle). Avant l'équivalence, chaque goutte de permanganate versée est consommée instantanément et se décolore ; la solution dans l'erlenmeyer prend la couleur des ions fer (vert pâle \ce{Fe^2+} / jaune \ce{Fe^3+}). \textbf{À l'équivalence}, la première goutte de \ce{MnO4-} en excès n'est plus consommée : elle donne à la solution une teinte \textbf{rose pâle persistante} (au moins 30 secondes). On passe donc de « incolore/vert/jaune » à « rose pâle persistant ».
\item \textbf{Indicateur coloré ajouté (cas \ce{I2} / \ce{S2O3^2-})} : le diiode \ce{I2} est brun/jaune en solution aqueuse. On ajoute de l'\cle{empois d'amidon} (solution de starch) qui forme avec \ce{I2} un complexe d'inclusion \textbf{bleu foncé} (noir bleuté). Avant l'équivalence, la solution est bleue. \textbf{À l'équivalence}, tout le diiode a réagi : on observe la \textbf{décoloration brutale} de la solution (passage du bleu foncé à l'incolore). L'indicateur est ajouté \emph{en fin de dosage} (quand la couleur jaune pâlit) pour ne pas piéger l'iode dans l'amidon trop tôt.
\end{itemize}

\begin{attention}[Rédaction attendue à l'écrit]
Dans ta copie, rédige toujours une phrase complète, par exemple :
\begin{itemize}
\item Pour \ce{Fe^2+}/\ce{MnO4-} : « L'équivalence est repérée par l'apparition d'une coloration \textbf{rose pâle persistante} due à l'excès d'ions \ce{MnO4-}, signe que tous les ions \ce{Fe^2+} ont été consommés. »
\item Pour \ce{I2}/\ce{S2O3^2-} : « L'équivalence est repérée par la \textbf{décoloration de la solution} (disparition du complexe bleu \ce{I2}-amidon), signe que tout le diiode a été réduit en iodures. »
\end{itemize}
\end{attention}

\courssub{Étape 4 : écrire la relation stœchiométrique à l'équivalence}

\begin{definition}[Équivalence d'un dosage]
L'\cle{équivalence} est atteinte lorsque les réactifs ont été mélangés dans les proportions stœchiométriques exactes de l'équation bilan : le titré et le titrant sont alors tous deux entièrement consommés (quantités de matière nulles pour les deux réactifs limités).
\end{definition}

Si l'équation bilan s'écrit $a\,\text{A} + b\,\text{B} \rightarrow$ produits, où A est l'espèce titrée (volume $V_A$, concentration $C_A$) et B l'espèce titrante (volume $V_E$ à l'équivalence, concentration $C_B$), alors à l'équivalence :
\[ \frac{n_{\text{A}}}{a} = \frac{n_{\text{B}}}{b} \quad \text{soit} \quad \frac{C_{\text{A}} \times V_{\text{A}}}{a} = \frac{C_{\text{B}} \times V_{E}}{b} \]

\paragraph{Application aux deux dosages de référence}

\begin{itemize}
\item \textbf{\ce{Fe^2+} / \ce{MnO4-}} ($a=5$, $b=1$) :
      \[ \frac{C_{\ce{Fe^2+}} \times V_{\ce{Fe^2+}}}{5} = C_{\ce{MnO4-}} \times V_{E} \quad \Longrightarrow \quad C_{\ce{Fe^2+}} = \frac{5 \, C_{\ce{MnO4-}} \, V_{E}}{V_{\ce{Fe^2+}}} \]
\item \textbf{\ce{I2} / \ce{S2O3^2-}} ($a=1$, $b=2$) :
      \[ \frac{C_{\ce{I2}} \times V_{\ce{I2}}}{1} = \frac{C_{\ce{S2O3^2-}} \times V_{E}}{2} \quad \Longrightarrow \quad C_{\ce{I2}} = \frac{C_{\ce{S2O3^2-}} \, V_{E}}{2 \, V_{\ce{I2}}} \]
\end{itemize}

\begin{attention}[Piège majeur : oublier les coefficients stœchiométriques]
L'erreur la plus fréquente (et la plus sanctionnée) est d'écrire $n_{\text{titré}} = n_{\text{titrant}}$ (relation 1/1) alors que les coefficients ne sont pas égaux. Pour \ce{Fe^2+}/\ce{MnO4-}, il faut \textbf{5 moles de \ce{Fe^2+}} pour \textbf{1 mole de \ce{MnO4-}} : oublier le facteur 5 fausse le résultat d'un facteur 5. \textbf{Recopie toujours l'équation bilan juste au-dessus de ta relation d'équivalence} pour visualiser les coefficients $a$ et $b$.
\end{attention}

\courssub{Étape 5 : calculer la concentration inconnue}

Le calcul se mène avec rigueur :
\begin{enumerate}
\item Volumes : on peut garder les \si{mL} des deux côtés car ils se simplifient (ou convertir en \si{L} pour obtenir directement \si{mol.L^{-1}}).
\item Résultat en \si{mol.L^{-1}} (ou \si{g.L^{-1}} si masse molaire demandée).
\item Chiffres significatifs : on garde autant de chiffres significatifs que la donnée la moins précise (souvent 3 au Bac : ex. \SI{2,00e-2}{mol.L^{-1}}, \SI{12,4}{mL} $\rightarrow$ 3 chiffres significatifs).
\end{enumerate}

\paragraph{Cas d'une dilution préalable (très fréquent au Bac)}
Il arrive très souvent que la solution commerciale (solution mère $S_0$) soit trop concentrée pour être dosée directement (dépassement de l'échelle de la burette, ou réaction trop violente). On la \cle{dilue} d'abord selon le protocole rigoureux :
\begin{enumerate}
\item Prélèvement de $V_0$ de $S_0$ à la \textbf{pipette jaugée} (ex. \SI{10,0}{mL} ou \SI{20,0}{mL}).
\item Transfert dans une \textbf{fiole jaugée} de volume $V_f$ (ex. \SI{100,0}{mL} ou \SI{200,0}{mL}).
\item Complétion à l'eau distillée jusqu'au trait de jauge (bas du ménisque au niveau du trait). On obtient la solution fille $S_1$ (diluée).
\end{enumerate}

Le \cle{facteur de dilution} est :
\[ F = \frac{V_f}{V_0} = \frac{C_0}{C_1} \]
où $C_1$ est la concentration de la solution diluée (celle qui est réellement dosée) et $C_0$ la concentration de la solution mère (celle de l'étiquette). Une fois $C_1$ calculée par le dosage, on remonte à la solution mère :
\[ C_0 = F \times C_1 \]

\begin{attention}[Erreur fatale : oublier le facteur de dilution]
Quand l'énoncé mentionne « on dilue 10 fois », « on prélève \SI{10,0}{mL} que l'on complète à \SI{100,0}{mL} », « fiole jaugée \SI{200,0}{mL} », le dosage ne donne que la concentration de la solution \emph{diluée} ($C_1$). Si tu oublies de multiplier par $F$ à la fin, ta concentration $C_0$ sera $F$ fois trop faible et ta conclusion fausse. \textbf{Souligne dans l'énoncé toute mention de fiole jaugée ou de dilution, et prévois la multiplication par $F$ dès le début de ta rédaction.}
\end{attention}

\begin{popfigure}
\begin{center}
\begin{tikzpicture}[scale=0.9, every node/.style={font=\small}]
% Styles
\tikzset{
  flask/.style={draw=popblue, thick, fill=popblue!10},
  pipette/.style={draw=popgreen, thick, fill=popgreen!10},
  beaker/.style={draw=poppurple, thick, fill=poppurple!10},
  arrow/.style={-{Stealth[length=3mm]}, very thick, poporange},
  label/.style={popdark, align=center, font=\scriptsize}
}
% Fiole jaugée (Solution diluée S1)
\draw[flask] (0,0) .. controls (-1.2,0.2) and (-1.2,2.0) .. (-0.3,2.8);
\draw[flask] (0,0) .. controls (1.2,0.2) and (1.2,2.0) .. (0.3,2.8);
\draw[flask] (-0.3,2.8) -- (-0.3,3.8);
\draw[flask] (0.3,2.8) -- (0.3,3.8);
\draw[popblue] (-0.3,3.8) -- (-0.4,4.0);
\draw[popblue] (0.3,3.8) -- (0.4,4.0);
\draw[popblue, thick] (-0.3,3.4) -- (0.3,3.4) node[midway, above=2pt, label] {trait de jauge $V_f$};
\fill[popblue!30] (0,0.05) .. controls (-1.15,0.25) and (-1.15,1.95) .. (-0.28,2.75) -- (0.28,2.75) .. controls (1.15,1.95) and (1.15,0.25) .. (0,0.05);
\node[label] at (0,1.4) {Solution diluée $S_1$\\($C_1$, $V_f$)};
\node[label] at (0,-0.7) {Fiole jaugée $V_f$};
% Pipette (Prélèvement $V_0$ de $S_0$)
\draw[pipette] (4.5,1.0) -- (4.5,3.5);
\draw[pipette] (4.9,1.0) -- (4.9,3.5);
\draw[pipette] (4.5,1.0) -- (4.7,0.3) -- (4.9,1.0);
\fill[popgreen!30] (4.52,1.05) -- (4.52,2.8) -- (4.88,2.8) -- (4.88,1.05) -- (4.7,0.4) -- cycle;
\node[label] at (4.7,4.1) {Pipette jaugée $V_0$\\(Solution mère $S_0$, $C_0$)};
% Flèche Pipette -> Fiole
\draw[arrow] (3.8,2.5) -- (1.0,2.5) node[midway, above, label] {On verse $V_0$\\dans la fiole};
% Bécher de dosage (Prélèvement $V_p$ de $S_1$)
\draw[beaker] (7.5,0) -- (7.5,2.2);
\draw[beaker] (9.5,0) -- (9.5,2.2);
\draw[beaker] (7.5,0) -- (9.5,0);
\fill[poppurple!20] (7.55,0.05) -- (7.55,1.5) -- (9.45,1.5) -- (9.45,0.05) -- cycle;
\node[label] at (8.5,0.8) {Prélèvement $V_p$ de $S_1$\\(dosage dans l'erlenmeyer)};
\node[label] at (8.5,-0.7) {Erlenmeyer / Bécher};
% Flèche Fiole -> Bécher
\draw[arrow] (1.2,0.5) -- (6.8,0.5) node[midway, above, label] {On prélève $V_p$ à la pipette\\dans la solution diluée};
\end{tikzpicture}
\end{center}
\legende{Chaîne de dilution suivie d'un prélèvement pour dosage : la solution mère $S_0$ ($C_0$) est diluée dans une fiole jaugée ($V_f$) pour donner $S_1$ ($C_1 = C_0/F$), puis un volume $V_p$ de $S_1$ est prélevé pour être dosé.}
\end{popfigure}

\courssub{Étape 6 : conclure par l'écart relatif}

Un dosage sert souvent à \emph{vérifier une indication portée sur une étiquette} (concentration d'eau de Javel, titre acide d'un vinaigre, teneur en fer d'un complément alimentaire, pureté d'un réactif). Pour comparer la valeur mesurée $C_{\text{mes}}$ (issue de ton calcul) à la valeur annoncée $C_{\text{étq}}$, on calcule l'\cle{écart relatif} :
\[ \varepsilon = \frac{|C_{\text{mes}} - C_{\text{étq}}|}{C_{\text{étq}}} \times 100 \quad (\text{en } \%) \]

En pratique, au laboratoire comme au Bac, on considère généralement que :
\begin{itemize}
\item Si $\varepsilon \le 5\,\%$ : l'étiquette est \textbf{conforme} (les écarts s'expliquent par les incertitudes expérimentales inévitables : lecture burette, trait de jauge, température, pureté des réactifs, dégradation naturelle du produit).
\item Si $\varepsilon > 5\,\%$ : l'étiquette est \textbf{non conforme} (on suspecte une erreur de manipulation majeure, une mauvaise conservation du produit — lumière, chaleur — ou une indication mensongère).
\end{itemize}

\begin{exemplebox}[Vérification de l'étiquette d'une eau de Javel (contexte camerounais)]
Une eau de Javel commerciale achetée au marché de Douala annonce une concentration en ions hypochlorite $C_{\text{étq}} = \SI{0,50}{mol.L^{-1}}$. Après dilution au 1/10 ($F=10$) puis dosage iodométrique (les ions \ce{ClO-} oxydent l'excès d'iodure \ce{I-} en diiode \ce{I2}, que l'on dose par le thiosulfate \ce{S2O3^2-}), l'exploitation des mesures donne pour la solution diluée $C_1 = \SI{0,047}{mol.L^{-1}}$ (en \ce{ClO-} équivalent).
Concentration de la solution mère retrouvée :
\[ C_{\text{mes}} = F \times C_1 = 10 \times \SI{0,047}{mol.L^{-1}} = \SI{0,47}{mol.L^{-1}} \]
Écart relatif :
\[ \varepsilon = \frac{|0,47 - 0,50|}{0,50} \times 100 = 6\,\% \]
\textbf{Conclusion :} L'écart dépasse légèrement 5\,\%. L'eau de Javel s'est probablement partiellement décomposée sous l'effet de la chaleur et de la lumière (les ions \ce{ClO-} sont instables, d'où l'importance de les stocker au frais et dans l'obscurité). Le produit est un peu moins concentré que ne l'annonce l'étiquette : \textbf{non conforme}.
\end{exemplebox}

\courssub{Dispositif expérimental de dosage (rappel schématique)}

Le savoir-faire « Schématiser et décrire le dispositif expérimental du dosage » est exigible. Voici le montage type pour un dosage au permanganate (auto-indicateur) ou au thiosulfate (avec amidon).

\begin{popfigure}
\begin{center}
\begin{tikzpicture}[scale=0.85, every node/.style={font=\small}]
% Support
\draw[thick, popdark] (0,0) -- (8,0); % Paillasse
\draw[thick, popdark] (1.5,0) -- (1.5,5.5); % Montant
\draw[thick, popdark] (1.5,5.5) -- (3.5,5.5); % Bras
\draw[thick, popdark] (3.5,5.5) -- (3.5,5.0); % Descente pince
% Pince
\draw[thick, popdark] (3.5,5.0) -- (3.2,4.7);
\draw[thick, popdark] (3.5,5.0) -- (3.8,4.7);
% Burette
\draw[thick, popblue] (3.5,4.7) -- (3.5,1.5);
\draw[thick, popblue] (3.9,4.7) -- (3.9,1.5);
\draw[thick, popblue] (3.5,4.7) -- (3.9,4.7); % Haut
\draw[thick, popblue] (3.5,1.5) -- (3.7,1.0) -- (3.9,1.5); % Pointe
\fill[popblue!30] (3.52,1.55) -- (3.52,4.2) -- (3.88,4.2) -- (3.88,1.55) -- cycle;
\node[popdark, right=0.2cm] at (3.9,3.5) {Burette};
\node[popdark, right=0.2cm] at (3.9,3.0) {Titrant ($C$ connu)};
\node[popdark, right=0.2cm] at (3.9,2.5) {$V_E$ lu au bas du ménisque};
% Erlenmeyer
\draw[thick, poppurple] (5.5,0.5) .. controls (5.0,0.5) and (5.0,1.5) .. (5.5,1.5);
\draw[thick, poppurple] (6.5,0.5) .. controls (7.0,0.5) and (7.0,1.5) .. (6.5,1.5);
\draw[thick, poppurple] (5.5,1.5) -- (5.5,2.5);
\draw[thick, poppurple] (6.5,1.5) -- (6.5,2.5);
\draw[thick, poppurple] (5.5,2.5) -- (5.3,2.8);
\draw[thick, poppurple] (6.5,2.5) -- (6.7,2.8);
\fill[poppurple!15] (5.55,0.55) .. controls (5.1,0.55) and (5.1,1.45) .. (5.55,1.45) -- (6.45,1.45) .. controls (6.9,1.45) and (6.9,0.55) .. (6.45,0.55) -- cycle;
\node[popdark, align=center] at (6.0,1.0) {Titré ($C$ inconnu)\\+ indicateur si besoin};
\node[popdark, below=0.1cm] at (6.0,0.5) {Erlenmeyer sur carreau blanc};
% Agitateur magnétique (optionnel)
\draw[popmuted, thick] (6.0,0.3) circle (0.4);
\node[popmuted, font=\tiny] at (6.0,0.3) {M};
% Légendes
\node[label, anchor=north west] at (0.2,5.2) {Montage type : Burette + Erlenmeyer sur fond blanc (carreau de papier)};
\end{tikzpicture}
\end{center}
\legende{Dispositif expérimental de dosage : burette fixée sur support, erlenmeyer contenant la solution titrée (et indicateur le cas échéant) posé sur un fond blanc pour bien voir le changement de couleur.}
\end{popfigure}

\courssub{Exercice résolu : application complète de la méthode}

\begin{exoresolu}[Dosage d'une solution de fer(II) avec dilution — Type Bac TI]
\textbf{Énoncé.} Une solution $S_0$ de sulfate de fer(II) \ce{FeSO4} est trop concentrée pour être dosée directement. On prélève $V_0 = \SI{20,0}{mL}$ de $S_0$ à la pipette jaugée que l'on introduit dans une fiole jaugée de $V_f = \SI{200,0}{mL}$ et l'on complète avec de l'eau distillée jusqu'au trait de jauge : on obtient la solution $S_1$. On prélève ensuite $V_1 = \SI{10,0}{mL}$ de $S_1$ à la pipette jaugée, on les verse dans un erlenmeyer, on acidifie avec \SI{20}{mL} d'acide sulfurique \SI{1}{mol.L^{-1}} et on dose par une solution de permanganate de potassium de concentration $C = \SI{2,00e-2}{mol.L^{-1}}$. L'équivalence est obtenue pour $V_E = \SI{12,4}{mL}$. L'étiquette du flacon $S_0$ annonce $C_{\text{étq}} = \SI{1,20}{mol.L^{-1}}$.
\begin{enumerate}
\item Écrire l'équation bilan de la réaction de dosage.
\item Décrire le repérage de l'équivalence.
\item Calculer la concentration $C_1$ de $S_1$, puis la concentration $C_0$ de $S_0$.
\item Vérifier la conformité de l'étiquette.
\end{enumerate}
\tcblower
\textbf{Étape 1 — Identification.}
Espèce titrée : \ce{Fe^2+} (dans l'erlenmeyer, solution $S_1$, concentration $C_1$ inconnue). Couple : \ce{Fe^3+}/\ce{Fe^2+}.
Espèce titrante : \ce{MnO4-} (dans la burette, concentration $C = \SI{2,00e-2}{mol.L^{-1}}$ connue). Couple : \ce{MnO4-}/\ce{Mn^2+}.

\textbf{Étape 2 — Équation bilan.}
Milieu acide (\ce{H2SO4} apportant \ce{H+}).
Demi-équation réduction : $\ce{MnO4- + 8H+ + 5e- -> Mn^2+ + 4H2O}$
Demi-équation oxydation : $\ce{Fe^2+ -> Fe^3+ + e-}$ (×5)
Équation bilan :
\[ \ce{MnO4- + 5Fe^2+ + 8H+ -> Mn^2+ + 5Fe^3+ + 4H2O} \]
Réaction totale, rapide, unique : convenable pour un dosage.

\textbf{Étape 3 — Repérage de l'équivalence.}
Avant l'équivalence, les ions \ce{MnO4-} violets versés sont entièrement consommés par les ions \ce{Fe^2+} ; la solution reste jaune pâle (couleur \ce{Fe^3+}). À l'équivalence, la première goutte de \ce{MnO4-} en excès n'est plus consommée et donne à la solution une coloration \textbf{rose pâle persistante} (durant au moins 30 s). C'est le repère visuel. On note le volume $V_E = \SI{12,4}{mL}$.

\textbf{Étape 4 — Relation à l'équivalence.}
D'après les coefficients de l'équation bilan ($a=5$ pour \ce{Fe^2+}, $b=1$ pour \ce{MnO4-}) :
\[ \frac{n_{\ce{Fe^2+}}}{5} = n_{\ce{MnO4-}} \quad \Longrightarrow \quad \frac{C_1 \times V_1}{5} = C \times V_E \]

\textbf{Étape 5 — Calculs.}
Calcul de $C_1$ (concentration de la solution diluée $S_1$) :
\[ C_1 = \frac{5 \times C \times V_E}{V_1} = \frac{5 \times \SI{2,00e-2}{mol.L^{-1}} \times \SI{12,4}{mL}}{\SI{10,0}{mL}} = \SI{0,124}{mol.L^{-1}} \]
(3 chiffres significatifs cohérents avec les données).

Facteur de dilution : $F = \dfrac{V_f}{V_0} = \dfrac{\SI{200,0}{mL}}{\SI{20,0}{mL}} = 10$.

Concentration de la solution mère $S_0$ :
\[ C_0 = F \times C_1 = 10 \times \SI{0,124}{mol.L^{-1}} = \SI{1,24}{mol.L^{-1}} \]

\textbf{Étape 6 — Conclusion (conformité étiquette).}
Valeur mesurée : $C_{\text{mes}} = \SI{1,24}{mol.L^{-1}}$. Valeur étiquette : $C_{\text{étq}} = \SI{1,20}{mol.L^{-1}}$.
Écart relatif :
\[ \varepsilon = \frac{|1,24 - 1,20|}{1,20} \times 100 = \frac{0,04}{1,20} \times 100 \approx 3,3\,\% \]
Comme $\varepsilon = 3,3\,\% < 5\,\%$, \textbf{l'étiquette est conforme}. La solution $S_0$ est bien titrée à \SI{1,20}{mol.L^{-1}} dans les limites des incertitudes expérimentales.
\end{exoresolu}

\courssub{Sources d'erreur à connaître et à citer}

Même avec une méthode parfaite, la mesure expérimentale comporte des imperfections. Sache citer les principales \cle{sources d'erreur} et leur effet sur le résultat (surestimation ou sous-estimation de $C_{\text{titré}}$) :

\begin{itemize}
\item \textbf{Dépassement de l'équivalence (goutte en trop)} : $V_E$ mesuré $>$ $V_E$ réel $\Rightarrow$ $n_{\text{titrant}}$ surestimé $\Rightarrow$ $C_{\text{titré}}$ calculée \textbf{trop grande}. \emph{Parade : dosage rapide puis goutte à goutte près de l'équivalence.}
\item \textbf{Erreur de lecture de la burette (parallaxe)} : Œil pas au niveau du bas du ménisque. Lecture trop haute $\Rightarrow$ $V_E$ surestimé $\Rightarrow$ $C_{\text{titré}}$ \textbf{trop grande}. Lecture trop basse $\Rightarrow$ $C_{\text{titré}}$ \textbf{trop faible}.
\item \textbf{Perte de solution titrée} : Gouttes projetées hors de l'erlenmeyer lors de l'agitation, ou solution restée sur les parois non rincées. \textbf{Aucune incidence} sur le résultat si la perte a lieu \emph{après} le prélèvement pipeté (le nombre de moles dans l'erlenmeyer ne change pas). En revanche, si on perd du titré \emph{avant} le prélèvement (ex. rinçage pipette mal fait), $C_{\text{titré}}$ est faussée.
\item \textbf{Solution titrante mal conservée / altérée} :
      \begin{itemize}
      \item \ce{S2O3^2-} : s'oxyde à l'air en \ce{S4O6^2-} $\Rightarrow$ $C_{\text{titrant}}$ réelle $<$ $C$ annoncée $\Rightarrow$ $V_E$ plus grand $\Rightarrow$ $C_{\text{titré}}$ calculée \textbf{trop grande}.
      \item \ce{MnO4-} : se décompose à la lumière (réduction par la matière organique/eau) en \ce{MnO2} (marron) $\Rightarrow$ $C_{\text{titrant}}$ réelle $<$ $C$ annoncée $\Rightarrow$ $C_{\text{titré}}$ calculée \textbf{trop grande}.
      \end{itemize}
\item \textbf{Erreur sur le facteur de dilution} : Trait de jauge dépassé (volume $V_f$ réel $>$ nominal $\Rightarrow$ $F$ réel $>$ $F$ théorique $\Rightarrow$ $C_0$ calculée \textbf{trop grande}), pipette mal rincée (restes d'eau $\Rightarrow$ dilution involontaire $\Rightarrow$ $C_0$ \textbf{trop faible}).
\item \textbf{Oubli de l'acidification (pour \ce{MnO4-})} : En milieu neutre/basique, \ce{MnO4-} se réduit en \ce{MnO2} (marron, 3 e-) au lieu de \ce{Mn^2+} (5 e-). La stœchiométrie change ($5/3$ au lieu de $5/1$) $\Rightarrow$ résultat \textbf{totalement faux}.
\end{itemize}

\begin{savaistu}[Le dosage, star des épreuves du Bac TI]
Au Baccalauréat TI, l'exercice de dosage d'oxydoréduction rapporte presque chaque année entre 4 et 6 points sur 20, et il suit \emph{exactement} la trame en six étapes que tu viens d'apprendre : équation bilan (1 pt), relation à l'équivalence (1 pt), calcul de concentration (1 à 2 pts), prise en compte de la dilution (1 pt) et conclusion avec écart relatif (1 pt). Les correcteurs récompensent la \emph{rédaction soignée} : une relation littérale bien posée \emph{avant} l'application numérique, des unités à chaque ligne, et une phrase de conclusion complète (« L'écart relatif est de X\%, inférieur à 5\%, l'étiquette est donc conforme »). En maîtrisant cette méthode, tu sécurises une part précieuse et prévisible de ta note de chimie. C'est aussi la démarche utilisée quotidiennement dans les laboratoires de contrôle qualité de la SONARA (raffinerie), des brasseries, des savonneries artisanales (dosage de la soude) ou des industries agroalimentaires : vérifier qu'un produit est conforme à ce qu'il annonce est un geste professionnel fondamental.
\end{savaistu}

\begin{aretenir}[L'essentiel de la section — Fiche de révision Bac]
\begin{itemize}
\item Tout exercice de dosage se résout en \textbf{six étapes} immuables : \textbf{1. Identifier} (titré/titrant/couples), \textbf{2. Équilibrer} (demi-équations $\rightarrow$ bilan), \textbf{3. Repérer} (couleur + explication), \textbf{4. Relier} ($n_{\text{titré}}/a = n_{\text{titrant}}/b$), \textbf{5. Calculer} ($C$, dilution $F$, sig. fig.), \textbf{6. Conclure} ($\varepsilon$ vs 5\%).
\item À l'équivalence : $\dfrac{n_{\text{titré}}}{a} = \dfrac{n_{\text{titrant}}}{b}$, où $a$ et $b$ sont les coefficients stœchiométriques de l'équation bilan (\textbf{ne jamais écrire $n_{\text{titré}} = n_{\text{titrant}}$ sans vérifier les coefficients}).
\item \textbf{Auto-indicateur \ce{MnO4-}} : passage incolore/jaune $\rightarrow$ \textbf{rose pâle persistant}. \textbf{Indicateur amidon/\ce{I2}} : passage \textbf{bleu foncé} $\rightarrow$ \textbf{incolore} (ajout amidon en fin de dosage).
\item Si dilution préalable : $C_{\text{mère}} = F \times C_{\text{diluée}}$ avec $F = V_f/V_0$. \textbf{Oublier $F$ est l'erreur n°1 au Bac.}
\item Conformité étiquette : $\varepsilon = \dfrac{|C_{\text{mes}} - C_{\text{étq}}|}{C_{\text{étq}}} \times 100$. $\varepsilon \le 5\,\% \Rightarrow$ conforme ; $\varepsilon > 5\,\% \Rightarrow$ non conforme.
\item Sources d'erreur principales : dépassement équivalence ($C \uparrow$), parallaxe burette, titrant altéré ($C \uparrow$), oubli acidification \ce{MnO4-} (stœchiométrie fausse), erreur facteur dilution.
\end{itemize}
\end{aretenir}

\newpage

\coursec{Travaux pratiques}

\coursec{Leçon 11 : Dosages d'oxydoréduction}

\courssub{1. Généralités sur les dosages}

\begin{definition}[Dosage, solution titrée, solution titrante, équivalence]
Un \cle{dosage} est une opération qui permet de déterminer la concentration d'une espèce chimique (la \cle{solution titrée}) en la faisant réagir avec une solution de concentration connue (la \cle{solution titrante}) versée goutte à goutte depuis une burette.
Le \cle{point d'équivalence} (noté $E$) est atteint lorsque les réactifs ont été introduits dans les proportions stœchiométriques exactes de l'équation de la réaction de dosage. On note $V_E$ le volume de solution titrante versé à l'équivalence.
\end{definition}

\begin{propriete}[Conditions d'un bon dosage d'oxydoréduction]
Pour qu'un dosage soit exploitable avec précision, la réaction mise en jeu doit satisfaire quatre conditions fondamentales :
\begin{enumerate}
\item \textbf{Rapidité :} La réaction est quasi instantanée à température ambiante (pas de chauffage nécessaire, pas d'attente).
\item \textbf{Caractère complet :} L'équilibre de la réaction est fortement déplacé vers les produits (constante d'équilibre $K$ très grande), ce qui garantit une variation brutale de la grandeur suivie (couleur, potentiel) au voisinage de l'équivalence.
\item \textbf{Unicité :} Une seule réaction se produit entre le titré et le titrant (pas de réactions parasites).
\item \textbf{Repérage de l'équivalence :} Il existe un moyen sensible de détecter $E$ : changement de couleur brutal (auto-indication ou indicateur coloré), saut de potentiel (potentiométrie), ou conductimétrie.
\end{enumerate}
\end{propriete}

\begin{aretenir}[Vocabulaire du Bac TI]
\begin{itemize}
\item \cle{Solution titrée} : solution de concentration inconnue (dans l'erlenmeyer).
\item \cle{Solution titrante} : solution de concentration connue $C_0$ (dans la burette).
\item \cle{Volume prélevé} $V_P$ : volume de solution titrée pipeté.
\item \cle{Volume à l'équivalence} $V_E$ : volume de solution titrante versé à l'équivalence.
\item \cle{Auto-indication} : l'un des réactifs (ou un produit) est intensément coloré ; sa disparition ou son apparition signale l'équivalence sans ajout d'indicateur.
\end{itemize}
\end{aretenir}

\courssub{2. Conditions d'un bon dosage et dispositif expérimental}

\begin{methode}[Montage et préparation du matériel pour un dosage]
\begin{enumerate}
\item \textbf{Choix de la verrerie :} Burette graduée (classe B, \SI{25}{mL} ou \SI{50}{mL}), pipette jaugée (\SI{10}{mL}, \SI{20,0}{mL} ou \SI{25,0}{mL}), erlenmeyer (\SI{125}{mL} ou \SI{250}{mL}), propipette (ou poire), bécher, entonnoir, support et pince.
\item \textbf{Rinçage de la burette :} 2 à 3 rinçages avec la solution titrante (quelques mL à chaque fois, faire couler par le robinet). \emph{But :} éliminer l'eau résiduelle qui diluerait la solution titrante et fausserait $C_0$.
\item \textbf{Rinçage de la pipette :} 2 à 3 rinçages avec la solution titrée. \emph{But :} éviter la dilution du volume prélevé $V_P$.
\item \textbf{Remplissage de la burette :} Verser la solution titrante au-dessus de la graduation zéro, puis ouvrir le robinet pour chasser les bulles d'air dans la pointe et ajuster le ménisque sur le zéro (lecture en bas du ménisque pour les liquides clairs, \textbf{en haut du ménisque} pour le permanganate violet foncé).
\item \textbf{Prélèvement $V_P$ :} Aspirer la solution titrée au-dessus du trait de jauge, ajuster le ménisque sur le trait (œil au niveau du trait), laisser s'écouler dans l'erlenmeyer, \textbf{ne pas souffler} la dernière goutte (la pipette est étalonnée "à livrer").
\item \textbf{Acidification / Milieu :} Ajouter le réactif nécessaire (ex. \ce{H2SO4} pour \ce{MnO4-}) et de l'eau distillée si besoin (ne change pas $n_{\text{titré}}$).
\end{enumerate}
\end{methode}

\begin{popfigure}
\begin{center}
\begin{tikzpicture}[scale=0.9, every node/.style={font=\small}]
% Support
\draw[very thick, popdark] (-0.2,0) -- (3.4,0);
\draw[very thick, popdark] (0.2,0) -- (0.2,7.2);
% Pince
\draw[very thick, popdark] (0.2,5.6) -- (1.5,5.6);
\draw[very thick, popdark] (1.5,5.35) rectangle (1.9,5.85);
% Burette
\draw[thick, popblue] (1.6,7.0) -- (1.6,2.6);
\draw[thick, popblue] (1.8,7.0) -- (1.8,2.6);
\draw[thick, popblue] (1.6,2.6) -- (1.68,2.2) -- (1.72,2.2) -- (1.8,2.6);
% Robinet
\draw[thick, popdark] (1.5,2.5) -- (1.9,2.5);
\fill[popdark] (1.7,2.5) circle (0.05);
% Liquide violet dans la burette
\fill[poppurple!60] (1.62,6.4) rectangle (1.78,2.62);
\draw[popmuted] (1.62,6.4) -- (1.78,6.4);
% Graduations
\foreach \y in {3.2,3.8,4.4,5.0,5,6,6.2} \draw[popmuted] (1.6,\y) -- (1.68,\y);
% Gouttes
\fill[poppurple] (1.7,1.9) circle (0.06);
\fill[poppurple] (1.7,1.55) circle (0.05);
% Erlenmeyer
\draw[thick, popblue] (1.05,0.1) -- (1.45,1.0) -- (1.45,1.35);
\draw[thick, popblue] (2.35,0.1) -- (1.95,1.0) -- (1.95,1.35);
\draw[thick, popblue] (1.45,1.35) -- (1.45,1.5);
\draw[thick, popblue] (1.95,1.35) -- (1.95,1.5);
\draw[thick, popblue] (1.05,0.1) -- (2.35,0.1);
% Liquide dans l'erlenmeyer (jaune/vert pâle)
\fill[popgreen!25] (1.14,0.12) -- (1.42,0.75) -- (1.98,0.75) -- (2.26,0.12) -- cycle;
% Agitateur magnétique (optionnel)
\draw[popmuted, thick] (1.7,0.2) ellipse (0.15 and 0.05);
% Légendes
\node[align=left, anchor=west, popdark] at (3.7,5.0) {Burette : solution titrante\\ (ex. \ce{KMnO4}, $C_0$ connu)};
\node[align=left, anchor=west, popdark] at (3.7,0.6) {Erlenmeyer : $V_P$ solution titrée\\ + milieu réactionnel + indicateur};
\draw[->, popmuted] (3.65,5.0) -- (2.0,5.0);
\draw[->, popmuted] (3.65,0.6) -- (2.2,0.5);
\end{tikzpicture}
\end{center}
\legende{Dispositif général de dosage (ici illustration pour \ce{Fe^2+}/\ce{MnO4-}).}
\end{popfigure}

\courssub{3. Dosage des ions fer(II) par les ions permanganate en milieu acide}

\textbf{Contexte :} Vérification de la teneur en fer(II) d'une solution de sulfate de fer(II) \ce{FeSO4} (ex. contrôle qualité engrais, traitement eaux, pharmacopée). La réaction est une oxydoréduction en milieu acide fort (\ce{H2SO4}).

\begin{propriete}[Équations de la réaction de dosage]
\begin{itemize}
\item Couple oxydant : \ce{MnO4-}/\ce{Mn^2+} \quad \ce{MnO4- + 8H+ + 5e- -> Mn^2+ + 4H2O} \quad ($E^\circ = \SI{+1,51}{V}$)
\item Couple réducteur : \ce{Fe^3+}/\ce{Fe^2+} \quad \ce{Fe^2+ -> Fe^3+ + e-} \quad ($E^\circ = \SI{+0,77}{V}$)
\item \textbf{Équation globale du dosage} (électrons simplifiés, $\times 5$ pour le couple fer) :
\[ \ce{MnO4- + 5Fe^2+ + 8H+ -> Mn^2+ + 5Fe^3+ + 4H2O} \]
\end{itemize}
\end{propriete}

\begin{attention}[Pourquoi l'acide sulfurique \ce{H2SO4} et pas \ce{HCl} ni \ce{HNO3} ?]
\begin{itemize}
\item \ce{HCl} : les ions \ce{Cl-} sont oxydés par \ce{MnO4-} en \ce{Cl2} (réaction parasite) : $\ce{2Cl- -> Cl2 + 2e-}$ ($E^\circ = \SI{+1,36}{V}$).
\item \ce{HNO3} : agent oxydant lui-même (réduction en \ce{NO2} ou \ce{NO}), fausse le dosage.
\item \ce{H2SO4} : anion \ce{SO4^2-} non oxydable, non réducteur, fournit les \ce{H+} nécessaires (8 par \ce{MnO4-}) et empêche la précipitation de \ce{MnO2} brun (oxydation incomplète en milieu neutre/basique).
\end{itemize}
\end{attention}

\begin{aretenir}[Auto-indication par le permanganate]
Les ions \ce{MnO4-} sont \textbf{violets intenses} ; les produits \ce{Mn^2+} (incolore, rose très pâle très dilué), \ce{Fe^3+} (jaune pâle) et \ce{Fe^2+} (vert pâle) le sont beaucoup moins.
\begin{itemize}
\item \textbf{Avant $E$ :} \ce{MnO4-} est consommé instantanément $\rightarrow$ solution se décolore.
\item \textbf{À $E$ :} La première goutte de \ce{KMnO4} en excès n'est plus consommée $\rightarrow$ coloration \textbf{rose pâle persistante} (au moins \SI{30}{s}) dans tout l'erlenmeyer.
\end{itemize}
\textbf{Aucun indicateur n'est nécessaire.}
\end{aretenir}

\courssub{Objectifs du TP (Fe2+/MnO4-)}

À l'issue de cette séance, tu seras capable de :
\begin{itemize}
\item \cle{Schématiser} et \cle{monter} le dispositif expérimental d'un dosage d'oxydoréduction ;
\item \cle{Doser} une solution de sulfate de fer(II) par une solution de permanganate de potassium de concentration connue ;
\item \cle{Repérer} le point d'équivalence grâce à la persistance de la coloration violette (rose pâle) des ions permanganate ;
\item \cle{Calculer} la concentration molaire de la solution d'ions \ce{Fe^2+} et \cle{vérifier} l'indication portée sur l'étiquette du flacon ;
\item Adopter les \cle{gestes de sécurité} adaptés à la manipulation des produits oxydants et acides.
\end{itemize}

\courssub{Matériel et produits}

\begin{center}
\begin{tabularx}{\linewidth}{|l|X|}
\hline
\rowcolor{popblueL} \textbf{Matériel} & \textbf{Produits} \\
\hline
Burette graduée de \SI{25}{mL} ou \SI{50}{mL} & Solution de sulfate de fer(II) \ce{FeSO4} à doser (étiquette : $C \approx \SI{0,10}{mol.L^{-1}}$) \\
\hline
Pipette jaugée de \SI{20,0}{mL} + propipette (ou poire) & Solution de permanganate de potassium \ce{KMnO4} titrée, $C_0 = \SI{0,020}{mol.L^{-1}}$ (préparée par l'enseignant) \\
\hline
Erlenmeyer de \SI{250}{mL} & Acide sulfurique dilué \ce{H2SO4} à \SI{1}{mol.L^{-1}} \\
\hline
Éprouvette graduée de \SI{25}{mL} & Eau distillée \\
\hline
Support, pince, agitateur magnétique (ou agitation manuelle), bécher de \SI{100}{mL}, entonnoir & Gants, lunettes de protection, blouse \\
\hline
\end{tabularx}
\end{center}

\begin{attention}[Si un produit ou un matériel manque]
\begin{itemize}
\item Pas de propipette ? Aspiration \textbf{interdite à la bouche} : utilise une poire en caoutchouc ou demande à l'enseignant de pipetter.
\item Pas de burette ? Une pipette graduée de \SI{10}{mL} fixée verticalement peut dépanner, mais la précision sera moindre : à signaler dans la conclusion.
\item Le sulfate de fer(II) s'oxyde à l'air : la solution doit être \textbf{préparée la veille au plus tôt} et acidifiée. Une solution brunâtre (ions \ce{Fe^3+}) donnera un résultat trop faible.
\end{itemize}
\end{attention}

\courssub{Sécurité}

\begin{attention}[Pictogrammes et gestes à adopter]
\begin{itemize}
\item \textbf{Acide sulfurique dilué} : pictogramme \emph{corrosif} (liquide attaquant une main et une barre métallique). Port obligatoire des \textbf{lunettes} et des \textbf{gants}. En cas de contact avec la peau : rincer abondamment à l'eau froide pendant plusieurs minutes et prévenir l'enseignant.
\item \textbf{Permanganate de potassium} : pictogramme \emph{comburant} (flamme sur un cercle). Il tache durablement la peau et les vêtements (taches brunes de \ce{MnO2}) : manipuler au-dessus de la paillasse, jamais au-dessus de son cahier.
\item \textbf{Règles générales} : blouse fermée, cheveux attachés, ne jamais pipetter à la bouche, ne jamais goûter un produit, laver la verrerie après usage. Les solutions de \ce{KMnO4} usagées se versent dans le bidon de récupération prévu, pas dans l'évier.
\end{itemize}
\end{attention}

\courssub{Protocole expérimental}

\begin{methode}[Déroulement du dosage \ce{Fe^2+}/\ce{MnO4-}]
\begin{enumerate}
\item \textbf{Rinçage et remplissage de la burette} : rincer la burette avec un peu de solution de \ce{KMnO4} (2 fois \SI{5}{mL}), puis la remplir au-delà du zéro. Ajuster le zéro en ouvrant le robinet (chasser les bulles). Lire la graduation \textbf{en haut} du ménisque (liquide trop foncé). Noter $V_i$.
\item \textbf{Prélèvement} : à l'aide de la pipette jaugée munie de la propipette, prélever $V_P = \SI{20,0}{mL}$ de la solution de sulfate de fer(II) et les introduire dans l'erlenmeyer.
\item \textbf{Acidification} : ajouter environ \SI{10}{mL} d'acide sulfurique dilué \SI{1}{mol.L^{-1}} mesurés à l'éprouvette. \emph{Rôle :} fournir les \ce{H+} (8 par \ce{MnO4-}) et éviter la formation de \ce{MnO2} brun.
\item \textbf{Dosage rapide (essai)} : verser la solution de \ce{KMnO4} mL par mL en agitant vigoureusement. La coloration violette disparaît d'abord instantanément. Noter le volume $V_1$ où elle devient persistante (premier rose qui ne part pas).
\item \textbf{Dosage précis} : recommencer avec un nouveau prélèvement ($V_P$ identique). Verser rapidement jusqu'à $V_1 - \SI{2}{mL}$, puis \textbf{goutte à goutte} en agitant constamment jusqu'à la coloration \textbf{rose pâle persistante} (au moins 30 s). Noter $V_f$ et calculer $V_E = V_f - V_i$ à \SI{0,05}{mL} près.
\item \textbf{Concordance} : effectuer un deuxième dosage précis (et un troisième si nécessaire). Les volumes $V_E$ sont concordants s'ils diffèrent de moins de \SI{0,20}{mL}. Conserver la moyenne des volumes concordants.
\end{enumerate}
\end{methode}

\courssub{Observations et résultats attendus}

Dans l'erlenmeyer, la solution initialement verdâtre (ions \ce{Fe^2+}) jaunit progressivement (formation des ions \ce{Fe^3+}). Chaque goutte violette de permanganate se décolore tant que des ions \ce{Fe^2+} restent présents. À l'équivalence, la première goutte en excès colore durablement la solution en \textbf{rose pâle}.

\begin{center}
\begin{tabularx}{\linewidth}{|l|X|X|X|}
\hline
\rowcolor{popblueL} \textbf{Essai} & \textbf{Dosage 1 (rapide)} & \textbf{Dosage 2 (précis)} & \textbf{Dosage 3 (précis)} \\
\hline
$V_E$ (mL) & 19,8 & 19,55 & 19,50 \\
\hline
\end{tabularx}
\end{center}

Volume retenu (moyenne des dosages concordants 2 et 3) : $V_E = \SI{19,5}{mL}$.

\courssub{4. Dosage du diiode par les ions thiosulfate}

\textbf{Contexte :} Dosage de l'iode dans une teinture d'iode (désinfectant, pharmacopée), ou dosage indirect d'un oxydant (ex. \ce{Cl2}, \ce{Cr2O7^2-}, \ce{Cu^2+}) par libération d'iode (iodométrie).

\begin{propriete}[Équation de la réaction de dosage]
\begin{itemize}
\item Couple oxydant : \ce{I2}/\ce{I-} \quad \ce{I2 + 2e- -> 2I-} \quad ($E^\circ = \SI{+0,54}{V}$)
\item Couple réducteur : \ce{S4O6^2-}/\ce{S2O3^2-} \quad \ce{2S2O3^2- -> S4O6^2- + 2e-} \quad ($E^\circ = \SI{+0,08}{V}$)
\item \textbf{Équation globale} :
\[ \ce{I2 + 2S2O3^2- -> 2I- + S4O6^2-} \]
\end{itemize}
La réaction est rapide, complète et unique en milieu neutre ou faiblement acide.
\end{propriete}

\begin{aretenir}[Repérage de l'équivalence : indicateur amidon]
Le diiode \ce{I2} forme avec l'amidon (empois d'amidon, préparable localement à partir de manioc ou maïs) un complexe \textbf{bleu foncé} très intense, visible à très faible concentration.
\begin{itemize}
\item \textbf{Au cours du dosage :} La solution est jaune/brune (\ce{I2}). Le thiosulfate la décolore progressivement.
\item \textbf{Problème :} La décoloration totale est difficile à voir à l'œil nu (jaune très pâle $\rightarrow$ incolore).
\item \textbf{Solution :} On ajoute \textbf{quelques gouttes d'empois d'amidon} \textbf{vers la fin} (quand la couleur jaune devient très pâle, paille). Le complexe bleu apparaît/brutalement disparaît.
\item \textbf{À l'équivalence :} \textbf{Disparition brutale de la coloration bleue} $\rightarrow$ solution incolore.
\end{itemize}
\textbf{Pourquoi en fin de dosage ?} Le complexe \ce{I2}-amidon est très stable ; s'il est formé trop tôt, il ne libère pas tout l'\ce{I2} pour réagir avec le thiosulfate $\rightarrow$ erreur systématique (volume $V_E$ trop petit).
\end{aretenir}

\begin{methode}[Protocole type dosage \ce{I2}/\ce{S2O3^2-}]
\begin{enumerate}
\item Burette : solution de thiosulfate de sodium \ce{Na2S2O3} titrée $C_0$ (ex. \SI{0,05}{mol.L^{-1}}). Rinçage, remplissage, zéro.
\item Erlenmeyer : Pipeter $V_P$ de solution de diiode (teinture d'iode diluée). Ajouter \SI{20}{mL} eau distillée.
\item Dosage rapide : verser \ce{S2O3^2-} jusqu'à couleur jaune paille.
\item Ajout indicateur : verser \SI{1}{mL} d'empois d'amidon $\rightarrow$ coloration bleu foncé.
\item Dosage précis : goutte à goutte jusqu'à \textbf{décoloration complète (incolore)}. Noter $V_E$.
\item Concordance : 2 volumes concordants ($\Delta V_E < \SI{0,20}{mL}$).
\end{enumerate}
\end{methode}

\begin{savaistu}[Application locale : Teinture d'iode au Cameroun]
La teinture d'iode (solution d'\ce{I2} dans \ce{KI} alcoolique) est un antiseptique courant dans les pharmacies de Douala, Yaoundé ou les centres de santé ruraux. Son dosage par thiosulfate permet de vérifier la teneur en iode actif (norme pharmacopée : \SI{10}{\%} m/V environ). L'empois d'amidon peut être préparé au laboratoire en faisant bouillir \SI{1}{g} d'amidon de manioc (fécule) dans \SI{100}{mL} d'eau, puis en filtrant. C'est un exemple concret de chimie analytique au service de la santé publique.
\end{savaistu}

\courssub{5. Exploitation d'un dosage : méthode complète}

\begin{methode}[Méthode rigoureuse d'exploitation (à appliquer à tout dosage)]
\begin{enumerate}
\item \textbf{Écrire l'équation de la réaction de dosage} (bilan matière, charges équilibrées).
\item \textbf{Construire le tableau d'avancement} (État initial, État intermédiaire, État final à l'équivalence).
\begin{center}
\begin{tabular}{|c|c|c|c|c|}
\hline
 & \ce{MnO4-} & \ce{Fe^2+} & \ce{H+} & \ce{Mn^2+} / \ce{Fe^3+} \\
\hline
État initial & $n_0(\ce{MnO4-}) = C_0 V_E$ & $n_0(\ce{Fe^2+}) = C_X V_P$ & Excès & 0 \\
\hline
État équivalence & 0 & 0 & Excès & Produits formés \\
\hline
\end{tabular}
\end{center}
\item \textbf{Écrire la relation à l'équivalence} (réactifs introduits dans les proportions stœchiométriques) :
\[ \frac{n_{\text{titré}}(\text{initial})}{\nu_{\text{titré}}} = \frac{n_{\text{titrant}}(\text{versé à }E)}{\nu_{\text{titrant}}} \]
Soit $C_X V_P = \frac{\nu_{\text{titré}}}{\nu_{\text{titrant}}} \times C_0 V_E$.
\item \textbf{Calculer la concentration inconnue $C_X$} (attention aux unités : volumes en L ou mL cohérents).
\item \textbf{Calculer l'incertitude / écart relatif} et \textbf{conclure} (validation étiquette, pureté, conformité).
\end{enumerate}
\end{methode}

\begin{exoresolu}[Exploitation guidée : Dosage \ce{Fe^2+}/\ce{MnO4-}]
\textbf{Énoncé.} On dose $V_P = \SI{20,0}{mL}$ d'une solution de sulfate de fer(II) par une solution de \ce{KMnO4} de concentration $C_0 = \SI{0,020}{mol.L^{-1}}$. L'équivalence est obtenue pour $V_E = \SI{19,5}{mL}$. L'étiquette du flacon indique $C_{\text{étiquette}} = \SI{0,10}{mol.L^{-1}}$.
\begin{enumerate}
\item Écrire les demi-équations des couples \ce{MnO4-}/\ce{Mn^2+} et \ce{Fe^3+}/\ce{Fe^2+}, puis l'équation du dosage.
\item Établir la relation à l'équivalence et calculer $C(\ce{Fe^2+})$.
\item Comparer avec l'étiquette et conclure.
\end{enumerate}
\tcblower
\textbf{Solution.}
\textbf{1.} Couple \ce{MnO4-}/\ce{Mn^2+} : \ce{MnO4- + 8H+ + 5e- -> Mn^2+ + 4H2O}.
Couple \ce{Fe^3+}/\ce{Fe^2+} : \ce{Fe^2+ -> Fe^3+ + e-} (multiplier par 5).
D'où l'équation globale :
\[ \ce{MnO4- + 5Fe^2+ + 8H+ -> Mn^2+ + 5Fe^3+ + 4H2O} \]

\textbf{2.} À l'équivalence, les quantités de matière réagies sont proportionnelles aux coefficients stœchiométriques :
\[ n(\ce{Fe^2+})_{\text{initial}} = 5 \times n(\ce{MnO4-})_{\text{versé à }E} \]
\[ C(\ce{Fe^2+}) \times V_P = 5 \times C_0 \times V_E \]
\[ C(\ce{Fe^2+}) = \frac{5 \times C_0 \times V_E}{V_P} = \frac{5 \times \SI{0,020}{mol.L^{-1}} \times \SI{19,5}{mL}}{\SI{20,0}{mL}} \]
\[ C(\ce{Fe^2+}) = \SI{0,0975}{mol.L^{-1}} \]

\textbf{3.} Écart relatif par rapport à l'étiquette :
\[ E_r = \frac{|C_{\text{étiquette}} - C_{\text{exp}}|}{C_{\text{étiquette}}} \times 100 = \frac{|0,10 - 0,0975|}{0,10} \times 100 = 2,5\,\% \]
L'écart est faible (inférieur à 5\%), compatible avec l'incertitude expérimentale (pipette, burette, oxydation lente des \ce{Fe^2+} à l'air depuis la préparation de la solution).
\textbf{Conclusion :} L'indication de l'étiquette ($C \approx \SI{0,10}{mol.L^{-1}}$) est \textbf{validée}. La solution est conforme.
\end{exoresolu}

\begin{exercice}[Entraînement : Dosage d'une teinture d'iode]
On dose $V_P = \SI{10,0}{mL}$ d'une teinture d'iode (solution de diiode \ce{I2}) par une solution de thiosulfate de sodium \ce{Na2S2O3} de concentration $C_0 = \SI{0,100}{mol.L^{-1}}$.
L'équivalence est atteinte pour $V_E = \SI{12,4}{mL}$ (moyenne de deux valeurs concordantes).
La masse molaire de \ce{I2} est $M(\ce{I2}) = \SI{254}{g.mol^{-1}}$.
\begin{enumerate}
\item Rappeler l'équation de la réaction de dosage.
\item Calculer la concentration molaire $C(\ce{I2})$ de la teinture.
\item Calculer la concentration massique $\rho(\ce{I2})$ en \SI{}{g.L^{-1}} et en \%\ (m/V).
\item L'étiquette indique "Teinture d'iode à 10\% m/V". Le produit est-il conforme ? (On admet une tolérance de $\pm 5\%$ relatif).
\end{enumerate}
\end{exercice}

\begin{corrige}[Exercice : Dosage teinture d'iode]
\textbf{1.} $\ce{I2 + 2S2O3^2- -> 2I- + S4O6^2-}$.
\textbf{2.} Relation équivalence : $n(\ce{I2}) = \frac{1}{2} n(\ce{S2O3^2-})$.
$C(\ce{I2}) = \frac{C_0 \times V_E}{2 \times V_P} = \frac{0,100 \times 12,4}{2 \times 10,0} = \SI{0,0620}{mol.L^{-1}}$.
\textbf{3.} $\rho(\ce{I2}) = C \times M = 0,0620 \times 254 = \SI{15,75}{g.L^{-1}}$.
\% m/V = $\frac{15,75}{10} = \SI{1,575}{\%}$ (soit $\approx \SI{1,58}{\%}$ m/V).
\textbf{4.} Valeur étiquette : \SI{10}{\%} m/V. Valeur trouvée : \SI{1,58}{\%} m/V.
Écart relatif énorme ($> 80\%$). \textbf{Non conforme.} Hypothèse : la teinture a été diluée avant dosage (ex. dilution 1/10 ou 1/5 non mentionnée dans l'énoncé simplifié) ou l'étiquette fait référence à la teinture mère non diluée. En TP réel, vérifier le protocole de dilution initial.
\end{corrige}

\courssub{Conclusion et synthèse}

Ce chapitre t'a permis de maîtriser deux dosages d'oxydoréduction majeurs au programme TI :
\begin{itemize}
\item \textbf{\ce{Fe^2+}/\ce{MnO4-} en milieu acide :} Auto-indication (rose pâle), acidification \ce{H2SO4} impérative, stœchiométrie 1/5.
\item \textbf{\ce{I2}/\ce{S2O3^2-} :} Indicateur amidon (bleu $\rightarrow$ incolore), ajout en fin de dosage, stœchiométrie 1/2.
\end{itemize}

La démarche d'exploitation est universelle et doit être appliquée rigoureusement :
\textbf{Équation $\rightarrow$ Tableau d'avancement $\rightarrow$ Relation équivalence $\rightarrow$ Calcul $C_X$ $\rightarrow$ Incertitude/Comparaison $\rightarrow$ Conclusion.}

Ces techniques sont au cœur du contrôle qualité industriel (CIMENCAM, SONARA, brasseries, savonneries, laboratoires pharmaceutiques) et de la surveillance de l'environnement (eaux, sols). La rigueur du geste technique (pipetage, lecture ménisque, concordance) conditionne la fiabilité du résultat. Entraîne-toi sur les sujets de Bac blancs !

\newpage

\coursec{Méthodes et exercices résolus}

\courssub{Généralités sur les dosages}

\begin{definition}[Vocabulaire du dosage]
Un \cle{dosage} est une opération qui permet de déterminer la concentration d'une espèce chimique (la \cle{solution titrée}) en la faisant réagir totalement avec une solution de concentration connue (la \cle{solution titrante}).
\begin{itemize}
\item \textbf{Solution titrée} : solution de concentration inconnue, prélevée au volume précis $V_{\text{titré}}$ (pipette jaugée).
\item \textbf{Solution titrante} : solution de concentration connue $C_{\text{titrante}}$, versée depuis la burette graduée.
\item \textbf{Point d'équivalence} : instant où les réactifs ont été introduits dans les proportions stœchiométriques exactes (avancement maximal $x_E$). On note $V_E$ le volume de titrant versé à l'équivalence.
\end{itemize}
\end{definition}

\begin{propriete}[Conditions d'un bon dosage d'oxydoréduction]
Pour qu'un dosage redox soit exploitable, la réaction doit satisfaire trois conditions fondamentales :
\begin{enumerate}
\item \textbf{Totalité} : la réaction va jusqu'au bout (constante d'équilibre $K$ très grande, $E^\circ_{\text{ox}} - E^\circ_{\text{red}} > \SI{0,3}{V}$ environ).
\item \textbf{Rapidité} : la réaction est quasi instantanée à température ambiante (sinon chauffage, ex. oxalate/permanganate à \SI{60}{\celsius}).
\item \textbf{Unicité} : pas de réaction parasite (choix du milieu, ex. $\ce{H2SO4}$ et non $\ce{HCl}$ avec $\ce{MnO4-}$).
\end{enumerate}
On doit pouvoir \textbf{repérer l'équivalence} visuellement (changement de couleur du réactif en excès ou indicateur coloré).
\end{propriete}

\courssub{Schématiser et décrire le dispositif expérimental du dosage}

\begin{methode}[Réaliser et décrire un montage de dosage d'oxydoréduction]
Pour schématiser et décrire correctement le dispositif, suis toujours la même démarche :
\begin{enumerate}
\item \textbf{Identifier les deux solutions} : la solution \cle{titrante} (concentration connue, réactif en excès à la fin) va dans la \textbf{burette graduée} ; la solution \cle{titrée} (concentration inconnue) va dans le \textbf{bécher} ou l'erlenmeyer.
\item \textbf{Prélever un volume précis} $V_{\text{titré}}$ de solution titrée avec une \textbf{pipette jaugée} (jamais une éprouvette : précision insuffisante).
\item \textbf{Placer la burette} au-dessus du bécher, fixée sur une potence, avec un agitateur magnétique si possible.
\item \textbf{Acidifier si nécessaire} : pour le dosage par \ce{MnO4-}, ajouter de l'acide sulfurique dilué (les ions \ce{H+} interviennent dans la demi-équation).
\item \textbf{Verser la solution titrante} mL par mL en agitant, puis goutte à goutte près de l'équivalence.
\item \textbf{Repérer l'équivalence} par le changement de couleur (auto-indicateur avec \ce{MnO4-} ou indicateur coloré comme l'empois d'amidon pour \ce{I2}).
\item \textbf{Lire le volume équivalent} $V_E$ au bas du ménisque, à $\pm \SI{0,05}{mL}$ près, et refaire un deuxième dosage concordant.
\end{enumerate}
Réflexe rédaction : dans la description, cite toujours \textbf{la verrerie de précision} (pipette jaugée, burette graduée) et \textbf{le critère de l'équivalence}.
\end{methode}

\begin{exoresolu}[Montage du dosage du fer(II) par le permanganate]
Un laborantin de la SONARA veut contrôler la concentration d'une solution de sulfate de fer(II). Il prélève $V_1 = \SI{20,0}{mL}$ de cette solution qu'il place dans un erlenmeyer, ajoute \SI{10}{mL} d'acide sulfurique dilué, puis dose par une solution de permanganate de potassium de concentration $C_2 = \SI{2,00e-2}{mol.L^{-1}}$ placée dans la burette.\\
\textbf{1.} Faire le schéma annoté du dispositif.\\
\textbf{2.} Justifier l'ajout d'acide sulfurique et le choix de la verrerie.
\tcblower
\textbf{1. Schéma annoté du dispositif}

\begin{popfigure}
\begin{tikzpicture}[scale=0.95, every node/.style={font=\small}]
% Burette
\draw[thick, popblue] (0,4.6) -- (0,1.6);
\draw[thick, popblue] (0.5,4.6) -- (0.5,1.6);
\draw[thick, popblue] (0,1.6) -- (0.18,1.2) -- (0.18,0.9);
\draw[thick, popblue] (0.5,1.6) -- (0.32,1.2) -- (0.32,0.9);
\draw[thick, popblue] (0.18,0.9) -- (0.32,0.9);
% robinet
\draw[thick, popdark] (0.05,1.35) -- (0.45,1.35);
\fill[popdark] (0.25,1.35) circle (0.05);
% liquide violet dans la burette
\fill[poppurple!40] (0.02,4.5) rectangle (0.48,2.2);
\draw[popmuted] (0.02,2.2) -- (0.48,2.2);
% graduations
\foreach \y in {2.4,2.8,3.2,3.6,4.0,4.4} \draw[popmuted] (0.38,\y) -- (0.48,\y);
% goutte
\fill[poppurple] (0.25,0.55) circle (0.06);
% erlenmeyer
\draw[thick, popblue] (-0.9,-1.6) -- (-0.25,-0.2) -- (-0.25,0.15);
\draw[thick, popblue] (1.4,-1.6) -- (0.75,-0.2) -- (0.75,0.15);
\draw[thick, popblue] (-0.9,-1.6) -- (1.4,-1.6);
\draw[thick, popblue] (-0.25,0.15) -- (0.75,0.15);
% liquide dans l'erlenmeyer
\fill[popgreen!25] (-0.78,-1.55) -- (1.28,-1.55) -- (0.98,-0.85) -- (-0.48,-0.85) -- cycle;
% barreau magnétique
\fill[popdark] (0.25,-1.45) ellipse (0.28 and 0.07);
% annotations
\node[align=left, poppurple, anchor=west] at (1.6,3.4) {Burette graduée :\\ solution titrante \ce{MnO4-}\\ ($C_2 = \SI{2,00e-2}{mol.L^{-1}}$)};
\draw[-{Stealth}, poppurple] (1.6,3.1) -- (0.55,3.0);
\node[align=left, popdark, anchor=west] at (1.6,1.35) {Robinet};
\draw[-{Stealth}, popdark] (1.6,1.35) -- (0.5,1.35);
\node[align=left, popgreen!60!black, anchor=west] at (1.6,-0.9) {Erlenmeyer :\\ $V_1 = \SI{20,0}{mL}$ de solution\\ titrée \ce{Fe^2+} + \ce{H2SO4} dilué};
\draw[-{Stealth}, popgreen!60!black] (1.6,-1.1) -- (1.0,-1.2);
\node[align=left, popmuted, anchor=east] at (-1.1,-1.45) {Barreau\\ aimanté};
\draw[-{Stealth}, popmuted] (-1.1,-1.3) -- (-0.1,-1.42);
\end{tikzpicture}
\legende{Dispositif de dosage des ions \ce{Fe^2+} par les ions permanganate \ce{MnO4-}}
\end{popfigure}

\textbf{2. Justifications}

\emph{Pourquoi acidifier ?} La demi-équation du couple \ce{MnO4-}/\ce{Mn^2+} s'écrit :
\[ \ce{MnO4- + 8H+ + 5e- -> Mn^2+ + 4H2O} \]
Les ions \ce{H+} sont des \textbf{réactifs} : sans acidification, la réaction ne se produit pas correctement (risque de formation de \ce{MnO2} brun au lieu de \ce{Mn^2+} incolore). On utilise l'acide \textbf{sulfurique} (et non l'acide chlorhydrique, dont les ions \ce{Cl-} pourraient être oxydés par \ce{MnO4-}).

\emph{Choix de la verrerie :}
\begin{itemize}
\item le volume $V_1 = \SI{20,0}{mL}$ est prélevé à la \textbf{pipette jaugée} (précision $\pm \SI{0,05}{mL}$) car il intervient dans le calcul de concentration ;
\item la solution titrante est placée dans la \textbf{burette graduée}, qui permet de mesurer le volume versé au dixième de millilitre ;
\item l'acide sulfurique est mesuré au \textbf{cylindre gradué} (ou versé directement) : il est en excès, son volume exact n'intervient pas dans le calcul.
\end{itemize}
\textbf{Conclusion :} le montage associe burette graduée (titrant \ce{MnO4-}), pipette jaugée (titré \ce{Fe^2+}) et milieu acidifié à l'acide sulfurique.
\end{exoresolu}

\courssub{Doser une solution oxydante de permanganate de potassium}

\begin{methode}[Exploiter un dosage d'oxydoréduction : la méthode générale]
Tout calcul de dosage repose sur la \textbf{stœchiométrie à l'équivalence}. Démarche en 5 étapes :
\begin{enumerate}
\item \textbf{Écrire les deux demi-équations} des couples mis en jeu, en équilibrant les électrons (milieu acide : équilibrer O avec \ce{H2O}, puis H avec \ce{H+}).
\item \textbf{Combiner} les demi-équations pour éliminer les électrons et obtenir l'\textbf{équation de la réaction de dosage} (elle doit être totale et rapide).
\item \textbf{Écrire la relation à l'équivalence} : les réactifs sont introduits dans les proportions stœchiométriques. Si l'équation est $a\,\text{Ox}_1 + b\,\text{Red}_2 \rightarrow$ produits, alors :
\[ \frac{n_1(\text{versé à l'équivalence})}{a} = \frac{n_2(\text{présent})}{b} \]
\item \textbf{Passer aux concentrations} : $n = C \times V$, d'où $C_{\text{inconnue}} = \dfrac{b}{a}\times\dfrac{C_{\text{titrante}}\,V_E}{V_{\text{titré}}}$. Attention aux \textbf{coefficients stœchiométriques} !
\item \textbf{Vérifier} : unités (volumes dans la même unité, elles se simplifient), chiffres significatifs (3 en général), ordre de grandeur plausible.
\end{enumerate}
Réflexe : en cas de doute, construis un \textbf{tableau d'avancement} à l'équivalence ($x_E$ rend nuls les deux réactifs).
\end{methode}

\begin{exoresolu}[Étalonage d'une solution de permanganate de potassium]
Pour vérifier la concentration d'une solution de permanganate de potassium préparée au laboratoire, on dose $V_{\text{ox}} = \SI{10,0}{mL}$ d'une solution étalon d'oxalate de sodium (\ce{Na2C2O4}) de concentration $C_{\text{red}} = \SI{5,00e-2}{mol.L^{-1}}$, acidifiée à l'acide sulfurique et chauffée vers \SI{60}{\celsius}. L'équivalence est atteinte pour un volume versé $V_E = \SI{10,4}{mL}$ de solution de permanganate.\\
Couples : \ce{MnO4-}/\ce{Mn^2+} et \ce{CO2}/\ce{C2O4^2-}.\\
\textbf{1.} Écrire l'équation de la réaction de dosage.\\
\textbf{2.} Déterminer la concentration $C_{\text{ox}}$ de la solution de permanganate.
\tcblower
\textbf{1. Équation de la réaction de dosage}

Demi-équations électroniques :
\begin{align*}
\ce{MnO4- + 8H+ + 5e-} &\ce{-> Mn^2+ + 4H2O} \quad (\times 2)\\
\ce{C2O4^2-} &\ce{-> 2CO2 + 2e-} \quad (\times 5)
\end{align*}
On multiplie la première par 2 et la seconde par 5 pour échanger le même nombre d'électrons (10), puis on additionne :
\[ \boxed{\ce{2MnO4- + 5C2O4^2- + 16H+ -> 2Mn^2+ + 10CO2 + 8H2O}} \]

\textbf{2. Calcul de la concentration}

Tableau d'avancement à l'équivalence (seules les lignes utiles) :
\begin{center}
\begin{tabularx}{\linewidth}{|l|X|X|X|}
\hline
\rowcolor{popblueL} État & $n(\ce{MnO4-})$ & $n(\ce{C2O4^2-})$ & Avancement \\
\hline
Initial & $C_{\text{ox}}V_E$ & $C_{\text{red}}V_{\text{ox}}$ & $0$ \\
\hline
Équivalence & $C_{\text{ox}}V_E - 2x_E = 0$ & $C_{\text{red}}V_{\text{ox}} - 5x_E = 0$ & $x_E$ \\
\hline
\end{tabularx}
\end{center}
À l'équivalence, les deux réactifs sont totalement consommés :
\[ \frac{n(\ce{MnO4-})}{2} = \frac{n(\ce{C2O4^2-})}{5} \quad\Longrightarrow\quad \frac{C_{\text{ox}}V_E}{2} = \frac{C_{\text{red}}V_{\text{ox}}}{5} \]
D'où :
\[ C_{\text{ox}} = \frac{2\,C_{\text{red}}\,V_{\text{ox}}}{5\,V_E} = \frac{2 \times \SI{5,00e-2}{} \times 10{,}0}{5 \times 10{,}4} = \SI{1,92e-2}{mol.L^{-1}} \]
\textbf{Conclusion :} la solution de permanganate de potassium a pour concentration $C_{\text{ox}} \approx \SI{1,92e-2}{mol.L^{-1}}$. Elle est maintenant \textbf{étalonnée} et peut servir de solution titrante pour d'autres dosages.
\end{exoresolu}

\courssub{Doser une solution d'ions \ce{Fe^2+} par les ions \ce{MnO4-}}

\begin{methode}[Dosage manganimétrique des ions fer(II)]
\begin{enumerate}
\item \textbf{Écrire l'équation de dosage} à partir des couples \ce{MnO4-}/\ce{Mn^2+} et \ce{Fe^3+}/\ce{Fe^2+} :
\[ \ce{MnO4- + 5Fe^2+ + 8H+ -> Mn^2+ + 5Fe^3+ + 4H2O} \]
(1 ion permanganate réagit avec \textbf{5} ions fer(II) : réflexe à mémoriser.)
\item \textbf{Acidifier} le prélèvement avec de l'acide sulfurique dilué.
\item \textbf{Repérer l'équivalence sans indicateur} : \ce{MnO4-} est \textbf{violet}, tous les autres ions sont quasi incolores (\ce{Fe^3+} est jaune pâle). Avant l'équivalence, le permanganate versé est consommé et se décolore ; \textbf{à l'équivalence, la première goutte en excès donne une coloration violette persistante}. On dit que \ce{MnO4-} joue le rôle d'\cle{auto-indicateur}.
\item \textbf{Appliquer la relation à l'équivalence} :
\[ n(\ce{Fe^2+}) = 5\,n(\ce{MnO4-})_E \quad\Longrightarrow\quad C(\ce{Fe^2+}) = \frac{5\,C(\ce{MnO4-})\,V_E}{V(\ce{Fe^2+})} \]
\item \textbf{Convertir si besoin} en concentration massique : $C_m = C \times M$.
\end{enumerate}
\end{methode}

\begin{exoresolu}[Contrôle d'un complément alimentaire en fer]
Un comprimé contre l'anémie, vendu en pharmacie à Douala, contient du sulfate de fer(II) \ce{FeSO4}. On dissout un comprimé dans de l'eau distillée acidifiée et on complète à $V_0 = \SI{100,0}{mL}$ (solution $S$). On prélève $V_1 = \SI{20,0}{mL}$ de $S$ que l'on dose par une solution de \ce{KMnO4} à $C_2 = \SI{1,00e-2}{mol.L^{-1}}$. La coloration violette persistante apparaît pour $V_E = \SI{12,0}{mL}$.\\
Données : $M(\ce{Fe}) = \SI{55,8}{g.mol^{-1}}$.\\
\textbf{1.} Écrire l'équation de la réaction de dosage.\\
\textbf{2.} Comment repère-t-on l'équivalence ?\\
\textbf{3.} Calculer la masse de fer contenue dans un comprimé. L'étiquette annonce \SI{65}{mg} : le comprimé est-il conforme (tolérance $\pm 5\,\%$) ?
\tcblower
\textbf{1. Équation de la réaction de dosage}

Couples : \ce{MnO4-}/\ce{Mn^2+} et \ce{Fe^3+}/\ce{Fe^2+}.
\begin{align*}
\ce{MnO4- + 8H+ + 5e-} &\ce{-> Mn^2+ + 4H2O}\\
\ce{Fe^2+} &\ce{-> Fe^3+ + e-} \quad (\times 5)
\end{align*}
\[ \boxed{\ce{MnO4- + 5Fe^2+ + 8H+ -> Mn^2+ + 5Fe^3+ + 4H2O}} \]

\textbf{2. Repérage de l'équivalence}

Avant l'équivalence, chaque goutte de \ce{MnO4-} (violet) versée est consommée par les ions \ce{Fe^2+} : la solution reste incolore à jaune pâle. Dès que tous les ions \ce{Fe^2+} ont disparu, la goutte de permanganate en excès n'est plus consommée : \textbf{l'équivalence est repérée par l'apparition d'une coloration violette persistante}. Le permanganate est son propre indicateur (auto-indicateur).

\textbf{3. Masse de fer dans un comprimé}

À l'équivalence :
\[ n(\ce{Fe^2+})_{\text{prélevé}} = 5\,n(\ce{MnO4-})_E = 5\,C_2 V_E = 5 \times \SI{1,00e-2}{} \times \SI{12,0e-3}{} = \SI{6,00e-4}{mol} \]
Cette quantité provient de \SI{20,0}{mL} de $S$, soit le cinquième du volume total. Dans le comprimé entier :
\[ n(\ce{Fe})_{\text{comprimé}} = 6{,}00\times 10^{-4} \times \frac{100{,}0}{20{,}0} = \SI{3,00e-3}{mol} \]
Masse de fer :
\[ m = n \times M = \SI{3,00e-3}{} \times 55{,}8 = \SI{0,167}{g} = \SI{167}{mg} \]
\textbf{Conclusion :} le comprimé contient \SI{167}{mg} de fer, très différent des \SI{65}{mg} annoncés (écart de $157\,\%$, largement hors tolérance de $5\,\%$). Le comprimé n'est \textbf{pas conforme} à l'étiquette : ce type de contrôle qualité est précisément le travail des laboratoires de vérification des normes.
\end{exoresolu}

\courssub{Doser une solution de diiode par les ions thiosulfate}

\begin{methode}[Dosage iodométrique : \ce{I2} par \ce{S2O3^2-}]
\begin{enumerate}
\item \textbf{Écrire l'équation de dosage} à partir des couples \ce{I2}/\ce{I-} et \ce{S4O6^2-}/\ce{S2O3^2-} :
\[ \ce{I2 + 2S2O3^2- -> 2I- + S4O6^2-} \]
(1 diiode réagit avec \textbf{2} thiosulfates : le facteur 2 est l'erreur classique.)
\item \textbf{Repérer l'équivalence} : la solution de \ce{I2} est brune (ou jaune-orangé diluée) ; à l'équivalence elle devient \textbf{incolore} (décoloration). Pour plus de précision, ajouter quelques gouttes d'\textbf{empois d'amidon} (ou thiodène) en fin de dosage : la solution est bleu foncé et devient incolore à l'équivalence.
\item \textbf{Appliquer la relation à l'équivalence} :
\[ n(\ce{I2}) = \frac{1}{2}\,n(\ce{S2O3^2-})_E \quad\Longrightarrow\quad C(\ce{I2}) = \frac{C(\ce{S2O3^2-})\,V_E}{2\,V(\ce{I2})} \]
\item \textbf{Cas du dosage en retour (indirect)} : si le diiode a été \emph{produit} par une réaction préalable (ex. oxydation de \ce{I-} par l'espèce à doser), écris d'abord la relation stœchiométrique de cette réaction préalable, puis remonte à l'espèce dosée.
\end{enumerate}
\end{methode}

\begin{exoresolu}[Dosage direct d'une solution de Lugol]
Le Lugol est un antiseptique à base de diiode utilisé dans les dispensaires. On prélève $V_1 = \SI{10,0}{mL}$ d'une solution de Lugol diluée 10 fois et on la dose par une solution de thiosulfate de sodium de concentration $C_2 = \SI{5,00e-2}{mol.L^{-1}}$. En présence d'empois d'amidon, la décoloration survient pour $V_E = \SI{8,40}{mL}$.\\
\textbf{1.} Écrire l'équation de la réaction de dosage.\\
\textbf{2.} Déterminer la concentration en diiode de la solution diluée, puis celle du Lugol commercial.
\tcblower
\textbf{1. Équation de la réaction de dosage}

Couples : \ce{I2}/\ce{I-} et \ce{S4O6^2-}/\ce{S2O3^2-}.
\begin{align*}
\ce{I2 + 2e-} &\ce{-> 2I-}\\
\ce{2S2O3^2-} &\ce{-> S4O6^2- + 2e-}
\end{align*}
\[ \boxed{\ce{I2 + 2S2O3^2- -> 2I- + S4O6^2-}} \]

\textbf{2. Concentrations}

À l'équivalence, les réactifs sont dans les proportions stœchiométriques :
\[ n(\ce{I2}) = \frac{n(\ce{S2O3^2-})_E}{2} \quad\Longrightarrow\quad C_1 V_1 = \frac{C_2 V_E}{2} \]
\[ C_1 = \frac{C_2 V_E}{2 V_1} = \frac{\SI{5,00e-2}{} \times 8{,}40}{2 \times 10{,}0} = \SI{2,10e-2}{mol.L^{-1}} \]
La solution commerciale est 10 fois plus concentrée :
\[ C_{\text{Lugol}} = 10 \times C_1 = \SI{2,10e-1}{mol.L^{-1}} \]
\textbf{Conclusion :} la solution diluée est à $\SI{2,10e-2}{mol.L^{-1}}$ en diiode et le Lugol commercial à $\SI{0,210}{mol.L^{-1}}$. Ne jamais oublier, en conclusion, de \textbf{tenir compte du facteur de dilution}.
\end{exoresolu}

\courssub{Exploiter un dosage en retour (dosage indirect)}

\begin{methode}[Méthode du dosage en retour]
Quand l'espèce à doser ne réagit pas directement avec le titrant, on passe par un intermédiaire :
\begin{enumerate}
\item \textbf{Réaction 1 (préalable)} : l'espèce $X$ à doser réagit avec un réactif en excès connu (souvent \ce{I-}, qui devient \ce{I2}, ou un excès de titrant).
\item \textbf{Réaction 2 (dosage)} : on dose le produit formé (ou l'excès restant) par la solution titrante.
\item \textbf{Enchaîner les relations stœchiométriques} : écris $n(\text{produit dosé})$ en fonction de $n_E(\text{titrant})$, puis $n(X)$ en fonction de $n(\text{produit})$ grâce à l'équation 1.
\item \textbf{Conclure} avec la concentration ou la masse demandée, en tenant compte des dilutions et des prélèvements.
\end{enumerate}
Réflexe : fais un schéma-bilan des deux réactions avant tout calcul, et vérifie chaque coefficient.
\end{methode}

\begin{exoresolu}[Degré chlorométrique d'une eau de Javel]
L'eau de Javel doit son pouvoir désinfectant à l'ion hypochlorite \ce{ClO-}. On dilue 20 fois une eau de Javel commerciale (solution $S$). À $V_1 = \SI{10,0}{mL}$ de $S$, on ajoute un excès d'iodure de potassium en milieu acide : le diiode formé est ensuite dosé par une solution de thiosulfate à $C_2 = \SI{0,100}{mol.L^{-1}}$. L'équivalence est atteinte pour $V_E = \SI{12,5}{mL}$.\\
Réaction préalable : \ce{ClO- + 2I- + 2H+ -> Cl- + I2 + H2O}.\\
\textbf{1.} Écrire l'équation du dosage du diiode.\\
\textbf{2.} Calculer la concentration en \ce{ClO-} de la solution $S$, puis de l'eau de Javel commerciale.
\tcblower
\textbf{1. Équation du dosage du diiode}
\[ \boxed{\ce{I2 + 2S2O3^2- -> 2I- + S4O6^2-}} \]

\textbf{2. Concentration en ions hypochlorite}

\emph{Étape A — dosage du diiode.} À l'équivalence :
\[ n(\ce{I2}) = \frac{C_2 V_E}{2} = \frac{\SI{0,100}{} \times \SI{12,5e-3}{}}{2} = \SI{6,25e-4}{mol} \]

\emph{Étape B — remontée à \ce{ClO-}.} D'après la réaction préalable, 1 mole de \ce{ClO-} produit 1 mole de \ce{I2} :
\[ n(\ce{ClO-}) = n(\ce{I2}) = \SI{6,25e-4}{mol} \]
Cette quantité était contenue dans $V_1 = \SI{10,0}{mL}$ de $S$ :
\[ C_S(\ce{ClO-}) = \frac{n(\ce{ClO-})}{V_1} = \frac{\SI{6,25e-4}{}}{\SI{10,0e-3}{}} = \SI{6,25e-2}{mol.L^{-1}} \]

\emph{Étape C — solution commerciale} (diluée 20 fois) :
\[ C_{\text{commerciale}} = 20 \times C_S = \SI{1,25}{mol.L^{-1}} \]
\textbf{Conclusion :} l'eau de Javel commerciale contient $\SI{1,25}{mol.L^{-1}}$ d'ions hypochlorite. Ce dosage en retour est la méthode officielle de contrôle qualité des eaux de Javel vendues sur les marchés.
\end{exoresolu}

\courssub{Vérifier une indication portée sur une étiquette}

\begin{methode}[Confronter un résultat de dosage à une norme ou une étiquette]
\begin{enumerate}
\item \textbf{Calculer la grandeur expérimentale} (concentration molaire, concentration massique $C_m = C \times M$, pourcentage massique, masse par unité) dans \textbf{l'unité de l'étiquette}.
\item \textbf{Tenir compte de toutes les dilutions et prélèvements} : remonte toujours à l'échantillon commercial initial.
\item \textbf{Calculer l'écart relatif} :
\[ \text{écart} = \frac{|x_{\text{exp}} - x_{\text{étiquette}}|}{x_{\text{étiquette}}} \times 100 \ (\%) \]
\item \textbf{Comparer à la tolérance} (souvent 5 \%) et \textbf{conclure par une phrase} : conforme ou non conforme.
\end{enumerate}
\end{methode}

\begin{exoresolu}[Teneur en peroxyde d'hydrogène d'une eau oxygénée]
Une eau oxygénée pharmaceutique est étiquetée « 10 volumes », ce qui correspond à une concentration en \ce{H2O2} de $C_{\text{étiquette}} = \SI{0,89}{mol.L^{-1}}$. On dilue 10 fois cette eau oxygénée (solution $S$), on prélève $V_1 = \SI{10,0}{mL}$ de $S$ que l'on acidifie, puis on dose par du permanganate à $C_2 = \SI{2,00e-2}{mol.L^{-1}}$. Équivalence : $V_E = \SI{11,1}{mL}$.\\
Couples : \ce{MnO4-}/\ce{Mn^2+} et \ce{O2}/\ce{H2O2}.\\
\textbf{1.} Écrire l'équation de dosage.\\
\textbf{2.} Calculer la concentration de l'eau oxygénée commerciale et conclure sur la conformité (tolérance 5 \%).
\tcblower
\textbf{1. Équation de la réaction de dosage}
\begin{align*}
\ce{MnO4- + 8H+ + 5e-} &\ce{-> Mn^2+ + 4H2O} \quad (\times 2)\\
\ce{H2O2} &\ce{-> O2 + 2H+ + 2e-} \quad (\times 5)
\end{align*}
\[ \boxed{\ce{2MnO4- + 5H2O2 + 6H+ -> 2Mn^2+ + 5O2 + 8H2O}} \]

\textbf{2. Concentration et conformité}

À l'équivalence :
\[ \frac{n(\ce{MnO4-})}{2} = \frac{n(\ce{H2O2})}{5} \quad\Longrightarrow\quad n(\ce{H2O2}) = \frac{5}{2}\,C_2 V_E \]
\[ C_S(\ce{H2O2}) = \frac{5\,C_2 V_E}{2\,V_1} = \frac{5 \times \SI{2,00e-2}{} \times 11{,}1}{2 \times 10{,}0} = \SI{5,55e-2}{mol.L^{-1}} \]
Solution commerciale (diluée 10 fois) :
\[ C_{\text{exp}} = 10 \times C_S = \SI{0,555}{mol.L^{-1}} \]
Écart relatif avec l'étiquette :
\[ \text{écart} = \frac{|0{,}555 - 0{,}89|}{0{,}89} \times 100 \approx 38\,\% \]
\textbf{Conclusion :} l'écart ($38\,\%$) dépasse largement la tolérance de 5 \% : le flacon est \textbf{non conforme}. Cela s'explique souvent : l'eau oxygénée se décompose lentement (\ce{2H2O2 -> 2H2O + O2}) surtout à la chaleur et à la lumière ; un flacon ancien ou mal conservé perd son titre, d'où l'intérêt des contrôles avant usage médical.
\end{exoresolu}

\begin{attention}[Les 5 erreurs qui coûtent des points au Bac]
\begin{enumerate}
\item \textbf{Oublier les coefficients stœchiométriques} : écrire $n(\ce{Fe^2+}) = n(\ce{MnO4-})$ au lieu de $n(\ce{Fe^2+}) = 5\,n(\ce{MnO4-})$, ou $n(\ce{I2}) = n(\ce{S2O3^2-})$ au lieu de la moitié. Écris toujours l'équation \emph{avant} la relation à l'équivalence.
\item \textbf{Inverser titrant et titré} dans le montage ou dans la formule : la solution de concentration \emph{connue} est dans la burette, l'inconnue est prélevée à la pipette jaugée.
\item \textbf{Mal décrire l'équivalence} : pour \ce{MnO4-}, c'est l'\emph{apparition} d'une teinte violette persistante ; pour \ce{I2}, c'est la \emph{décoloration} (disparition du bleu de l'empois d'amidon). Ne pas confondre les deux sens.
\item \textbf{Oublier la dilution ou le prélèvement} : le résultat du dosage concerne l'échantillon \emph{dosé} ; il faut remonter à la solution commerciale (facteur de dilution) ou au comprimé entier (rapport des volumes).
\item \textbf{Acidifier avec le mauvais acide ou pas du tout} : le dosage au permanganate exige un milieu acide (acide \emph{sulfurique}, jamais chlorhydrique dont \ce{Cl-} serait oxydé). Sans \ce{H+}, la demi-équation de \ce{MnO4-} est impossible.
\end{enumerate}
Bonus : soigne les unités (volumes convertis en L dans $n = CV$) et donne le résultat avec 3 chiffres significatifs et une phrase de conclusion.
\end{attention}

\newpage

\coursec{Exercices}

\coursec{Leçon 11 : Dosages d'oxydoréduction}

\courssub{Je vérifie mes connaissances}

\begin{aretenir}[Rappels essentiels : définitions et vocabulaire]
\textbf{Dosage} : Opération qui permet de déterminer la concentration d'une espèce chimique (solution \cle{titrée}) en la faisant réagir avec une solution de concentration connue (solution \cle{titrante}).
\textbf{Point d'équivalence} : Moment où les réactifs ont été introduits dans les proportions \cle{stœchiométriques} de l'équation de la réaction de dosage. On note $V_E$ le volume équivalent de solution titrante versé.
\textbf{Solution titrante} : Solution de concentration connue, placée dans la burette.
\textbf{Solution titrée} : Solution de concentration inconnue, prélevée à la pipette jaugée et placée dans l'erlenmeyer.
\end{aretenir}

\begin{aretenir}[Conditions d'un bon dosage d'oxydoréduction]
La réaction doit être :
\begin{enumerate}
\item \textbf{Rapide} : pas d'attente entre l'ajout et la réaction.
\item \textbf{Totale} : avance maximale (réaction « quasi-complète »), $K$ très grand.
\item \textbf{Unique} : pas de réaction parasite.
\item \textbf{Repérable} : changement brutal d'une grandeur physique (couleur, potentiel $E$, conductivité) à l'équivalence.
\end{enumerate}
\end{aretenir}

\begin{attention}[Pièges classiques]
\begin{itemize}
\item Ne pas confondre \emph{solution titrante} (burette, concentration connue) et \emph{solution titrée} (erlenmeyer, concentration cherchée).
\item L'équivalence n'est pas la « fin de réaction » (les réactifs ont disparu), c'est le moment où ils sont en proportions stœchiométriques exactes. S'il y a un excès de titrant, l'un des réactifs est en excès \emph{après} l'équivalence.
\item Lire le volume au \textbf{bas du ménisque} (pour les solutions aqueuses transparentes), l'œil au niveau du liquide.
\end{itemize}
\end{attention}

\begin{exercice}[QCM : les bons réflexes]
Pour chaque question, choisis la (ou les) bonne(s) réponse(s).
\begin{enumerate}
\item Une réaction de dosage doit être :
\begin{enumerate}
\item lente et totale ;
\item \textbf{rapide, totale et unique} ;
\item rapide et limitée ;
\item lente, totale et unique.
\end{enumerate}
\item Dans le dosage des ions \ce{Fe^2+} par les ions \ce{MnO4-}, la solution titrante est :
\begin{enumerate}
\item la solution de \ce{Fe^2+} ;
\item \textbf{la solution de \ce{MnO4-}} ;
\item l'acide sulfurique ;
\item l'eau distillée.
\end{enumerate}
\item À l'équivalence du dosage de \ce{Fe^2+} par \ce{MnO4-} :
\begin{enumerate}
\item $n(\ce{Fe^2+}) = n(\ce{MnO4-})$ ;
\item \textbf{$n(\ce{Fe^2+}) = 5\,n(\ce{MnO4-})$} ;
\item $n(\ce{MnO4-}) = 5\,n(\ce{Fe^2+})$ ;
\item les deux réactifs ont disparu.
\end{enumerate}
\item Dans le dosage du diiode par les ions thiosulfate, l'équivalence est repérée par :
\begin{enumerate}
\item l'apparition d'une coloration violette persistante ;
\item \textbf{la décoloration de la solution} ;
\item l'apparition d'un précipité ;
\item un dégagement gazeux.
\end{enumerate}
\end{enumerate}
\end{exercice}

\begin{exercice}[Vrai ou faux ?]
Indique si chaque affirmation est vraie ou fausse, puis justifie ta réponse en une ou deux phrases.
\begin{enumerate}
\item On peut acidifier le milieu d'un dosage par le permanganate avec l'acide chlorhydrique.
\item Le permanganate de potassium joue à la fois le rôle de réactif titrant et d'indicateur de fin de réaction.
\item À l'équivalence, les réactifs ont été mélangés dans les proportions stœchiométriques de l'équation de la réaction de dosage.
\item Une solution de thiosulfate de sodium peut être conservée indéfiniment sans précaution.
\item Le volume équivalent se lit toujours au sommet de la burette.
\end{enumerate}
\end{exercice}

\begin{exercice}[Texte à trous]
Complète le texte suivant avec les mots ou expressions de la liste : \emph{équivalence ; titrante ; burette graduée ; titrée ; stœchiométriques ; décoloration ; violette ; thiosulfate}.

\medskip
Dans un dosage d'oxydoréduction, la solution de concentration connue est appelée solution \textbf{titrante} ; elle est placée dans la \textbf{burette graduée}. La solution de concentration inconnue est dite solution \textbf{titrée}. L'\textbf{équivalence} est atteinte lorsque les réactifs ont été mélangés dans les proportions \textbf{stœchiométriques}. Lors du dosage des ions \ce{Fe^2+} par les ions \ce{MnO4-}, l'équivalence est repérée par l'apparition d'une teinte \textbf{violette} persistante. Lors du dosage du diiode, l'équivalence est repérée par la \textbf{décoloration} de la solution, et le réactif titrant est la solution de \textbf{thiosulfate} de sodium.
\end{exercice}

\begin{exercice}[Définitions et couples en jeu]
\begin{enumerate}
\item \textbf{Dosage} : Opération quantitative consistant à déterminer la concentration d'une espèce (titrée) par réaction complète avec une espèce de concentration connue (titrante). \textbf{Point d'équivalence} : État du système où les quantités de matière des réactifs introduits sont dans les rapports stœchiométriques de l'équation bilan.
\item Couples mis en jeu : \ce{MnO4-}/\ce{Mn^2+} (oxydant/réducteur) et \ce{Fe^3+}/\ce{Fe^2+} (oxydant/réducteur).
\item Demi-équations :
  \begin{align*}
  \text{Oxydation : } & \ce{2 S2O3^2- -> S4O6^2- + 2 e-} \\
  \text{Réduction : } & \ce{I2 + 2 e- -> 2 I-} \\
  \text{Bilan : } & \ce{I2 + 2 S2O3^2- -> 2 I- + S4O6^2-}
  \end{align*}
\end{enumerate}
\end{exercice}

\courssub{Je m'entraîne}

\begin{experience}[Dispositif expérimental type : dosage par titrage]
\begin{popfigure}
\begin{tikzpicture}[scale=0.75, every node/.style={font=\small}]
% Support
\draw[thick, popdark] (0,0) -- (0,5.5) -- (3.5,5.5) -- (3.5,5.0);
% Burette
\draw[thick, popblue] (2.5, 5.0) -- (2.5, 1.2) -- (3.0, 1.2) -- (3.0, 5.0);
\draw[thick, popblue] (2.5, 1.2) -- (2.75, 0.5) -- (2.75, 0.0);
% Robinet
\draw[fill=popgray, thick] (2.5, 1.2) rectangle (3.0, 0.9);
% Graduations burette
\foreach \y in {1.5, 2.0, ..., 4.5} \draw[popblue] (2.4, \y) -- (3.1, \y);
\node[left, popblue] at (2.4, 4.8) {0};
\node[left, popblue] at (2.4, 1.5) {50};
% Erlenmeyer
\draw[thick, popdark] (5.5,0) -- (6.0, 0.8) -- (9.0, 0.8) -- (9.5, 0) -- (9.5, -0.3) -- (9.2, -0.8) -- (6.3, -0.8) -- (6.0, -0.3) -- (6.0, 0) -- cycle;
\draw[thick, popdark] (7.25, 0.8) -- (7.25, 1.8) -- (8.25, 1.8) -- (8.25, 0.8);
% Liquide erlenmeyer
\fill[popblue!15] (6.1, -0.7) -- (9.1, -0.7) -- (9.1, 0.5) -- (6.1, 0.5) -- cycle;
% Agitateur magnétique
\draw[poporange, thick] (7.25, -0.7) -- (7.25, -1.3);
\draw[poporange, thick] (7.0, -1.3) -- (7.5, -1.3);
% Plaque chauffante/agitatrice
\fill[popgray!50] (5.0, -1.5) rectangle (10.0, -1.3);
% Étiquettes
\node[left, align=right] at (2.4, 3.0) {Burette graduée\\(solution titrante)};
\node[right] at (3.1, 0.6) {Robinet};
\node[right, align=left] at (9.6, 0.2) {Erlenmeyer\\(solution titrée + H$_2$SO$_4$)};
\node[below] at (7.25, -1.6) {Agitateur magnétique};
\end{tikzpicture}
\legende{Dispositif de dosage : burette (titrant), erlenmeyer (titré + indicateur/milieu), agitateur.}
\end{popfigure}
\end{experience}

\begin{exercice}[Schéma du dispositif]
Un technicien du laboratoire de la SONARA doit doser une solution de sulfate de fer(II) par une solution de permanganate de potassium.
\begin{enumerate}
\item Schématise et légende le dispositif expérimental du dosage (burette, erlenmeyer, agitateur magnétique, solutions). \emph{(Voir figure ci-dessus)}
\item Précise la verrerie utilisée pour prélever \SI{20,0}{mL} de la solution de \ce{Fe^2+} et justifie ce choix.
\item Explique pourquoi on ajoute de l'acide sulfurique dans l'erlenmeyer et pourquoi l'acide chlorhydrique est proscrit.
\end{enumerate}
\end{exercice}

\begin{exercice}[Dosage direct des ions fer(II)]
On dose $V_1 = \SI{20,0}{mL}$ d'une solution de sulfate de fer(II) acidifiée par une solution de permanganate de potassium de concentration $C_2 = \SI{0,020}{mol.L^{-1}}$. L'équivalence est obtenue pour $V_E = \SI{12,5}{mL}$.
\begin{enumerate}
\item Écris l'équation de la réaction de dosage.
\item Établis la relation entre les quantités de matière à l'équivalence.
\item Calcule la concentration $C_1$ de la solution de \ce{Fe^2+}.
\item Calcule la concentration massique en fer de cette solution.
\end{enumerate}
\emph{Donnée :} $M(\ce{Fe}) = \SI{55,8}{g.mol^{-1}}$.
\end{exercice}

\begin{exercice}[Dosage d'une solution de diiode]
On dose $V_1 = \SI{10,0}{mL}$ d'une solution de diiode par une solution de thiosulfate de sodium de concentration $C_2 = \SI{0,050}{mol.L^{-1}}$. Le volume versé à l'équivalence est $V_E = \SI{8,4}{mL}$.
\begin{enumerate}
\item Écris l'équation de la réaction de dosage sachant que les couples en jeu sont \ce{I2}/\ce{I-} et \ce{S4O6^2-}/\ce{S2O3^2-}.
\item Calcule la quantité de matière de diiode dosée, puis la concentration $C_1$ de la solution de diiode.
\item Décris précisément le changement de couleur observé à l'équivalence.
\end{enumerate}
\end{exercice}

\begin{exercice}[Exploitation d'une courbe potentiométrique]
On dose $V_1 = \SI{10,0}{mL}$ d'une solution de \ce{Fe^2+} par une solution de \ce{MnO4-} de concentration $C_2 = \SI{0,010}{mol.L^{-1}}$, en suivant le potentiel $E$ d'une électrode de platine plongée dans le mélange. Les mesures sont consignées ci-dessous.

\begin{center}
\begin{tabularx}{\linewidth}{|l|*{9}{>{\centering\arraybackslash}X|}}
\hline
\rowcolor{popblueL}
$V$ (mL) & 0 & 2 & 4 & 6 & 8 & 9 & 10 & 11 & 12 \\
\hline
$E$ (mV) & 780 & 800 & 830 & 870 & 950 & 1080 & 1420 & 1460 & 1470 \\
\hline
\end{tabularx}
\end{center}

\begin{enumerate}
\item Trace la courbe $E = f(V)$ sur papier millimétré (échelles : \SI{1}{cm} pour \SI{1}{mL} ; \SI{1}{cm} pour \SI{100}{mV}).
\item Détermine graphiquement le volume équivalent $V_E$ par la méthode des tangentes.
\item Calcule la concentration $C_1$ de la solution de \ce{Fe^2+}.
\end{enumerate}
\end{exercice}

\begin{methode}[Exploitation complète d'un dosage d'oxydoréduction]
\textbf{Étape 1 :} Écrire l'équation bilan de la réaction de dosage (demi-équations + bilan).
\textbf{Étape 2 :} Noter les données : $V_1$ (volume titré), $C_2$ (concentration titrant), $V_E$ (volume équivalent).
\textbf{Étape 3 :} Écrire la relation à l'équivalence : $n_{\text{titré}} = \frac{\nu_{\text{titré}}}{\nu_{\text{titrant}}} \times n_{\text{titrant}}$.
\textbf{Étape 4 :} Calculer $n_{\text{titrant}} = C_2 \times V_E$ (en L), puis $n_{\text{titré}}$.
\textbf{Étape 5 :} Calculer $C_1 = \frac{n_{\text{titré}}}{V_1}$ (en L).
\textbf{Étape 6 :} Si demandé, remonter à la masse, au degré (iodométrique, chlorométrique) ou au titre de la solution mère (facteur de dilution).
\textbf{Étape 7 :} Conclure en comparant à une norme (écart relatif $\%$).
\end{methode}

\courssub{Je me prépare au Bac}

\begin{savaistu}[Contexte camerounais : Santé publique et industrie]
\begin{itemize}
\item \textbf{Anémie et fer :} Au Cameroun, l'anémie touche près de 40\% des femmes enceintes (ENSP 2018). Le contrôle de la teneur en fer des médicaments (comprimés de sulfate ferreux) est vital. La pharmacopée impose une teneur en principe actif entre 95\% et 105\% de la dose annoncée.
\item \textbf{Teinture d'iode :} Antiseptique majeur dans les formations sanitaires (CSP, hôpitaux). Son degré iodométrique garantit l'efficacité désinfectante.
\item \textbf{Eau de Javel (SONARA/usines locales) :} Produite par électrolyse de saumure. Le degré chlorométrique (°chl) définit sa puissance désinfectante. L'OMS recommande 0,5\% de chlore actif (env. 15 °chl) pour la désinfection de l'eau de boisson.
\end{itemize}
\end{savaistu}

\begin{exercice}[Type Bac : comprimé de fer contre l'anémie]
Pour lutter contre l'anémie, fréquente chez les femmes enceintes suivies à l'hôpital de district de Bafoussam, on prescrit des comprimés de fer dont l'étiquette indique : « fer : \SI{60}{mg} par comprimé ». Un élève de Terminale TI veut vérifier cette indication.

Il dissout un comprimé dans une fiole jaugée de \SI{100}{mL} contenant de l'acide sulfurique dilué, puis complète au trait de jauge avec de l'eau distillée : il obtient la solution $S$. Il prélève $V_1 = \SI{20,0}{mL}$ de $S$ qu'il dose par une solution de permanganate de potassium de concentration $C_2 = \SI{0,010}{mol.L^{-1}}$. L'équivalence est obtenue pour $V_E = \SI{15,8}{mL}$.

\emph{Données :} $M(\ce{Fe}) = \SI{55,8}{g.mol^{-1}}$ ; couples : \ce{MnO4-}/\ce{Mn^2+} et \ce{Fe^3+}/\ce{Fe^2+}.

\begin{enumerate}
\item Écris les demi-équations électroniques puis l'équation bilan de la réaction de dosage.
\item Schématise et légende le dispositif de dosage.
\item Comment repère-t-on expérimentalement l'équivalence ? Justifie.
\item Calcule la quantité de matière d'ions \ce{MnO4-} versés à l'équivalence, puis la quantité d'ions \ce{Fe^2+} présents dans la prise d'essai.
\item En déduis la concentration molaire $C_1$ de la solution $S$, puis la masse de fer contenue dans un comprimé.
\item Compare le résultat à l'indication de l'étiquette en calculant l'écart relatif en \%. Le fabricant respecte-t-il la norme (tolérance de 5 \%) ?
\end{enumerate}
\end{exercice}

\begin{exercice}[Type Bac : degré iodométrique d'une teinture d'iode]
La teinture d'iode, antiseptique courant dans les pharmacies de Yaoundé, est une solution alcoolique de diiode. Son \cle{degré iodométrique} est la masse de diiode (en grammes) contenue dans \SI{100}{mL} de solution. L'étiquette d'un flacon annonce un degré de 5.

Un laborantin dilue d'abord la teinture 10 fois : il obtient une solution $S'$. Il prélève $V_1 = \SI{20,0}{mL}$ de $S'$ qu'il dose par une solution de thiosulfate de sodium de concentration $C_2 = \SI{0,100}{mol.L^{-1}}$. L'équivalence est obtenue pour $V_E = \SI{15,6}{mL}$.

\emph{Données :} $M(\ce{I}) = \SI{126,9}{g.mol^{-1}}$ ; couples : \ce{I2}/\ce{I-} et \ce{S4O6^2-}/\ce{S2O3^2-}.

\begin{enumerate}
\item Écris l'équation de la réaction de dosage.
\item Décris le changement de couleur observé à l'équivalence. Pourquoi peut-on ajouter de l'empois d'amidon pour un repérage plus précis ?
\item Calcule la quantité de matière d'ions thiosulfate versés à l'équivalence, puis la quantité de diiode dosée.
\item En déduis la concentration molaire du diiode dans $S'$, puis dans la teinture commerciale.
\item Calcule le degré iodométrique de la teinture. L'indication de l'étiquette est-elle correcte (tolérance de 5 \%) ?
\end{enumerate}
\end{exercice}

\begin{exercice}[Type Bac : degré chlorométrique d'une eau de Javel]
L'eau de Javel, très utilisée pour la désinfection dans les ménages camerounais, contient des ions hypochlorite \ce{ClO-}. Son \cle{degré chlorométrique} (°chl) est le volume de dichlore gazeux (en litres, mesuré dans les CNTP) que peut libérer un litre d'eau de Javel selon :
\[ \ce{ClO- + 2H+ + Cl- -> Cl2 + H2O} \]

On procède à un \cle{dosage en retour} (ou dosage indirect) :
\begin{itemize}
\item On dilue l'eau de Javel commerciale 20 fois : solution $S$.
\item Dans un erlenmeyer, on introduit $V_0 = \SI{10,0}{mL}$ de $S$, un excès d'iodure de potassium et de l'acide éthanoïque. Les ions \ce{ClO-} oxydent les ions \ce{I-} : $\ce{ClO- + 2I- + 2H+ -> Cl- + I2 + H2O}$.
\item Le diiode formé est dosé par une solution de thiosulfate de sodium de concentration $C_2 = \SI{0,100}{mol.L^{-1}}$. L'équivalence est obtenue pour $V_E = \SI{11,2}{mL}$.
\end{itemize}

\emph{Données :} volume molaire dans les CNTP : $V_m = \SI{22,4}{L.mol^{-1}}$ ; couples : \ce{ClO-}/\ce{Cl-}, \ce{I2}/\ce{I-}, \ce{S4O6^2-}/\ce{S2O3^2-}.

\begin{enumerate}
\item Justifie le qualificatif de « dosage en retour » (ou dosage indirect) donné à cette méthode.
\item Écris l'équation de la réaction de dosage du diiode par les ions thiosulfate.
\item Calcule la quantité de matière de diiode formée dans l'erlenmeyer.
\item En déduis la quantité d'ions \ce{ClO-} présents dans la prise de \SI{10,0}{mL}, puis la concentration en ions \ce{ClO-} de la solution $S$ et de l'eau de Javel commerciale.
\item Montre que la quantité de dichlore libérable par un litre d'eau de Javel est égale à la quantité d'ions \ce{ClO-} qu'il contient. Calcule alors le degré chlorométrique de cette eau de Javel.
\item L'étiquette du flacon indique « 12 °chl ». Calcule l'écart relatif et conclus.
\end{enumerate}
\end{exercice}

\begin{attention}[Rappels sécurité et rigueur pour le Bac]
\begin{itemize}
\item \textbf{Permanganate + HCl :} Interdit ! $\ce{2 MnO4- + 10 Cl- + 16 H+ -> 2 Mn^2+ + 5 Cl2 + 8 H2O}$ (fausse réaction, consommation de titrant, gaz toxique).
\item \textbf{Thiosulfate :} Solution instable (décomposition acide/light/bactéries). Toujours étalonner avant usage. Ne pas chauffer.
\item \textbf{Amidon :} Ajouté \textbf{uniquement} près de l'équivalence (coloration jaune paille) pour éviter la formation d'un complexe amidon-iode trop stable qui fausserait $V_E$.
\item \textbf{Dosage en retour :} L'espèce dosée (\ce{ClO-}) ne réagit pas directement avec le titrant (\ce{S2O3^2-}). Elle génère une espèce intermédiaire (\ce{I2}) qui est dosée. Relation de chaînage : $n(\ce{ClO-}) = n(\ce{I2}) = \frac{1}{2}n(\ce{S2O3^2-})$.
\end{itemize}
\end{attention}

\newpage

\coursec{Corrigés des exercices}

\coursec{Généralités sur les dosages d'oxydoréduction}

\courssub{Définitions et vocabulaire fondamental}

\begin{definition}[Dosage]
Opération quantitative qui consiste à déterminer la concentration d'une espèce chimique (la \cle{solution titrée}) en la faisant réagir de façon \cle{totale}, \cle{rapide} et \cle{unique} avec une espèce de concentration connue (la \cle{solution titrante}).
\end{definition}

\begin{definition}[Solution titrante et solution titrée]
\begin{itemize}
    \item \textbf{Solution titrante} (notée $S_2$, concentration $C_2$) : solution de concentration \emph{connue avec précision}, placée dans la \cle{burette graduée}.
    \item \textbf{Solution titrée} (notée $S_1$, concentration $C_1$) : solution de concentration \emph{inconnue}, prélevée à l'aide d'une \cle{pipette jaugée} et versée dans l'\cle{erlenmeyer}.
\end{itemize}
\end{definition}

\begin{definition}[Point d'équivalence]
État du système atteint lorsque les quantités de matière des réactifs introduits sont exactement dans les rapports \cle{stœchiométriques} définis par l'équation-bilan de la réaction de dosage. On note $V_E$ le volume de solution titrante versé à l'équivalence.
\end{definition}

\begin{aretenir}[Notations usuelles au Bac TI]
\begin{center}
\begin{tabularx}{\linewidth}{|l|X|l|}\hline
\textbf{Symbole} & \textbf{Sens} & \textbf{Unité} \\\hline
$C_1$ & Concentration molaire de la solution titrée & \SI{}{mol.L^{-1}} \\\hline
$V_1$ & Volume de la prise d'essai (pipette jaugée) & \SI{}{mL} ou \SI{}{L} \\\hline
$C_2$ & Concentration molaire de la solution titrante & \SI{}{mol.L^{-1}} \\\hline
$V_E$ & Volume équivalent (lu sur la burette) & \SI{}{mL} ou \SI{}{L} \\\hline
\end{tabularx}
\end{center}
\end{aretenir}

\courssub{Conditions d'un bon dosage d'oxydoréduction}

\begin{propriete}[Les trois conditions indispensables]
Pour qu'un dosage soit rigoureux, la réaction chimique mise en jeu doit être :
\begin{enumerate}
    \item \textbf{Totale} : la constante d'équilibre $K$ est très grande ($K > 10^4$). La réaction va « jusqu'au bout », ce qui garantit la relation stœchiométrique exacte à l'équivalence.
    \item \textbf{Rapide} : la vitesse de réaction est quasi instantanée à température ambiante. Le repérage de l'équivalence est immédiat après chaque ajout de titrant.
    \item \textbf{Unique} : aucune \cle{réaction parasite} ne consomme le titrant ou la substance dosée. La stœchiométrie de l'équation-bilan est la seule qui compte.
\end{enumerate}
\end{propriete}

\begin{attention}[Pièges fréquents]
\begin{itemize}
    \item Une réaction \emph{lente} (ex. oxydation des ions \ce{Mn^2+} par \ce{MnO4-} en milieu neutre) rend le repérage impossible : il faut attendre, la couleur change après coup.
    \item Une réaction \emph{limitée} (ex. dosage d'un acide faible par une base forte sans indicateur adapté) fausse le calcul : $n_{\text{réel}} \neq n_{\text{stœch}}$.
    \item Une \emph{réaction parasite} (ex. \ce{Cl-} oxydé par \ce{MnO4-} en milieu chlorhydrique) consomme du titrant « pour rien » : $V_E$ mesuré $>$ $V_E$ théorique $\Rightarrow$ $C_1$ calculée surestimée.
\end{itemize}
\end{attention}

\courssub{Dispositif expérimental et verrerie de précision}

\begin{experience}[Montage type d'un dosage]
\begin{itemize}
    \item \textbf{Support} : pied à pince + pince à burette.
    \item \textbf{Burette graduée} (classe A, \SI{25}{mL} ou \SI{50}{mL}) : contient la solution titrante. \textbf{Prélavage} : 3 fois avec quelques mL de la solution titrante (évite la dilution).
    \item \textbf{Erlenmeyer} (conique, \SI{100}{mL} ou \SI{250}{mL}) : contient la prise d'essai $V_1$ de solution titrée + réactifs annexes (acide, indicateur). Forme conique = pas de projection au mélange.
    \item \textbf{Agitateur magnétique} + barreau aimanté : mélange homogène et continu sans éclaboussure.
    \item \textbf{Pipette jaugée} (classe A, \SI{10,0}{mL}, \SI{20,0}{mL} ou \SI{25,0}{mL}) + \textbf{propipette} : prélève $V_1$ avec exactitude (incertitude $\approx \pm \SI{0,03}{mL}$ pour \SI{20,0}{mL}). \textbf{Jamais} d'éprouvette graduée (incertitude $\approx \pm \SI{0,2}{mL}$).
    \item \textbf{Fiole jaugée} (classe A, \SI{100}{mL}, \SI{250}{mL}...) : pour la préparation de la solution titrée par dissolution/dilution.
\end{itemize}
\end{experience}

\begin{popfigure}
\begin{tikzpicture}[scale=0.85, every node/.style={transform shape}]
% Support
\draw[popdark, thick] (0,0) -- (0,5.5) -- (3.5,5.5) -- (3.5,5) -- (0.5,5) -- (0.5,0) -- cycle;
% Burette
\draw[popblue, thick] (2.5, 5) -- (2.5, 1.2) -- (3.5, 1.2) -- (3.5, 5);
% Liquid in burette
\fill[popblue!20] (2.5, 1.2) rectangle (3.5, 4.2);
% Meniscus
\draw[popblue, thick] (2.5, 4.2) .. controls (2.8, 4.15) and (3.2, 4.15) .. (3.5, 4.2);
% Graduations
\foreach \y in {1.5, 2.0, ..., 4.5} {
  \draw[popdark] (2.4, \y) -- (2.5, \y);
  \draw[popdark] (3.5, \y) -- (3.6, \y);
  \ifdim\y pt=2.0pt\node[left, font=\tiny, popdark] at (2.3, \y) {20};\fi
  \ifdim\y pt=3.0pt\node[left, font=\tiny, popdark] at (2.3, \y) {30};\fi
  \ifdim\y pt=4.0pt\node[left, font=\tiny, popdark] at (2.3, \y) {40};\fi
}
% Tap (robinet)
\draw[popdark, thick] (2.5, 1.2) -- (3.0, 0.7) -- (3.5, 1.2);
\draw[popdark, thick] (3.0, 0.7) -- (3.0, 0.4);
% Drop falling
\draw[popblue, thick] (3.0, 0.4) -- (3.0, 0.2);
\shade[ball color=popblue] (3.0, 0.1) circle (0.08);
% Erlenmeyer
\draw[popgreen, thick] (1.2, 0) -- (1.6, 1.2) -- (4.4, 1.2) -- (4.8, 0);
% Liquid in erlenmeyer
\fill[popgreen!15] (1.7, 1.15) -- (4.3, 1.15) -- (4.0, 0.2) -- (2.0, 0.2) -- cycle;
% Stirrer bar
\draw[poporange, thick, rotate=15] (3.0, 0.3) ellipse (0.4 and 0.1);
% Magnetic stirrer base
\draw[popmuted, thick, rounded corners=2pt] (0.8, -0.6) rectangle (5.2, -0.1);
\node[popmuted, font=\small] at (3, -0.35) {Agitateur magnétique};
% Labels
\node[popblue, right, font=\small] at (3.7, 3.5) {Burette : solution titrante ($S_2$)};
\node[popgreen, right, font=\small] at (5.0, 0.5) {Erlenmeyer : solution titrée ($S_1$) + $\ce{H2SO4}$};
\node[poporange, above, font=\small] at (3.0, 0.8) {Barreau aimanté};
\end{tikzpicture}
\legende{Dispositif expérimental de dosage : burette graduée (titrant), erlenmeyer sur agitateur magnétique (titré acidifié). Lecture du ménisque au bas du creux, œil à hauteur.}
\end{popfigure}

\begin{methode}[Prélavage et mise en place]
\begin{enumerate}
    \item \textbf{Burette} : 3 prélavages avec $\approx \SI{5}{mL}$ de solution titrante $S_2$. Vider par le robinet pour chasser les bulles d'air dans la pointe.
    \item \textbf{Pipette jaugée} : 1 prélavage avec la solution titrée $S_1$ (ou l'eau distillée si étalonnage « à sec » mais préférence à $S_1$). Prélèvement de $V_1$ : aspiration au-dessus du trait de jauge, réglage au bas du ménisque, écoulement libre dans l'erlenmeyer (ne pas souffler la dernière goutte, la pipette est étalonnée « à l'écoulement »).
    \item \textbf{Erlenmeyer} : rincé à l'eau distillée uniquement (pas de prélavage avec $S_1$ car cela ajouterait de la matière). Verser $V_1$, ajouter l'acide sulfurique (voir §3), poser sur l'agitateur.
    \item \textbf{Remplissage burette} : au-dessus du zéro, chasse de la bulle d'air, lecture initiale $V_i$ au bas du ménisque, œil à hauteur.
\end{enumerate}
\end{methode}

\begin{savaistu}[Contexte camerounais : contrôle qualité]
Au Cameroun, les laboratoires du \textbf{LANORM} (Laboratoire National de la Normalisation) et ceux des industries agro-alimentaires (brasseries \textbf{UCB} pour le contrôle du degré d'alcool, \textbf{SOCAPALM} pour l'acidité de l'huile de palme, \textbf{SONARA} pour la teneur en soufre des carburants) réalisent quotidiennement des dosages redox. La rigueur du prélavage et du choix de la verrerie classe A est une exigence normative (normes ISO/CEI 17025).
\end{savaistu}

\coursec{Dosage des ions fer(II) par les ions permanganate}

\courssub{Couples redox en jeu et équation-bilan}

\begin{propriete}[Couples mis en jeu]
\begin{itemize}
    \item Couple oxydant/réducteur du titrant : \ce{MnO4-}/\ce{Mn^2+} ($E^\circ = \SI{+1,51}{V}$).
    \item Couple oxydant/réducteur du titré : \ce{Fe^3+}/\ce{Fe^2+} ($E^\circ = \SI{+0,77}{V}$).
\end{itemize}
Comme $E^\circ(\ce{MnO4-}/\ce{Mn^2+}) > E^\circ(\ce{Fe^3+}/\ce{Fe^2+})$, les ions \ce{MnO4-} oxydent \emph{spontanément} et \emph{totalement} les ions \ce{Fe^2+} en milieu acide.
\end{propriete}

\begin{methode}[Écriture de l'équation-bilan (méthode des demi-équations)]
\begin{enumerate}
    \item \textbf{Demi-équation de réduction} (oxydant $\ce{MnO4-}$) : bilan d'atomes O par $\ce{H2O}$, bilan d'atomes H par $\ce{H+}$, bilan de charges par $\ce{e-}$.
    \[ \ce{MnO4- + 8H+ + 5e- -> Mn^2+ + 4H2O} \]
    \item \textbf{Demi-équation d'oxydation} (réducteur \ce{Fe^2+}) :
    \[ \ce{Fe^2+ -> Fe^3+ + e-} \]
    \item \textbf{Égalisation des électrons} : multiplier la 2ème par 5.
    \item \textbf{Somme} : simplifier les $\ce{e-}$.
    \[ \ce{MnO4- + 5Fe^2+ + 8H+ -> Mn^2+ + 5Fe^3+ + 4H2O} \]
\end{enumerate}
\end{methode}

\begin{aretenir}[Relation stœchiométrique à l'équivalence]
D'après l'équation-bilan : \textbf{1 mol \ce{MnO4-} réagit avec 5 mol \ce{Fe^2+}}.
\[ n(\ce{Fe^2+}) = 5 \times n(\ce{MnO4-}) \quad \Longleftrightarrow \quad C_1 V_1 = 5\, C_2 V_E \]
\end{aretenir}

\courssub{Repérage de l'équivalence : l'auto-indication}

\begin{propriete}[Indicateur coloré intrinsèque]
L'ion permanganate \ce{MnO4-} est \textbf{violet intense} (couleur propre). L'ion manganèse(II) \ce{Mn^2+} formé est \textbf{quasi incolore} (très pâle rose, invisible en solution diluée).
\begin{itemize}
    \item \textbf{Avant l'équivalence} : tout \ce{MnO4-} versé est consommé instantanément $\rightarrow$ solution incolore (ou vert pâle dû à \ce{Fe^2+}/\ce{Fe^3+}).
    \item \textbf{À l'équivalence} : la première goutte de \ce{MnO4-} en excès n'est plus réduite $\rightarrow$ apparition d'une \textbf{teinte rose-violette persistante} (environ \SI{30}{\second} sans s'estomper).
\end{itemize}
\cle{Le permanganate est à la fois titrant et indicateur.} Aucun indicateur externe n'est nécessaire.
\end{propriete}

\courssub{Rôle de l'acide sulfurique et interdiction de l'acide chlorhydrique}

\begin{attention}[Milieu acide obligatoire : \ce{H2SO4} uniquement !]
\begin{itemize}
    \item \textbf{Rôle des $\ce{H+}$} : 8 $\ce{H+}$ sont nécessaires par $\ce{MnO4-}$ réduit (demi-équation). Sans acide, \ce{MnO4-} se réduit en \ce{MnO2} (dioxyde de manganèse), précipité \textbf{brun} insoluble : $\ce{MnO4- + 4H+ + 3e- -> MnO2 v + 2H2O}$. La stœchiométrie passe de 5 à 3 électrons $\rightarrow$ résultats faux.
    \item \textbf{Pourquoi pas \ce{HCl} ?} Les ions \ce{Cl-} sont oxydés par \ce{MnO4-} (réaction parasite) :
    \[ \ce{2MnO4- + 10Cl- + 16H+ -> 2Mn^2+ + 5Cl2 ^ + 8H2O} \]
    Conséquences : (1) Surconsommation de titrant $\rightarrow$ $V_E$ trop grand $\rightarrow$ $C_1$ surestimée. (2) Dégagement de \ce{Cl2}, gaz \textbf{toxique} et corrosif (danger pour l'opérateur).
    \item \textbf{Pourquoi pas \ce{HNO3} ?} L'ion nitrate \ce{NO3-} peut être réduit en \ce{NO} ou \ce{NO2} (gaz brun toxique) par \ce{Fe^2+} à chaud, ou perturber le potentiel.
    \item \textbf{Acide sulfurique \ce{H2SO4}} : l'ion \ce{SO4^2-} n'est pas oxydable par \ce{MnO4-} (soufre déjà au degré max +VI). Il fournit les $\ce{H+}$ sans réaction parasite.
\end{itemize}
\end{attention}

\begin{exemplebox}[Application numérique directe]
On dose \SI{20,0}{mL} d'une solution $S_1$ d'ions \ce{Fe^2+} par une solution $S_2$ de \ce{KMnO4} à $C_2 = \SI{0,0200}{mol.L^{-1}}$. Volume équivalent $V_E = \SI{12,5}{mL}$.
\begin{align*}
n(\ce{MnO4-}) &= C_2 V_E = 0,0200 \times 12,5 \times 10^{-3} = \SI{2,50e-4}{mol} \\
n(\ce{Fe^2+}) &= 5 \times n(\ce{MnO4-}) = \SI{1,25e-3}{mol} \\
C_1 &= \frac{n(\ce{Fe^2+})}{V_1} = \frac{1,25 \times 10^{-3}}{20,0 \times 10^{-3}} = \SI{0,0625}{mol.L^{-1}} \\
\text{Concentration massique en Fe} &= C_1 \times M(\ce{Fe}) = 0,0625 \times 55,8 = \SI{3,49}{g.L^{-1}}
\end{align*}
\end{exemplebox}

\coursec{Dosage du diiode par les ions thiosulfate}

\courssub{Couples redox et équation-bilan}

\begin{propriete}[Couples mis en jeu]
\begin{itemize}
    \item Couple \ce{I2}/\ce{I-} : $E^\circ = \SI{+0,54}{V}$ (diiode = oxydant).
    \item Couple \ce{S4O6^2-}/\ce{S2O3^2-} (tétrathionate/thiosulfate) : $E^\circ = \SI{+0,08}{V}$ (thiosulfate = réducteur).
\end{itemize}
La réaction est spontanée et totale ($E^\circ_{\text{cell}} = 0,46\,\text{V} > 0$).
\end{propriete}

\begin{methode}[Équation-bilan du dosage iodométrique]
\begin{enumerate}
    \item Réduction du diiode : $\ce{I2 + 2e- -> 2I-}$
    \item Oxydation du thiosulfate : $\ce{2S2O3^2- -> S4O6^2- + 2e-}$ (formation de l'ion tétrathionate).
    \item Bilan :
    \[ \ce{I2 + 2S2O3^2- -> 2I- + S4O6^2-} \]
\end{enumerate}
\end{methode}

\begin{aretenir}[Relation stœchiométrique à l'équivalence]
\textbf{1 mol \ce{I2} réagit avec 2 mol \ce{S2O3^2-}}.
\[ n(\ce{I2}) = \frac{1}{2}\, n(\ce{S2O3^2-}) \quad \Longleftrightarrow \quad C_1 V_1 = \frac{1}{2}\, C_2 V_E \]
\end{aretenir}

\courssub{Repérage de l'équivalence : décoloration et empois d'amidon}

\begin{propriete}[Observations visuelles]
\begin{itemize}
    \item \textbf{Sans indicateur} : La solution titrée contient du diiode $\rightarrow$ couleur \textbf{brun-jaune} (plus ou moins foncé selon la concentration). Au fur et à mesure de l'ajout de thiosulfate, la teinte s'atténue (jaune paille). À l'équivalence, tout le \ce{I2} est consommé $\rightarrow$ \textbf{décoloration complète} (solution incolore). \emph{Difficulté} : le passage du jaune très pâle à l'incolore est subjectif.
    \item \textbf{Avec empois d'amidon} (solution d'amidon $\approx 1\%$) : L'amidon forme avec \ce{I2} un \textbf{complexe d'inclusion bleu foncé} (très intense, visible à très faible concentration). On ajoute l'amidon \textbf{près de l'équivalence} (quand la solution devient jaune paille).
    \begin{itemize}
        \item Avant équivalence : excès \ce{I2} $\rightarrow$ \textbf{bleu foncé}.
        \item À l'équivalence : dernier \ce{I2} consommé $\rightarrow$ \textbf{passage brutal du bleu foncé à l'incolore}. Repérage très précis ($V_E$ au \SI{0,05}{mL} près).
    \end{itemize}
\end{itemize}
\end{propriete}

\begin{attention}[Pourquoi ne pas ajouter l'amidon dès le début ?]
En présence d'un \textbf{excès important de diiode}, le complexe amidon-\ce{I2} est très stable. Il « piège » le diiode et le libère tardivement lors du dosage, ce qui \textbf{retarde le virage} et fausse $V_E$ (volume mesuré trop grand). On attend que la concentration en \ce{I2} soit faible (jaune paille) pour ajouter l'amidon.
\end{attention}

\begin{attention}[Stabilité de la solution de thiosulfate]
La solution de \ce{Na2S2O3} n'est \textbf{pas une solution étalon primaire} (elle ne se conserve pas). Elle se décompose par :
\begin{itemize}
    \item \textbf{Photolyse} (lumière) : $\ce{2S2O3^2- -> S4O6^2- + 2e-}$ (auto-oxydation).
    \item \textbf{Acidité} : $\ce{S2O3^2- + 2H+ -> SO2 ^ + S v + H2O}$.
    \item \textbf{Bactéries} : métabolisme microbien.
\end{itemize}
\cle{Conservation} : flacon ambré (abri de la lumière), addition de quelques gouttes de \ce{NaOH} ou \ce{CH3COONa} (tampon pH $\approx 9$), renouvellement fréquent. \cle{Étalonnage obligatoire} avant utilisation (par dosage d'une solution étalon de \ce{K2Cr2O7} ou \ce{KIO3}).
\end{attention}

\courssub{Degré iodométrique : application industrielle}

\begin{definition}[Degré iodométrique]
Masse de diiode (\ce{I2}), en \textbf{grammes}, contenue dans \textbf{\SI{100}{mL}} de solution commerciale (teinture d'iode, eau de Javel après acidification/ajout de KI).
\[ \text{Degré iodométrique} = \frac{m(\ce{I2}) \text{ dans } \SI{100}{mL}}{\SI{1}{g}} \quad (\text{sans unité, noté } ^\circ\text{I}) \]
\end{definition}

\begin{exemplebox}[Calcul du degré iodométrique]
Solution commerciale diluée 10 fois ($S'$). Dosage de \SI{20,0}{mL} de $S'$ par \ce{Na2S2O3} \SI{0,100}{M} : $V_E = \SI{15,6}{mL}$.
\begin{align*}
n(\ce{S2O3^2-}) &= 0,100 \times 15,6 \times 10^{-3} = \SI{1,56e-3}{mol} \\
n(\ce{I2})_{\text{dans } 20\,\text{mL } S'} &= \frac{1,56 \times 10^{-3}}{2} = \SI{7,80e-4}{mol} \\
C'(S') &= \frac{7,80 \times 10^{-4}}{20,0 \times 10^{-3}} = \SI{3,90e-2}{mol.L^{-1}} \\
C(\text{commerciale}) &= 10 \times C' = \SI{0,390}{mol.L^{-1}} \\
m(\ce{I2}) \text{ par L} &= 0,390 \times 253,8 = \SI{99,0}{g.L^{-1}} \quad (M(\ce{I2}) = 2 \times 126,9) \\
\text{Degré iodométrique} &= \frac{99,0}{10} = \mathbf{9,9^\circ\text{I}} \quad (\text{soit } \SI{9,9}{g} \text{ pour } \SI{100}{mL})
\end{align*}
\end{exemplebox}

\coursec{Exploitation complète d'un dosage d'oxydoréduction}

\courssub{Méthode pas à pas (type Bac)}

\begin{methode}[Protocole complet de résolution]
\begin{enumerate}
    \item \textbf{Écrire l'équation-bilan} de la réaction de dosage (demi-équations + somme). Identifier les couples.
    \item \textbf{Schéma du dispositif} : burette (titrant), erlenmeyer (titré + réactifs), agitateur. Légendes obligatoires.
    \item \textbf{Repérage de l'équivalence} : description précise de l'observation (couleur persistante, décoloration, virage amidon, saut potentiométrique).
    \item \textbf{Calcul des quantités de matière} : $n_{\text{titrant}} = C_2 V_E$ (en L). Déduire $n_{\text{titré}}$ par la stœchiométrie.
    \item \textbf{Concentration de la solution titrée} : $C_1 = n_{\text{titré}} / V_1$ (en L).
    \item \textbf{Remontée à l'échantillon initial} : dilution, masse molaire, degré (iodométrique, chlorométrique), pureté, teneur.
    \item \textbf{Conclusion critique} : écart relatif par rapport à la norme/étiquette, décision (conforme / non conforme), enjeu sécurité/qualité.
\end{enumerate}
\end{methode}

\courssub{Tableau d'avancement : outil indispensable}

\begin{exemplebox}[Tableau d'avancement pour \ce{Fe^2+} / \ce{MnO4-}]
On dose $V_1 = \SI{20,0}{mL}$ de $S_1$ (\ce{Fe^2+}) par $S_2$ (\ce{MnO4-}, $C_2$). Avancement $x$ en mol.
\[ \ce{MnO4- + 5Fe^2+ + 8H+ -> Mn^2+ + 5Fe^3+ + 4H2O} \]
\begin{center}
\begin{tabular}{|c|c|c|c|c|c|}\hline
 & \ce{MnO4-} & \ce{Fe^2+} & \ce{H+} & \ce{Mn^2+} & \ce{Fe^3+} \\\hline
État initial & $C_2 V$ & $C_1 V_1$ & Excès & 0 & 0 \\\hline
État intermédiaire & $C_2 V - x$ & $C_1 V_1 - 5x$ & ... & $x$ & $5x$ \\\hline
\textbf{Équivalence} ($V=V_E$) & \textbf{0} & \textbf{0} & ... & $C_2 V_E$ & $5 C_2 V_E$ \\\hline
État final ($V>V_E$) & $C_2(V-V_E)$ & 0 & ... & ... & ... \\\hline
\end{tabular}
\end{center}
À l'équivalence : $C_1 V_1 - 5 C_2 V_E = 0 \Rightarrow C_1 V_1 = 5 C_2 V_E$.
\end{exemplebox}

\courssub{Repérage potentiométrique : méthode des tangentes}

\begin{experience}[Dosage potentiométrique]
\begin{itemize}
    \item \textbf{Dispositif} : même montage, mais ajout d'une \textbf{électrode de référence} (électrode calomel saturée ECS ou Ag/AgCl/KCl sat) et d'une \textbf{électrode indicatrice} (électrode de platine Pt) plongées dans l'erlenmeyer, reliées à un \textbf{voltmètre haute impédance} (pH-mètre en mode mV ou potentiomètre).
    \item \textbf{Mesure} : $E = f(V)$ après chaque ajout de titrant (ou en continu avec burette motorisée).
    \item \textbf{Courbe} : forme sigmoïde. Branche montante lente (zone tampon \ce{Fe^3+}/\ce{Fe^2+}), \textbf{saut brutal de potentiel} au voisinage de $V_E$, palier final (zone tampon \ce{MnO4-}/\ce{Mn^2+}).
\end{itemize}
\end{experience}

\begin{methode}[Méthode des tangentes (graphique)]
\begin{enumerate}
    \item Tracer la courbe $E = f(V)$.
    \item Tracer les \textbf{deux tangentes} aux branches quasi-linéaires (avant et après le saut).
    \item Tracer la \textbf{droite équidistante} (parallèle aux deux tangentes, passant au milieu verticalement).
    \item L'intersection de cette droite médiane avec la courbe donne le \textbf{point d'équivalence} ($V_E$, $E_E$).
    \item \textbf{Justification} : au point d'équivalence, la pente $\frac{dE}{dV}$ est maximale (point d'inflexion). Pour une courbe symétrique, l'inflexion est au milieu du saut.
\end{enumerate}
\end{methode}

\begin{popfigure}
\begin{tikzpicture}
\begin{axis}[
    width=\linewidth, height=9cm,
    xlabel={Volume de titrant $V$ (\SI{}{mL})}, ylabel={Potentiel $E$ (\SI{}{mV})},
    xmin=0, xmax=20, ymin=700, ymax=1550,
    grid=major, grid style={popline},
    axis lines=middle,
    enlargelimits=0.05,
    tick label style={font=\small},
    label style={font=\small},
    clip=false
]
% Experimental points / smooth curve for Fe2+/MnO4-
% Pre-equivalence: E = E_Fe + 0.059 log([Fe3+]/[Fe2+]) ~ 770 + 59 log(x/(1-x))
% Equivalence: E_eq = (E_Mn + 5E_Fe)/6 ~ (1510 + 5*770)/6 = 1150 mV approx
% Post-equivalence: E = E_Mn + 0.059/5 log([MnO4-]/[Mn2+])
% We plot a smooth sigmoid shape.
\addplot[popcurve, very thick, domain=0:20, samples=300] 
    ({x}, 
    {780 + 700/(1+exp(-12*(x-10)))} 
    );
% Tangents
\addplot[poporange, dashed, thick, domain=0:9, samples=2] {780 + 19*x}; % approx tangent before
\addplot[poporange, dashed, thick, domain=11:20, samples=2] {1470 + 2*(x-11)}; % approx tangent after
% Median line (equidistant)
\addplot[poppurple, dashed, thick, domain=0:20, samples=2] {1125 + 10.5*(x-10)}; % rough middle slope
% Vertical line at VE
\draw[poppurple, very thick, dashed] (axis cs:10,700) -- (axis cs:10,1500);
\node[poppurple, anchor=south, font=\small\bfseries] at (axis cs:10, 1510) {$V_E = \SI{10,0}{mL}$};
% Annotations
\node[popblue, anchor=south east, font=\small] at (axis cs:5, 900) {Branche \ce{Fe^3+}/\ce{Fe^2+}};
\node[popgreen, anchor=north west, font=\small] at (axis cs:15, 1450) {Branche \ce{MnO4-}/\ce{Mn^2+}};
\node[poporange, anchor=south, font=\small] at (axis cs:5, 850) {Tangentes};
\end{axis}
\end{tikzpicture}
\legende{Courbe potentiométrique $E = f(V)$ pour le dosage \ce{Fe^2+}/\ce{MnO4-}. Détermination graphique de $V_E$ par la méthode des tangentes (droite équidistante en violet).}
\end{popfigure}

\coursec{Exercices corrigés type Bac}

\begin{corrige}[Exercice 1 — QCM : les bons réflexes]
\begin{enumerate}
\item \textbf{Réponse b.} Une réaction de dosage doit être \cle{rapide} (résultat immédiat à chaque ajout), \cle{totale} (constante d'équilibre très grande) et \cle{unique} (aucune réaction parasite). Une réaction lente rend le repérage de l'équivalence impossible ; une réaction limitée fausse la relation stœchiométrique.
\item \textbf{Réponse b.} La solution titrante est celle de concentration connue, placée dans la burette : c'est la solution de \ce{MnO4-}. La solution de \ce{Fe^2+} est la solution titrée (erlenmeyer). L'acide sulfurique ne sert qu'à acidifier le milieu.
\item \textbf{Réponse b.} D'après l'équation $\ce{MnO4- + 5Fe^2+ + 8H+ -> Mn^2+ + 5Fe^3+ + 4H2O}$, il faut 5 mol de \ce{Fe^2+} pour 1 mol de \ce{MnO4-} : à l'équivalence, $n(\ce{Fe^2+}) = 5\,n(\ce{MnO4-})$. La réponse d est fausse : à l'équivalence les réactifs sont en proportions stœchiométriques, pas nécessairement « disparus » au sens strict.
\item \textbf{Réponse b.} Avant l'équivalence, le diiode (brun-jaune) est en excès ; à l'équivalence, tout le diiode a été consommé : on observe la \cle{décoloration} de la solution. La coloration violette persistante caractérise, elle, le dosage par le permanganate.
\end{enumerate}
\end{corrige}

\begin{corrige}[Exercice 2 — Vrai ou faux ?]
\begin{enumerate}
\item \textbf{Faux.} L'acide chlorhydrique est proscrit car les ions \ce{Cl-} sont oxydés par \ce{MnO4-} : $\ce{2MnO4- + 10Cl- + 16H+ -> 2Mn^2+ + 5Cl2 ^ + 8H2O}$. Cette réaction parasite consomme du titrant (et dégage un gaz toxique) : on acidifie avec l'acide sulfurique.
\item \textbf{Vrai.} L'ion \ce{MnO4-} est violet alors que \ce{Mn^2+} est quasi incolore : la première goutte de permanganate en excès colore la solution en rose-violet persistant. Le permanganate est donc à la fois titrant et indicateur coloré (auto-indication).
\item \textbf{Vrai.} C'est la définition même de l'équivalence : les réactifs ont été mélangés dans les proportions stœchiométriques de l'équation de la réaction de dosage.
\item \textbf{Faux.} Le thiosulfate se décompose sous l'action de la lumière, de l'acidité et des bactéries : sa solution doit être conservée à l'abri de la lumière (flacon ambré), tamponnée à pH basique et étalonnée avant usage.
\item \textbf{Faux.} Le volume équivalent se lit par différence entre le niveau final et le niveau initial du liquide dans la burette, au bas du ménisque, l'œil à hauteur de ce ménisque (erreur de parallaxe évitée).
\end{enumerate}
\end{corrige}

\begin{corrige}[Exercice 3 — Texte à trous : synthèse du vocabulaire]
Dans un dosage d'oxydoréduction, la solution de concentration connue est appelée solution \textbf{titrante} ; elle est placée dans la \textbf{burette graduée}. La solution de concentration inconnue est dite solution \textbf{titrée}. L'\textbf{équivalence} est atteinte lorsque les réactifs ont été mélangés dans les proportions \textbf{stœchiométriques}. Lors du dosage des ions \ce{Fe^2+} par les ions \ce{MnO4-}, l'équivalence est repérée par l'apparition d'une teinte \textbf{violette} persistante. Lors du dosage du diiode, l'équivalence est repérée par la \textbf{décoloration} de la solution, et le réactif titrant est la solution de \textbf{thiosulfate} de sodium.
\end{corrige}

\begin{corrige}[Exercice 4 — Définitions, couples et équations]
\begin{enumerate}
\item \textbf{Dosage} : opération quantitative consistant à déterminer la concentration d'une espèce (solution titrée) en la faisant réagir de façon totale avec une espèce de concentration connue (solution titrante). \textbf{Point d'équivalence} : état du système où les quantités de matière des réactifs introduits sont exactement dans les rapports stœchiométriques de l'équation bilan.
\item \textbf{Couples mis en jeu} :
\begin{itemize}
    \item \ce{MnO4-}/\ce{Mn^2+} (oxydant/réducteur) pour le titrant.
    \item \ce{Fe^3+}/\ce{Fe^2+} (oxydant/réducteur) pour le titré.
\end{itemize}
\item \textbf{Demi-équations et bilan pour \ce{I2}/\ce{S2O3^2-}} :
\begin{align*}
\text{Oxydation : } & \ce{2S2O3^2- -> S4O6^2- + 2e-} \\
\text{Réduction : } & \ce{I2 + 2e- -> 2I-} \\
\text{Bilan : } & \ce{I2 + 2S2O3^2- -> 2I- + S4O6^2-}
\end{align*}
\end{enumerate}
\end{corrige}

\begin{corrige}[Exercice 5 — Schéma du dispositif et rôle de l'acide]
\begin{enumerate}
\item Le schéma attendu est celui de la figure du cours : burette graduée contenant la solution titrante de \ce{KMnO4}, fixée sur un support, robinet au-dessus d'un erlenmeyer contenant la solution titrée de \ce{Fe^2+} acidifiée, le tout sur un agitateur magnétique. Légendes obligatoires : burette (titrant), robinet, erlenmeyer (titré + \ce{H2SO4}), barreau aimanté, agitateur.
\item On prélève \SI{20,0}{mL} avec une \cle{pipette jaugée} de \SI{20,0}{mL} (éventuellement munie d'un propipette). C'est la verrerie de précision adaptée : le volume de la prise d'essai intervient dans le calcul de $C_1$, il doit donc être connu avec exactitude (une éprouvette graduée serait trop imprécise, incertitude $\approx 100$ fois plus grande).
\item L'acide sulfurique apporte les ions \ce{H+} nécessaires à la demi-équation de réduction du permanganate : $\ce{MnO4- + 8H+ + 5e- -> Mn^2+ + 4H2O}$. Sans acidification, \ce{MnO4-} serait réduit en \ce{MnO2} (précipité brun) et la stœchiométrie serait faussée (3 électrons au lieu de 5). L'acide chlorhydrique est proscrit car ses ions \ce{Cl-} seraient oxydés en \ce{Cl2} par \ce{MnO4-} : réaction parasite consommant du titrant et dégageant un gaz toxique.
\end{enumerate}
\end{corrige}

\begin{corrige}[Exercice 6 — Dosage direct des ions fer(II)]
\begin{enumerate}
\item \textbf{Demi-équations} :
\[ \ce{MnO4- + 8H+ + 5e- -> Mn^2+ + 4H2O} \qquad \ce{Fe^2+ -> Fe^3+ + e-} \ (\times 5) \]
\textbf{Équation bilan} :
\[ \ce{MnO4- + 5Fe^2+ + 8H+ -> Mn^2+ + 5Fe^3+ + 4H2O} \]
\item \textbf{Relation à l'équivalence} :
\[ \frac{n(\ce{Fe^2+})}{5} = \frac{n(\ce{MnO4-})}{1} \quad \Longleftrightarrow \quad n(\ce{Fe^2+}) = 5\,n(\ce{MnO4-}) \]
soit $C_1 V_1 = 5\,C_2 V_E$.
\item \textbf{Application numérique} : $C_2 = \SI{0,0200}{mol.L^{-1}}$, $V_E = \SI{12,5}{mL}$, $V_1 = \SI{20,0}{mL}$.
\[ C_1 = \frac{5\,C_2 V_E}{V_1} = \frac{5 \times 0,0200 \times 12,5}{20,0} = \SI{0,0625}{mol.L^{-1}} \]
\item \textbf{Concentration massique en fer} : $M(\ce{Fe}) = \SI{55,8}{g.mol^{-1}}$.
\[ t = C_1 \times M(\ce{Fe}) = 0,0625 \times 55,8 = \SI{3,49}{g.L^{-1}} \]
\end{enumerate}
\textbf{Vérification :} $n(\ce{MnO4-}) = 0,0200 \times 12,5\times10^{-3} = \SI{2,50e-4}{mol}$ ; $n(\ce{Fe^2+}) = 5 \times 2,50\times10^{-4} = \SI{1,25e-3}{mol}$ ; $C_1 = 1,25\times10^{-3}/0,0200 = \SI{0,0625}{mol.L^{-1}}$. Cohérent.
\end{corrige}

\begin{corrige}[Exercice 7 — Dosage d'une solution de diiode]
\begin{enumerate}
\item \textbf{Demi-équations} : $\ce{I2 + 2e- -> 2I-}$ et $\ce{2S2O3^2- -> S4O6^2- + 2e-}$. \textbf{Bilan} :
\[ \ce{I2 + 2S2O3^2- -> 2I- + S4O6^2-} \]
\item \textbf{Calcul de concentration} : $C_2 = \SI{0,0500}{mol.L^{-1}}$, $V_E = \SI{8,40}{mL}$, $V_1 = \SI{10,0}{mL}$.
\begin{align*}
n(\ce{S2O3^2-}) &= C_2 V_E = 0,0500 \times 8,40\times10^{-3} = \SI{4,20e-4}{mol} \\
n(\ce{I2}) &= \frac{1}{2}\,n(\ce{S2O3^2-}) = \SI{2,10e-4}{mol} \\
C_1 &= \frac{n(\ce{I2})}{V_1} = \frac{2,10\times10^{-4}}{10,0\times10^{-3}} = \SI{0,0210}{mol.L^{-1}}
\end{align*}
\item \textbf{Observations} :
\begin{itemize}
    \item Avant l'équivalence : solution \textbf{brun-jaune} (diiode en excès), teinte s'atténuant progressivement (jaune paille).
    \item À l'équivalence : dernière trace de \ce{I2} consommée $\rightarrow$ \textbf{décoloration complète} (passage à l'incolore).
    \item Avec empois d'amidon (ajouté près de l'équivalence) : passage du \textbf{bleu foncé} (complexe amidon-\ce{I2}) à l'\textbf{incolore}, beaucoup plus net.
\end{itemize}
\end{enumerate}
\end{corrige}

\begin{corrige}[Exercice 8 — Exploitation d'une courbe potentiométrique]
\begin{enumerate}
\item \textbf{Allure de la courbe} $E = f(V)$ : croissante. Montée douce de 780 à 950 mV (0 à 8 mL), \textbf{saut brutal de potentiel} entre 9 et 11 mV (1080 à 1460 mV), puis palier vers 1470 mV.
\item \textbf{Méthode des tangentes} : on trace les tangentes aux branches lentes (avant et après le saut), puis la droite équidistante (parallèle, milieu vertical). Elle coupe la courbe au point d'inflexion (pente max).
\[ V_E \approx \SI{10,0}{mL} \quad (\text{milieu du saut } 9 \leftrightarrow 11\,\text{mL}) \]
\item \textbf{Calcul de $C_1$} : $C_2 = \SI{0,0100}{mol.L^{-1}}$, $V_1 = \SI{10,0}{mL}$.
\[ C_1 = \frac{5 \times C_2 \times V_E}{V_1} = \frac{5 \times 0,0100 \times 10,0}{10,0} = \SI{0,0500}{mol.L^{-1}} \]
\end{enumerate}
\textbf{Point de vigilance :} le point d'équivalence correspond au point d'inflexion du saut de potentiel (pente maximale), pas au milieu arithmétique des valeurs de $E$.
\end{corrige}

\begin{corrige}[Exercice 9 — Type Bac : comprimé de fer contre l'anémie]
\begin{enumerate}
\item \textbf{Équation-bilan} :
\[ \ce{MnO4- + 5Fe^2+ + 8H+ -> Mn^2+ + 5Fe^3+ + 4H2O} \]
\item \textbf{Dispositif} : burette graduée contenant la solution de \ce{KMnO4} (titrante) au-dessus d'un erlenmeyer contenant $V_1 = \SI{20,0}{mL}$ de $S$ acidifiée (titrée), sur agitateur magnétique (voir schéma du cours).
\item \textbf{Repérage de l'équivalence} : Avant l'équivalence, \ce{MnO4-} versé est entièrement consommé et réduit en \ce{Mn^2+} quasi incolore : la solution reste incolore (verdâtre pâle dû à \ce{Fe^2+}/\ce{Fe^3+}). À l'équivalence, la première goutte de \ce{MnO4-} en excès n'est plus consommée : apparition d'une \cle{teinte rose-violette persistante} (environ 30 s). Le permanganate joue le rôle d'indicateur coloré (auto-indication).
\item \textbf{Quantités de matière} : $C_2 = \SI{0,0100}{mol.L^{-1}}$, $V_E = \SI{15,8}{mL}$.
\begin{align*}
n(\ce{MnO4-}) &= C_2 V_E = 0,0100 \times 15,8\times10^{-3} = \SI{1,58e-4}{mol} \\
n(\ce{Fe^2+}) &= 5\,n(\ce{MnO4-}) = 5 \times 1,58\times10^{-4} = \SI{7,90e-4}{mol} \text{ dans la prise de } \SI{20,0}{mL}
\end{align*}
\item \textbf{Concentration de $S$ et masse de fer} :
\[ C_1 = \frac{7,90\times10^{-4}}{20,0\times10^{-3}} = \SI{3,95e-2}{mol.L^{-1}} \]
Masse de fer dans la fiole mère de \SI{100}{mL} (un comprimé dissous) :
\[ m(\ce{Fe}) = C_1 \times V_{\text{fiole}} \times M(\ce{Fe}) = 0,0395 \times 0,100 \times 55,8 = \SI{0,220}{g} = \SI{220}{mg} \]
\item \textbf{Contrôle de conformité} : Étiquette \SI{60}{mg}.
\[ \text{Écart relatif} = \frac{|220 - 60|}{60} \times 100 \approx 267\,\% \]
L'écart dépasse très largement la tolérance de 5\,\% (pharmacopée) : \textbf{le fabricant ne respecte pas la norme}. Le comprimé est fortement surdosé (risque de surdosage thérapeutique) ou l'étiquette est mensongère : signalement au contrôle qualité (MINSANTE/LANORM).
\end{enumerate}
\textbf{Barème indicatif (sur 6 pts) :} équation bilan 1 pt ; schéma légendé 1 pt ; repérage de l'équivalence 0,5 pt ; $n(\ce{MnO4-})$ et $n(\ce{Fe^2+})$ 1,5 pt ; $C_1$ et masse de fer 1,5 pt ; écart relatif et conclusion 0,5 pt.

\textbf{Points de vigilance :} ne pas oublier le facteur 5 de la stœchiométrie ; convertir les volumes en litres ; tenir compte du facteur de dilution (la fiole fait \SI{100}{mL}, la prise \SI{20,0}{mL}) ; exprimer la masse finale en mg pour la comparaison.
\end{corrige}

\begin{corrige}[Exercice 10 — Type Bac : degré iodométrique d'une teinture d'iode]
\begin{enumerate}
\item \textbf{Équation de dosage} :
\[ \ce{I2 + 2S2O3^2- -> 2I- + S4O6^2-} \]
\item \textbf{Observation et rôle de l'amidon} : À l'équivalence, tout le diiode est consommé : la solution passe du brun-jaune à l'\cle{incolore} (décoloration). L'empois d'amidon forme avec \ce{I2} un complexe bleu foncé très visible : ajouté près de l'équivalence (teinte jaune paille), il permet de repérer avec précision le passage du bleu foncé à l'incolore. On ne l'ajoute pas dès le début car le complexe amidon-iode, trop stable en présence d'un excès de diiode, libérerait \ce{I2} trop tardivement et fausserait $V_E$ (valeur trop grande).
\item \textbf{Quantités de matière} : $C_2 = \SI{0,100}{mol.L^{-1}}$, $V_E = \SI{15,6}{mL}$, $V_1' = \SI{20,0}{mL}$ (solution diluée $S'$).
\begin{align*}
n(\ce{S2O3^2-}) &= C_2 V_E = 0,100 \times 15,6\times10^{-3} = \SI{1,56e-3}{mol} \\
n(\ce{I2})_{\text{dans } S'} &= \frac{1}{2}\,n(\ce{S2O3^2-}) = \SI{7,80e-4}{mol}
\end{align*}
\item \textbf{Concentrations} :
\[ C' = \frac{7,80\times10^{-4}}{20,0\times10^{-3}} = \SI{3,90e-2}{mol.L^{-1}} \quad (\text{solution diluée } S') \]
La teinture commerciale est 10 fois plus concentrée (dilution 10 fois) :
\[ C = 10 \times C' = \SI{0,390}{mol.L^{-1}} \]
\item \textbf{Degré iodométrique} : $M(\ce{I2}) = 2 \times 126,9 = \SI{253,8}{g.mol^{-1}}$.
\[ t = C \times M(\ce{I2}) = 0,390 \times 253,8 \approx \SI{99,0}{g.L^{-1}} \]
soit \SI{9,90}{g} pour \SI{100}{mL} : \textbf{degré iodométrique $\approx 9,9^\circ\text{I}$}.
Écart relatif par rapport à l'étiquette (5°I) :
\[ \text{écart} = \frac{|9,9 - 5|}{5} \times 100 \approx 98\,\% \]
Très supérieur à la tolérance de 5\,\% : \textbf{l'indication de l'étiquette est incorrecte} (solution environ deux fois plus concentrée qu'annoncé).
\end{enumerate}
\textbf{Barème indicatif (sur 5 pts) :} équation 0,5 pt ; observation + rôle de l'amidon 1 pt ; $n(\ce{S2O3^2-})$ et $n(\ce{I2})$ 1 pt ; concentrations $C'$ et $C$ 1,5 pt ; degré iodométrique et conclusion 1 pt.

\textbf{Points de vigilance :} facteur $\frac{1}{2}$ de la stœchiométrie ; ne pas oublier le facteur de dilution ($\times 10$) ; le degré iodométrique s'exprime en grammes de \ce{I2} pour \SI{100}{mL}, pas pour 1 L.
\end{corrige}

\begin{corrige}[Exercice 11 — Type Bac : degré chlorométrique d'une eau de Javel]
\begin{enumerate}
\item \textbf{Justification du dosage en retour (indirect)} : On parle de \cle{dosage en retour} car l'espèce à doser, l'hypochlorite \ce{ClO-}, ne réagit pas directement (ou trop lentement) avec le titrant \ce{S2O3^2-}. Elle oxyde d'abord un excès d'ions \ce{I-} (ajouté en grande quantité) en milieu acide, produisant du diiode \ce{I2} dosable. On remonte ensuite à \ce{ClO-} par chaînage des stœchiométries.
\item \textbf{Équation du dosage du diiode} (réaction rapide, totale) :
\[ \ce{I2 + 2S2O3^2- -> 2I- + S4O6^2-} \]
\item \textbf{Quantité de diiode formée} : $C_2 = \SI{0,100}{mol.L^{-1}}$, $V_E = \SI{11,2}{mL}$.
\begin{align*}
n(\ce{S2O3^2-}) &= C_2 V_E = 0,100 \times 11,2\times10^{-3} = \SI{1,12e-3}{mol} \\
n(\ce{I2}) &= \frac{1}{2}\,n(\ce{S2O3^2-}) = \SI{5,60e-4}{mol}
\end{align*}
\item \textbf{Remontée à l'hypochlorite} : Réaction de libération de l'iode :
\[ \ce{ClO- + 2I- + 2H+ -> Cl- + I2 + H2O} \]
1 mol \ce{ClO-} produit 1 mol \ce{I2} $\Rightarrow$ $n(\ce{ClO-}) = n(\ce{I2}) = \SI{5,60e-4}{mol}$ dans la prise de \SI{10,0}{mL} de solution diluée $S$.
\[ [\ce{ClO-}]_S = \frac{5,60\times10^{-4}}{10,0\times10^{-3}} = \SI{5,60e-2}{mol.L^{-1}} \]
Eau de Javel commerciale (diluée 20 fois : $V_{\text{mère}} = \SI{5,0}{mL} \rightarrow \SI{100}{mL}$) :
\[ [\ce{ClO-}]_{\text{commerciale}} = 20 \times 0,0560 = \SI{1,12}{mol.L^{-1}} \]
\item \textbf{Degré chlorométrique} : D'après $\ce{ClO- + 2H+ + Cl- -> Cl2 ^ + H2O}$, 1 mol \ce{ClO-} libère 1 mol \ce{Cl2}.
Pour 1 L d'eau de Javel commerciale : $n(\ce{Cl2}) = n(\ce{ClO-}) = \SI{1,12}{mol}$.
Volume de chlore gazeux aux CNTP ($V_m = \SI{22,4}{L.mol^{-1}}$) :
\[ V(\ce{Cl2}) = n \times V_m = 1,12 \times 22,4 \approx \SI{25,1}{L} \]
Le degré chlorométrique est donc d'environ \textbf{25,1 °chl}.
\item \textbf{Contrôle de conformité} : Étiquette 12 °chl.
\[ \text{écart} = \frac{|25,1 - 12|}{12} \times 100 \approx 109\,\% \]
L'écart est considérable : \textbf{l'étiquette est fausse}, l'eau de Javel est environ deux fois plus concentrée qu'annoncé. Une telle surconcentration présente un danger pour l'utilisateur (brûlures cutanées, dégagement de \ce{Cl2} toxique en cas de mélange accidentel avec un acide détartrant) et doit être signalée (Direction de la Concurrence, MINCOMMERCE).
\end{enumerate}
\textbf{Barème indicatif (sur 6 pts) :} justification du dosage indirect 0,5 pt ; équation de dosage 0,5 pt ; $n(\ce{I2})$ 1 pt ; $n(\ce{ClO-})$ et concentrations 1,5 pt ; degré chlorométrique 1,5 pt ; écart relatif et conclusion 1 pt.

\textbf{Points de vigilance :} bien enchaîner les relations $n(\ce{ClO-}) = n(\ce{I2}) = \frac{1}{2}n(\ce{S2O3^2-})$ ; ne pas oublier le facteur de dilution ($\times 20$) ; le degré chlorométrique se calcule pour \textbf{un litre} d'eau de Javel commerciale, pas pour la solution diluée $S$ ; utiliser $V_m = \SI{22,4}{L.mol^{-1}}$ (CNTP).
\end{corrige}

\newpage

\coursec{Fiche bilan}

\coursec{Généralités sur les dosages}

\begin{definition}[Vocabulaire du dosage]
Un \cle{dosage} est une opération qui permet de déterminer la concentration (ou la quantité de matière) d'une espèce chimique en solution, appelée \cle{solution titrée}, en la faisant réagir avec une solution de concentration connue, appelée \cle{solution titrante}. La réaction utilisée doit être \cle{totale}, \cle{rapide} et \cle{unique}.
Le \cle{point d'équivalence} (noté $E$) est l'instant où les réactifs ont été introduits dans les proportions stœchiométriques exactes. Il est repéré expérimentalement par un changement brutal d'une propriété physique (couleur, conductivité, pH, potentiel).
\end{definition}

\begin{propriete}[Relation fondamentale]
À l'équivalence, les quantités de matière des réactifs sont proportionnelles à leurs coefficients stœchiométriques. Pour une réaction :
\[ \ce{a A + b B -> produits} \]
on a :
\[ \frac{n_A(E)}{a} = \frac{n_B(E)}{b} \quad \text{soit} \quad n_A(E) = \frac{a}{b}\, n_B(E) \]
Avec $n = C \times V$ (volume en \si{\liter}), on en déduit la concentration inconnue.
\end{propriete}

\begin{methode}[Étapes générales d'un dosage]
\begin{enumerate}
  \item Choisir la réaction de dosage (totale, rapide, unique, repérable).
  \item Préparer le matériel : burette (titrante), erlenmeyer (titrée + indicateur si besoin), agitateur magnétique.
  \item Effectuer un dosage grossier pour localiser $V_E$ (facultatif mais recommandé).
  \item Effectuer au moins deux dosages précis (volume ajouté goutte à goutte près de $V_E$).
  \item Noter $V_E$ (volume à l'équivalence) lu au bas du ménisque, au \SI{0,1}{\milli\liter} près.
  \item Exploiter les résultats : équation-bilan $\to$ tableau d'avancement $\to$ relation à $E$ $\to$ calcul de $C_{\text{inconnue}}$ $\to$ conclusion.
\end{enumerate}
\end{methode}

\coursec{Conditions d'un bon dosage et dispositif}

\begin{propriete}[Quatre conditions indispensables]
Pour qu'un dosage soit fiable, la réaction chimique doit satisfaire :
\begin{enumerate}
  \item \textbf{Totale} : le rendement est quasi 100\,\% (constante d'équilibre $K$ très grande). Les réactifs sont totalement consommés à l'équivalence.
  \item \textbf{Rapide} : la réaction est quasi instantanée à température ambiante (pas d'attente entre l'ajout et la mesure).
  \item \textbf{Unique} : aucune réaction parasite ne consomme les réactifs (ex. : pas d'oxydation du \ce{Cl-} par \ce{MnO4-} en milieu chlorhydrique).
  \item \textbf{Équivalence facilement repérable} : changement de couleur net (indicateur coloré ou auto-indication) au voisinage immédiat de $E$.
\end{enumerate}
\end{propriete}

\begin{experience}[Dispositif expérimental type]
\begin{popfigure}
\begin{tikzpicture}[
  scale=0.9,
  every node/.style={transform shape},
  burette/.style={draw=popdark, thick, fill=popblue!10},
  erlenmeyer/.style={draw=popdark, thick, fill=popgreen!10},
  stand/.style={draw=popdark, very thick},
  liquid/.style={fill=popblue!60, opacity=0.8},
  liquid2/.style={fill=poporange!60, opacity=0.8},
]
% Support
\draw[stand] (-3,0) -- (-3,5.5) -- (1,5.5) -- (1,5);
\draw[stand] (-3,0) -- (-4,0) -- (-4,-0.5) -- (-2,-0.5) -- (-2,0);
% Burette
\draw[burette] (-2.5,5) -- (-2.5,1) -- (-1.5,1) -- (-1.5,5);
\draw[liquid] (-2.5,1) -- (-2.5,4.2) -- (-1.5,4.2) -- (-1.5,1) -- cycle;
\node at (-2,4.5) {\small \SI{0}{\milli\liter}};
\node at (-2,1.3) {\small \SI{50}{\milli\liter}};
\draw[thick] (-2.5,1) -- (-2,0.5) -- (-1.5,1); % robinet
\draw[popblue, thick, decorate, decoration={snake, amplitude=1pt, segment length=3pt}] (-2,0.5) -- (-2,0);
% Erlenmeyer
\draw[erlenmeyer] (2,0) -- (1.5,1.5) -- (0.5,1.5) -- (0,0) -- cycle;
\draw[erlenmeyer] (0.5,1.5) -- (0.5,2) -- (1,2.2) -- (1,2.5) -- (1.5,2.5) -- (1.5,1.5);
\draw[liquid2] (0.2,0.2) -- (1.3,0.2) .. controls (1.4,0.5) and (1.4,1.2) .. (1.3,1.5) -- (0.7,1.5) .. controls (0.6,1.2) and (0.6,0.5) .. (0.7,0.2) -- cycle;
% Agitateur
\draw[fill=popgray!50, draw=popdark] (2.5,-0.5) rectangle (4.5,0);
\draw[->, thick, poporange] (3.5,-0.2) arc (0:360:0.3 and 0.1);
\node[below] at (3.5,-0.7) {\small Agitateur magnétique};
% Légendes
\node[align=center, font=\bfseries\small] at (-2,5.3) {Burette\\(Solution titrante)};
\node[align=center, font=\bfseries\small] at (1,2.8) {Erlenmeyer\\(Solution titrée + indicateur)};
\end{tikzpicture}
\legende{Dispositif de dosage : burette graduée (classe B, \SI{50}{\milli\liter} ou \SI{25}{\milli\liter}), erlenmeyer, agitateur magnétique. La lecture du volume se fait au bas du ménisque, l'œil au niveau du liquide.}
\end{popfigure}
\end{experience}

\begin{attention}[Pièges classiques au montage]
\begin{itemize}
  \item Ne pas rincer la burette avec la solution titrante \textbf{après} l'avoir rincée à l'eau distillée (sinon dilution $\to$ $C_{\text{titrante}}$ faussée $\to$ erreur systématique).
  \item Chasser les bulles d'air dans la pointe de la burette avant de noter $V_i$.
  \item Placer l'erlenmeyer sur une feuille blanche (ou carrelage noir pour \ce{MnO4-}) pour mieux voir le virage.
  \item Ajouter la titrante goutte à goutte près de l'équivalence en secouant énergiquement l'erlenmeyer.
\end{itemize}
\end{attention}

\coursec{Dosage des ions fer(II) par les ions permanganate}

\begin{savaistu}[Contexte camerounais]
Le dosage des ions \ce{Fe^{2+}} est crucial dans le contrôle qualité des engrais (ex. : sulfate de fer(II) pour la correction des sols acides dans les plantations de café ou de cacao au Sud), dans l'analyse des eaux de forage (teneur en fer) ou dans l'industrie pharmaceutique locale (suppléments en fer). Le permanganate de potassium \ce{KMnO4} est l'oxydant de référence car il est son propre indicateur.
\end{savaistu}

\begin{propriete}[Couples rédox et demi-équations (milieu acide)]
\begin{itemize}
  \item Couple \ce{MnO4-}/\ce{Mn^{2+}} (oxydant) : $\ce{MnO4- + 8H+ + 5e- -> Mn^{2+} + 4H2O}$ \quad ($E^\circ = \SI{+1,51}{\volt}$)
  \item Couple \ce{Fe^{3+}}/\ce{Fe^{2+}} (réducteur) : $\ce{Fe^{3+} + e- -> Fe^{2+}}$ \quad ($E^\circ = \SI{+0,77}{\volt}$)
\end{itemize}
L'oxydation des \ce{Fe^{2+}} par \ce{MnO4-} est spontanée ($\Delta E^\circ > 0$).
\end{propriete}

\begin{propriete}[Équation-bilan de la réaction de dosage]
En combinant les demi-équations (multiplier la seconde par 5), on obtient :
\[ \ce{MnO4- + 5Fe^{2+} + 8H+ -> Mn^{2+} + 5Fe^{3+} + 4H2O} \]
\begin{itemize}
  \item \textbf{Milieu} : Acide sulfurique \ce{H2SO4} (\SI{1}{\mole\per\liter} environ). \textbf{Jamais} d'acide chlorhydrique (\ce{Cl-} oxydé en \ce{Cl2} par \ce{MnO4-}) ni d'acide nitrique (\ce{NO3-} oxydant parasite).
  \item \textbf{Couleur} : \ce{MnO4-} (violet foncé) $\to$ \ce{Mn^{2+}} (incolore). \ce{Fe^{2+}} (vert pâle) $\to$ \ce{Fe^{3+}} (jaune/brun).
  \item \textbf{Auto-indication} : Tant qu'il y a du \ce{Fe^{2+}}, le \ce{MnO4-} ajouté est décoloré instantanément. Au premier excès de \ce{MnO4-}, la solution vire au \textbf{rose pâle persistant} (environ \SI{30}{\second}).
\end{itemize}
\end{propriete}

\begin{methode}[Protocole opératoire : Dosage \ce{Fe^{2+}} par \ce{MnO4-}]
\begin{enumerate}
  \item Pipeter \SI{20,0}{\milli\liter} de solution de \ce{Fe^{2+}} (titrée) dans un erlenmeyer.
  \item Ajouter \SI{20}{\milli\liter} d'acide sulfurique \SI{1}{\mole\per\liter} (milieu acide + volume suffisant).
  \item Remplir la burette (rincée 2$\times$ avec la titrante) de solution de \ce{KMnO4} (titrante). Noter $V_i$.
  \item Titrer sous agitation magnétique jusqu'à l'apparition d'une coloration \textbf{rose pâle persistante} (\SI{30}{\second}). Noter $V_f$.
  \item Recommencer pour obtenir deux $V_E$ concordants (écart $< \SI{0,2}{\milli\liter}$).
\end{enumerate}
\end{methode}

\begin{exemplebox}[Exemple chiffré : Contrôle d'un lot de sulfate de fer(II) heptahydraté]
On dose \SI{20,0}{\milli\liter} d'une solution de \ce{Fe^{2+}} (préparée en dissolvant \SI{1,50}{\gram} de sel dans \SI{250}{\milli\liter}) par une solution de \ce{KMnO4} de concentration $C_{\text{ox}} = \SI{0,0200}{\mole\per\liter}$.
Les volumes à l'équivalence trouvés sont : $V_{E1} = \SI{12,4}{\milli\liter}$ ; $V_{E2} = \SI{12,5}{\milli\liter}$.
\begin{enumerate}
  \item \textbf{Moyenne} : $V_E = \SI{12,45}{\milli\liter} = \SI{1,245e-2}{\liter}$.
  \item \textbf{Quantité de matière de \ce{MnO4-}} : $n(\ce{MnO4-}) = C_{\text{ox}} \times V_E = \SI{0,0200}{\mole\per\liter} \times \SI{1,245e-2}{\liter} = \SI{2,49e-4}{\mole}$.
  \item \textbf{Tableau d'avancement} :
  \begin{center}
  \begin{tabular}{|c|c|c|c|c|c|}\hline
  & \ce{MnO4-} & \ce{5Fe^{2+}} & \ce{8H+} & \ce{Mn^{2+}} & \ce{5Fe^{3+}} \\ \hline
  État initial & $n_{\text{ox}}$ & $n_{\text{réd}}$ & excès & 0 & 0 \\ \hline
  État final & 0 & 0 & excès & $n_{\text{ox}}$ & $5n_{\text{ox}}$ \\ \hline
  \end{tabular}
  \end{center}
  À l'équivalence : $n(\ce{Fe^{2+}}) = 5 \times n(\ce{MnO4-}) = 5 \times \SI{2,49e-4}{\mole} = \SI{1,245e-3}{\mole}$.
  \item \textbf{Concentration molaire} : $C(\ce{Fe^{2+}}) = \frac{n(\ce{Fe^{2+}})}{V_{\text{réd}}} = \frac{\SI{1,245e-3}{\mole}}{\SI{20,0e-3}{\liter}} = \SI{0,0623}{\mole\per\liter}$.
  \item \textbf{Titre massique} (si demandé) : $M(\ce{FeSO4.7H2O}) = \SI{278,0}{\gram\per\mole}$.
  $t = C \times M = \SI{0,0623}{\mole\per\liter} \times \SI{278,0}{\gram\per\mole} = \SI{17,3}{\gram\per\liter}$.
  Masse dans \SI{250}{\milli\liter} = $t \times 0,250 = \SI{4,33}{\gram}$. (Écart par rapport aux \SI{1,50}{\gram} pesés $\to$ impuretés ou erreur de manipulation).
\end{enumerate}
\end{exemplebox}

\begin{attention}[Erreurs fréquentes \ce{Fe^{2+}}/\ce{MnO4-}]
\begin{itemize}
  \item Oublier l'acide sulfurique : formation de \ce{MnO2} (précipité brun) au lieu de \ce{Mn^{2+}}, réaction incomplète.
  \item Utiliser \ce{HCl} : dégagement de \ce{Cl2} (toxique, vert-jaune), fausse $V_E$.
  \item Lire le volume au sommet du ménisque (liquide coloré opaque) $\to$ lire au bas du ménisque en se plaçant de côté ou utiliser une burette à fond blanc gradué "en négatif".
  \item Ne pas attendre la persistance de la couleur rose (\SI{30}{\second}) $\to$ surdosage.
\end{itemize}
\end{attention}

\coursec{Dosage du diiode par les ions thiosulfate}

\begin{savaistu}[Applications locales]
Ce dosage (iodométrie) est la méthode de référence pour titrer le \textbf{chlore actif} dans l'eau de Javel (production locale type \textit{Socochim} ou traitement d'eau \textit{CAMWATER}), pour doser l'iode dans le sel iodé (lutte contre le goitre, programme santé publique), ou pour doser la vitamine C dans le jus de \textbf{bissap} (\textit{Hibiscus sabdariffa}) par iodométrie indirecte.
\end{savaistu}

\begin{propriete}[Couples rédox et demi-équations]
\begin{itemize}
  \item Couple \ce{I2}/\ce{I-} (oxydant) : $\ce{I2 + 2e- -> 2I-}$ \quad ($E^\circ = \SI{+0,54}{\volt}$)
  \item Couple \ce{S4O6^{2-}}/\ce{S2O3^{2-}} (réducteur) : $\ce{S4O6^{2-} + 2e- -> 2S2O3^{2-}}$ \quad ($E^\circ = \SI{+0,08}{\volt}$)
\end{itemize}
La réaction est totale et rapide en milieu neutre ou faiblement acide.
\end{propriete}

\begin{propriete}[Équation-bilan et indicateur]
\[ \ce{I2 + 2S2O3^{2-} -> 2I- + S4O6^{2-}} \]
\begin{itemize}
  \item \textbf{Indicateur} : \textbf{Empois d'amidon} (solution colloïdale d'amidon). L'amidon forme un complexe \textbf{bleu-noir intense} avec le diiode \ce{I2} (même à très faible concentration).
  \item \textbf{Protocole indicateur} : On n'ajoute l'empois d'amidon \textbf{qu'au voisinage de l'équivalence} (quand la couleur jaune/brun de l'iode devient pâle). Si on l'ajoute au début, le complexe \ce{I2}-amidon est trop stable et ne se décompose pas complètement $\to$ erreur sur $V_E$.
  \item \textbf{Repérage} : Décoloration brutale du bleu-noir $\to$ solution incolore.
\end{itemize}
\end{propriete}

\begin{methode}[Protocole opératoire : Dosage \ce{I2} par \ce{S2O3^{2-}}]
\begin{enumerate}
  \item Verser un volume connu $V(\ce{I2})$ de solution de diiode (titrée) dans l'erlenmeyer.
  \item Remplir la burette (rincée) de solution de thiosulfate de sodium \ce{Na2S2O3} (titrante). Noter $V_i$.
  \item Titrer sous agitation. La solution est jaune/brun (couleur \ce{I2}).
  \item Quand la couleur devient \textbf{jaune pâle (paille)}, ajouter \SI{1}{\milli\liter} à \SI{2}{\milli\liter} d'empois d'amidon $\to$ virée \textbf{bleu-noir}.
  \item Continuer goutte à goutte jusqu'à \textbf{décoloration complète} (solution incolore). Noter $V_f$.
  \item Recommencer pour concordance.
\end{enumerate}
\end{methode}

\begin{exemplebox}[Exemple chiffré : Titrage de l'eau de Javel (chlore actif)]
On prélève \SI{10,0}{\milli\liter} d'eau de Javel commerciale. On la dilue au \SI{100}{\milli\liter} (flacon jaugé). On prend \SI{20,0}{\milli\liter} de cette dilution, on ajoute un excès d'iodure de potassium \ce{KI} et de l'acide sulfurique : $\ce{ClO- + 2I- + 2H+ -> I2 + Cl- + H2O}$. Le \ce{I2} libéré est dosé par \ce{Na2S2O3} \SI{0,100}{\mole\per\liter}.
$V_E = \SI{15,2}{\milli\liter}$ (moyenne de deux essais concordants).
\begin{enumerate}
  \item $n(\ce{S2O3^{2-}}) = \SI{0,100}{\mole\per\liter} \times \SI{15,2e-3}{\liter} = \SI{1,52e-3}{\mole}$.
  \item Relation : $n(\ce{S2O3^{2-}}) = 2 \times n(\ce{I2}) \implies n(\ce{I2}) = \SI{7,60e-4}{\mole}$.
  \item $n(\ce{ClO-}) = n(\ce{I2}) = \SI{7,60e-4}{\mole}$ (dans les \SI{20,0}{\milli\liter} de dilution).
  \item Dans le flacon \SI{100}{\milli\liter} : $n(\ce{ClO-})_{\text{total}} = 5 \times \SI{7,60e-4}{\mole} = \SI{3,80e-3}{\mole}$.
  \item Dans l'aliquote \SI{10,0}{\milli\liter} initiale : $C(\ce{ClO-}) = \frac{\SI{3,80e-3}{\mole}}{\SI{10,0e-3}{\liter}} = \SI{0,380}{\mole\per\liter}$.
  \item Titre massique chlore actif (exprimé en \ce{Cl2} équivalent, $M(\ce{Cl2}) = \SI{71,0}{\gram\per\mol}$) :
  $t = C \times M = \SI{0,380}{\mole\per\liter} \times \SI{71,0}{\gram\per\mol} = \SI{27,0}{\gram\per\liter}$.
  \item Conclusion : L'eau de Javel titre \SI{27}{\gram\per\liter} de chlore actif (norme souvent \SI{25}{\gram\per\liter} à \SI{30}{\gram\per\liter}).
\end{enumerate}
\end{exemplebox}

\begin{attention}[Erreurs fréquentes \ce{I2}/\ce{S2O3^{2-}}]
\begin{itemize}
  \item Ajouter l'amidon trop tôt $\to$ complexation forte, équivalence non nette, $V_E$ surestimé.
  \item Solution de thiosulfate non stable (se décompose en \ce{SO2} et \ce{S} si acide ou bactéries). La préparer fraîchement ou ajouter un peu de \ce{Na2CO3} / \ce{NaOH} pour tamponner pH $\approx 9$. L'étalonner avant usage (sur \ce{KIO3}/\ce{KI} ou \ce{K2Cr2O7}/\ce{KI}).
  \item Oublier le facteur de dilution dans les calculs (ex. eau de Javel).
\end{itemize}
\end{attention}

\coursec{Exploitation d'un dosage : méthode complète}

\begin{methode}[Méthode rigoureuse en 5 étapes (type Bac)]
\begin{enumerate}
  \item \textbf{Équation-bilan} : Écrire l'équation équilibrée (atomes + charges). Identifier oxydant/réducteur.
  \item \textbf{Tableau d'avancement} : Construire le tableau (État initial / Intermédiaire / Final). Utiliser $x$ pour l'avancement.
  \item \textbf{Relation à l'équivalence} : Exprimer $x_{\text{max}}$ (ou $x_E$) et écrire la relation stœchiométrique entre $n_{\text{titré}}$ et $n_{\text{titrant}}$.
  \item \textbf{Calculs} :
  \begin{itemize}
    \item $n_{\text{titrant}} = C_{\text{titrant}} \times V_E$ ($V_E$ en \si{\liter}).
    \item $n_{\text{titré}}$ via la relation stœchiométrique.
    \item $C_{\text{titré}} = n_{\text{titré}} / V_{\text{titré}}$ ($V_{\text{titré}}$ en \si{\liter}).
    \item Si titre massique $t$ demandé : $t = C \times M$ (en \si{\gram\per\liter}).
  \end{itemize}
  \item \textbf{Conclusion} : Phrase complète : « La concentration molaire de la solution titrée est $C = \ldots\ \si{\mole\per\liter}$ » (arrondi au nombre de chiffres significatifs le plus faible des données, généralement 3).
\end{enumerate}
\end{methode}

\begin{exoresolu}[Dosage d'une solution de permanganate étalonnée au sel de Mohr]
\textbf{Énoncé.} On souhaite étalonner une solution de \ce{KMnO4} (titrante) à l'aide d'une solution étalon de \ce{Fe^{2+}} issue du sel de Mohr \ce{(NH4)2Fe(SO4)2,6H2O} (titrée).\\
On a préparé la solution titrée en dissolvant \SI{0,980}{\gram} de sel de Mohr dans de l'eau acidifiée (\ce{H2SO4}) et complété à \SI{250,0}{\milli\liter} (flacon jaugé).\\
On prélève \SI{20,0}{\milli\liter} de cette solution, on ajoute \SI{20}{\milli\liter} \ce{H2SO4} \SI{1}{\mole\per\liter}, on titre par la \ce{KMnO4}.\\
Résultats : $V_{E1} = \SI{22,35}{\milli\liter}$ ; $V_{E2} = \SI{22,40}{\milli\liter}$.
\textbf{Données :} $M(\ce{Fe}) = \SI{55,8}{\gram\per\mol}$ ; $M(\text{sel de Mohr}) = \SI{392,1}{\gram\per\mol}$.

\tcblower
\textbf{Solution détaillée.}
\begin{enumerate}
  \item \textbf{Équation-bilan :} $\ce{MnO4- + 5Fe^{2+} + 8H+ -> Mn^{2+} + 5Fe^{3+} + 4H2O}$.
  \item \textbf{Concentration de la solution titrée (Fe2+) :}
  $n(\text{sel de Mohr}) = \frac{\SI{0,980}{\gram}}{\SI{392,1}{\gram\per\mol}} = \SI{2,500e-3}{\mol}$.
  $C(\ce{Fe^{2+}}) = \frac{\SI{2,500e-3}{\mol}}{\SI{0,2500}{\liter}} = \SI{0,01000}{\mole\per\liter}$.
  \item \textbf{Volume moyen à l'équivalence :} $V_E = \frac{22,35 + 22,40}{2} = \SI{22,38}{\milli\liter} = \SI{2,238e-2}{\liter}$.
  \item \textbf{Tableau d'avancement (dans l'erlenmeyer) :}
  \begin{center}
  \begin{tabular}{|c|c|c|c|c|c|}\hline
  Avancement $x$ & \ce{MnO4-} & \ce{5Fe^{2+}} & \ce{8H+} & \ce{Mn^{2+}} & \ce{5Fe^{3+}} \\ \hline
  $x=0$ & $n_{\text{ox}}$ & $n_{\text{réd}}$ & excès & 0 & 0 \\ \hline
  $x=x_E$ & 0 & 0 & excès & $n_{\text{ox}}$ & $n_{\text{réd}}$ \\ \hline
  \end{tabular}
  \end{center}
  Relation : $n_{\text{réd}} = 5\, n_{\text{ox}} \implies C_{\text{réd}} V_{\text{réd}} = 5\, C_{\text{ox}} V_E$.
  \item \textbf{Calcul de $C_{\text{ox}}$ :}
  $C_{\text{ox}} = \frac{C_{\text{réd}} V_{\text{réd}}}{5 V_E} = \frac{\SI{0,01000}{\mole\per\liter} \times \SI{20,0e-3}{\liter}}{5 \times \SI{22,38e-3}{\liter}} = \frac{\SI{2,000e-4}{\mol}}{\SI{111,9e-3}{\liter}} = \SI{1,787e-3}{\mole\per\liter}$.
  \item \textbf{Conclusion :} La concentration molaire de la solution de permanganate de potassium est $\boldsymbol{C(\ce{KMnO4}) = \SI{1,79e-3}{\mole\per\liter}}$ (3 chiffres significatifs).
\end{enumerate}
\end{exoresolu}

\begin{exercice}[Entraînement : Dosage de vitamine C dans le jus de bissap]
On souhaite doser la vitamine C (acide ascorbique, \ce{C6H8O6}, $M = \SI{176,1}{\gram\per\mol}$) dans un jus de bissap artisanal.
\textbf{Protocole :} On introduit \SI{20,0}{\milli\liter} de jus de bissap filtré dans un erlenmeyer. On ajoute un excès d'iode (\ce{I2}) en solution connue ($C(\ce{I2}) = \SI{5,00e-3}{\mole\per\liter}$, volume ajouté $V_{\ce{I2}} = \SI{25,0}{\milli\liter}$). La vitamine C réduit l'iode : $\ce{C6H8O6 + I2 -> C6H6O6 + 2I- + 2H+}$.
L'iode en excès est dosé par une solution de thiosulfate de sodium \ce{Na2S2O3} de concentration $C(\ce{S2O3^{2-}}) = \SI{0,0100}{\mole\per\liter}$.
Volume de thiosulfate à l'équivalence : $V_E = \SI{18,5}{\milli\liter}$.
\textbf{Questions :}
\begin{enumerate}
  \item Écrire l'équation-bilan de la réaction de dosage de l'excès d'iode.
  \item Calculer la quantité de matière d'iode initialement introduite.
  \item Calculer la quantité de matière d'iode réagie avec la vitamine C.
  \item En déduire la concentration molaire et le titre massique (en \si{\gram\per\liter}) de vitamine C dans le jus.
\end{enumerate}
\end{exercice}

\begin{corrige}[Exercice : Dosage de vitamine C dans le jus de bissap]
\begin{enumerate}
  \item \textbf{Équation dosage excès :} $\ce{I2 + 2S2O3^{2-} -> 2I- + S4O6^{2-}}$.
  \item \textbf{Iode initial :} $n_{\text{init}}(\ce{I2}) = C(\ce{I2}) \times V_{\ce{I2}} = \SI{5,00e-3}{\mole\per\liter} \times \SI{25,0e-3}{\liter} = \SI{1,25e-4}{\mol}$.
  \item \textbf{Iode en excès (dosé) :}
  $n(\ce{S2O3^{2-}}) = \SI{0,0100}{\mole\per\liter} \times \SI{18,5e-3}{\liter} = \SI{1,85e-4}{\mol}$.
  $n_{\text{excès}}(\ce{I2}) = \frac{n(\ce{S2O3^{2-}})}{2} = \SI{9,25e-5}{\mol}$.
  \item \textbf{Iode consommé par vitamine C :}
  $n_{\text{réagi}}(\ce{I2}) = n_{\text{init}} - n_{\text{excès}} = \SI{1,25e-4}{\mol} - \SI{0,925e-4}{\mol} = \SI{3,25e-5}{\mol}$.
  \textbf{Vitamine C :} $\ce{C6H8O6 + I2 -> ...}$ $\implies$ $n(\text{Vit C}) = n_{\text{réagi}}(\ce{I2}) = \SI{3,25e-5}{\mol}$.
  \textbf{Concentration :} $C(\text{Vit C}) = \frac{\SI{3,25e-5}{\mol}}{\SI{20,0e-3}{\liter}} = \SI{1,63e-3}{\mole\per\liter}$.
  \textbf{Titre massique :} $t = C \times M = \SI{1,63e-3}{\mole\per\liter} \times \SI{176,1}{\gram\per\mol} = \SI{0,287}{\gram\per\liter} = \SI{28,7}{\milli\gram\per\centi\liter}$.
  \item \textbf{Conclusion :} Le jus de bissap contient \SI{28,7}{\milli\gram} de vitamine C par \SI{100}{\milli\liter} (valeur cohérente avec la littérature : \SI{20}{\milli\gram} à \SI{40}{\milli\gram}/\SI{100}{\milli\liter}).
\end{enumerate}
\end{corrige}

\begin{popfigure}
\begin{tikzpicture}[
  central/.style={rectangle, rounded corners=8pt, draw=popdark, fill=popblue!25, text width=3.4cm, align=center, font=\bfseries\small, inner sep=6pt},
  branche/.style={rectangle, rounded corners=5pt, text width=3.6cm, align=center, font=\footnotesize, inner sep=4pt},
  lien/.style={thick, -{Stealth[length=2mm]}}
]
\node[central] (C) {Dosages d'oxydoréduction};

\node[branche, draw=popblue, fill=popblueL, above left=1.6cm and 3.6cm of C] (B1) {\textbf{Généralités}\\ Dosage, titrée, titrante, équivalence};
\node[branche, draw=poporange, fill=poporangeL, above=2.2cm of C] (B2) {\textbf{Conditions}\\ Réaction totale, rapide, unique, repérage de $E$};
\node[branche, draw=popgreen, fill=popgreenL, above right=1.6cm and 3.6cm of C] (B3) {\textbf{Dispositif}\\ Burette, erlenmeyer, agitateur};
\node[branche, draw=poppurple, fill=poppurpleL, below left=1.6cm and 3.6cm of C] (B4) {\textbf{Fe$^{2+}$ / MnO$_4^-$}\\ Virage incolore $\to$ rose persistant\\ Milieu \ce{H2SO4} obligatoire};
\node[branche, draw=popgold, fill=popgoldL, below=2.2cm of C] (B5) {\textbf{I$_2$ / S$_2$O$_3^{2-}$}\\ Décoloration bleu-noir $\to$ incolore\\ Empois d'amidon près de $E$};
\node[branche, draw=poppink, fill=poppinkL, below right=1.6cm and 3.6cm of C] (B6) {\textbf{Exploitation}\\ $n_{\text{ox}} = n_{\text{réd}}$ à $E$\\ Tableau d'avancement $\to$ $C$, $t$};

\draw[lien, popblue] (C) -- (B1);
\draw[lien, poporange] (C) -- (B2);
\draw[lien, popgreen] (C) -- (B3);
\draw[lien, poppurple] (C) -- (B4);
\draw[lien, popgold] (C) -- (B5);
\draw[lien, poppink] (C) -- (B6);
\end{tikzpicture}
\legende{Carte mentale de la leçon : les six idées forces des dosages d'oxydoréduction.}
\end{popfigure}

\begin{aretenir}[Synthèse de la leçon]
\textbf{Définitions clés}
\begin{itemize}
  \item \cle{Dosage} : détermination de la concentration d'une espèce (titrée) par réaction avec une solution de concentration connue (titrante).
  \item \cle{Point d'équivalence $E$} : mélange stœchiométrique des réactifs. Repéré par virage coloré (indicateur ou auto-indication).
\end{itemize}

\textbf{Équations-types à connaître par cœur}
\begin{itemize}
  \item \textbf{Fer(II) / Permanganate (milieu \ce{H2SO4})} :
  \[ \ce{MnO4- + 5Fe^{2+} + 8H+ -> Mn^{2+} + 5Fe^{3+} + 4H2O} \]
  \item \textbf{Diiode / Thiosulfate (milieu neutre)} :
  \[ \ce{I2 + 2S2O3^{2-} -> 2I- + S4O6^{2-}} \]
\end{itemize}

\textbf{Relations à l'équivalence ($V_E$ en \si{\liter})}
\begin{center}
\begin{tabularx}{\linewidth}{|l|X|X|}
\hline
\rowcolor{popblueL} \textbf{Dosage} & \textbf{Relation stœchiométrique} & \textbf{Concentration cherchée} \\
\hline
\ce{MnO4-} (ox) / \ce{Fe^{2+}} (réd) & $n(\ce{Fe^{2+}}) = 5\,n(\ce{MnO4-})$ & $C(\ce{Fe^{2+}}) = \dfrac{5\,C(\ce{MnO4-})\,V_E}{V(\ce{Fe^{2+}})}$ \\
\hline
\ce{I2} (ox) / \ce{S2O3^{2-}} (réd) & $n(\ce{S2O3^{2-}}) = 2\,n(\ce{I2})$ & $C(\ce{I2}) = \dfrac{C(\ce{S2O3^{2-}})\,V_E}{2\,V(\ce{I2})}$ \\
\hline
\end{tabularx}
\end{center}

\textbf{Points de vigilance (sécurité \& réussite)}
\begin{itemize}
  \item \ce{MnO4-} : \textbf{Jamais \ce{HCl}} (oxydation \ce{Cl-} $\to$ \ce{Cl2} toxique). Toujours \ce{H2SO4}. Auto-indication : rose persistant \SI{30}{\second}.
  \item \ce{S2O3^{2-}} : Solution instable (à préparer fraîchement, conserver basique). Indicateur \textbf{empois d'amidon} ajouté \textbf{seulement} au voisinage de $E$ (jaune paille).
  \item Lecture burette : bas du ménisque, \SI{0,1}{\milli\liter}, concordance \SI{0,2}{\milli\liter} entre deux essais.
  \item Unités : Volumes en \si{\liter} dans $n = C \times V$. Titre massique $t = C \times M$ (\si{\gram\per\liter}).
\end{itemize}
\end{aretenir}

\begin{aretenir}[Checklist avant le Bac]
\begin{itemize}
  \item[$\square$] Je sais définir dosage, solution titrée, solution titrante et point d'équivalence.
  \item[$\square$] Je cite les quatre conditions d'un bon dosage d'oxydoréduction (totale, rapide, unique, repérable).
  \item[$\square$] Je sais schématiser le dispositif : burette (titrante), erlenmeyer (titrée), agitateur magnétique.
  \item[$\square$] J'écris sans hésiter les demi-équations des couples \ce{MnO4-}/\ce{Mn^{2+}} et \ce{Fe^{3+}}/\ce{Fe^{2+}} (milieu acide).
  \item[$\square$] Je sais que \ce{MnO4-} est violet et sert de son propre indicateur : l'équivalence est repérée par la teinte rose persistante.
  \item[$\square$] Je sais que le dosage \ce{MnO4-}/\ce{Fe^{2+}} se fait en milieu acide (acide sulfurique, jamais \ce{HCl} qui serait oxydé).
  \item[$\square$] J'écris l'équation \ce{I2 + 2S2O3^{2-} -> 2I- + S4O6^{2-}} et j'utilise l'empois d'amidon comme indicateur (ajout près de $E$).
  \item[$\square$] Je construis un tableau d'avancement et j'écris la relation stœchiométrique à l'équivalence.
  \item[$\square$] Je vérifie l'homogénéité des unités : volumes en litres dans $n = C \times V$.
  \item[$\square$] Je lis le volume à l'équivalence $V_E$ au bas du ménisque, au dixième de millilitre près.
  \item[$\square$] Je sais passer de la concentration molaire au titre massique : $t = C \times M$.
  \item[$\square$] Je conclus par une phrase : « la concentration de la solution titrée est \ldots » avec le bon nombre de chiffres significatifs.
\end{itemize}
\end{aretenir}

\findecours

\end{document}
